\chapter{Galois cohomology, local and global, and Selmer groups}\label{ch:galcoh}
\chaptermark{Galois cohomology and Selmer groups}

\Cref{ch:cohomology} built the machinery of group cohomology: cochains, inflation and restriction, continuous cohomology of profinite groups, procyclic groups and
the group-theoretic core of Lemma lem:tamfactor. This chapter specialises it to Galois groups of fields, local and global, and develops the arithmetic the paper uses:
Hilbert's Theorem 90 and Kummer theory in cohomological form; the cohomology of a lattice $T$, of $V=T\otimes\Q_\ell$ and of $W=T\otimes\Q_\ell/\Z_\ell$, and their Kummer duals
$T^D$, $V^D$, $W^D$; Tate's local duality and Euler characteristic formula; unramified cohomology and the exact orthogonality of unramified classes; the Kummer class of
the Tate parameter, which identifies the corner entry of the Frobenius matrix in Lemma lem:tamfactor; the groups $G_{\Q,S}$ and the real place; the classical Selmer
group of an elliptic curve; Selmer structures and their duals; the Poitou-Tate sequence and the formula of Greenberg and Wiles; and the exact sequence comparing two
Selmer structures, of which Lemma lem:dfg21 (Lemma 2.1 of Diamond, Flach and Guo) is a special case. The deep duality theorems are stated with references; everything
the paper uses on top of them is proved. \Cref{ch:blochkato} then supplies the Bloch-Kato local conditions, the Tamagawa factors and the theorem of Diamond, Flach and Guo
that the Selmer group of $V^D$ vanishes. The chapter uses \Cref{ch:modules} for $\Z_\ell$-modules and Pontryagin duality, \Cref{ch:galois} for infinite Galois theory and
Kummer characters, \Cref{ch:ant} for decomposition and inertia groups, \Cref{ch:local} for local fields, the tame character and Kummer cocycles, \Cref{ch:ellcurves,ch:reduction,ch:galrep}
for elliptic curves, the Tate curve and the shape of $T_\ell E$ at a multiplicative prime, and \Cref{ch:cohomology} throughout.

\section*{Core results at a glance}

Throughout, $\ell$ is a prime different from the characteristic, $T$ a free $\Z_\ell$-module of finite rank with a continuous Galois action, $V=T\otimes\Q_\ell$,
$W=T\otimes\Q_\ell/\Z_\ell$, $\iota\colon H^1(-,T)\to H^1(-,V)$, $\pi\colon H^1(-,V)\to H^1(-,W)$; $\Frob$ is an arithmetic Frobenius and $\Phi=\Frob^{-1}$; ${}^*$ is the Pontryagin dual and ${}^D$
the Kummer dual.

\begin{sloppypar}
\begin{enumerate}[label=\textbf{\thechapter.\arabic*}]
\item \textbf{Conventions.} $H^n(K,M)=H^n(G_K,M)$, continuous cochains, coefficients of type (D), (C), (V); $\Frob$ arithmetic in all cohomology, $\Phi$ only in $L$-factors,
$L_p(V,s)=\det(1-\Phi p^{-s}\mid V^{I_p})^{-1}$; the Kummer connecting map is $a\mapsto[\sigma\mapsto\sigma(\alpha)/\alpha]$, $\alpha^n=a$, with no sign.
\Cref{conv:galcoh:fields,conv:galcoh:frobenius,conv:galcoh:signs}.
\item \textbf{Galois modules.} $\overline K$, $\overline K^\times$, $\mu_n$, $\Z/n(r)$, $E[n]$, $E(\overline K)$, $\mu_{\ell^\infty}$ are discrete; $\Z_\ell(1)=\varprojlim\mu_{\ell^k}$ (action by $\chi_\ell$),
$T_\ell E$ are of type (C); $G_{\F_q}\simeq\hat\Z$. \Cref{ex:galcoh:galois-modules}.
\item \textbf{Hilbert 90.} For every Galois extension $L/K$ (finite or not), $H^1(\Gal(L/K),L^\times)=0$ and $H^1(\Gal(L/K),L)=0$; so $H^1(K,\overline K^\times)=H^1(K,\overline K)=0$.
\emph{(proved here; $H^n(K,\overline K)=0$ for $n\ge2$ stated)} \Cref{thm:galcoh:hilbert90}.
\item \textbf{Kummer theory.} For $n$ prime to $\operatorname{char}K$: $\delta\colon K^\times/(K^\times)^n\xrightarrow{\sim}H^1(K,\mu_n)$, $\delta(a)=[\sigma\mapsto\sigma(\alpha)/\alpha]$, natural in $K$;
$H^2(K,\mu_n)\simeq\operatorname{Br}(K)[n]$; if $\mu_n\subset K$, $\delta(a)$ is the Kummer character. \emph{(proved here; compatibility with corestriction stated)} \Cref{thm:galcoh:kummer}.
\item \textbf{Kummer class, Kummer cocycle} (definition): $\delta(a)$, a cocycle $\sigma\mapsto\sigma(\alpha)/\alpha$, and the $\ell$-adic cocycle $\kappa_a(\sigma)=(\sigma(a^{1/\ell^k})/a^{1/\ell^k})_k\in\Z_\ell(1)$.
\Cref{def:galcoh:kummer-class}.
\item \textbf{Properties.} $\delta_K(a)|_{G_L}=0$ iff $a\in(L^\times)^n$; for $\mu_n\subset K$, $\ker\chi_a=G_{K(\alpha)}$ and $[K(\alpha):K]$ is the order of $a$ in $K^\times/(K^\times)^n$; $\kappa_a$
reduces modulo $\ell^k$ to a Kummer cocycle. \emph{(proved here)} \Cref{prop:galcoh:kummer-class}.
\item \textbf{Finite fields.} $H^1(\F_q,\mu_{\ell^k})\simeq\F_q^\times/(\F_q^\times)^{\ell^k}\simeq\mu_{\ell^k}/\mu_{\ell^k}^{q-1}$, $u\mapsto u^{(q-1)/\ell^k}$ ($k\le v=v_\ell(q-1)$);
$H^1(\F_q,\Z_\ell(1))=\Z_\ell(1)/(q-1)\simeq\F_q^\times\otimes\Z_\ell\simeq\Z/\ell^v$; $v_\ell$ of a class, with $v_\ell(0)=v$. \emph{(proved here, given \Cref{ch:cohomology})}
\Cref{ex:galcoh:kummer-finite-field}.
\item \textbf{Brauer groups.} $\operatorname{Br}(\F_q)=\operatorname{Br}(\C)=0$, $\operatorname{Br}(\R)=\Z/2$; $\operatorname{inv}_K\colon\operatorname{Br}(K)\simeq\Q/\Z$ for $K/\Q_p$ finite, so $H^2(K,\mu_n)\simeq\Z/n$; the sequence
$0\to\operatorname{Br}(\Q)\to\bigoplus_v\operatorname{Br}(\Q_v)\xrightarrow{\sum\operatorname{inv}_v}\Q/\Z\to0$ is exact. \emph{(stated; reference)} \Cref{thm:galcoh:brauer}.
\item \textbf{Lattices and duals} (convention and definition): $W[\ell^k]\simeq T/\ell^kT$; $M^D=\Hom(M,\mu)$ for finite $M$; $T^D=\Hom(T,\Z_\ell(1))$, $V^D=T^D\otimes\Q_\ell$, $W^D=V^D/T^D$.
\Cref{conv:galcoh:duals,def:galcoh:kummer-duals}.
\item \textbf{Kummer duality.} The pairings $T\times T^D\to\Z_\ell(1)$, $V\times V^D\to\Q_\ell(1)$, $W\times T^D\to\Q_\ell/\Z_\ell(1)$, $T\times W^D\to\Q_\ell/\Z_\ell(1)$, $M\times M^D\to\mu$ are perfect and
equivariant; $W^D\simeq T^*(1)$, $(T/\ell^k)^D\simeq T^D/\ell^k$, $(T^D)^D=T$; $E[n]^D\simeq E[n]$ by the Weil pairing. \emph{(proved here)} \Cref{prop:galcoh:kummer-duals}.
\item \textbf{The cohomology of $T$, $V$, $W$.} If $H^1(G,T/\ell)$ and $H^2(G,T/\ell)$ are finite: $H^n(G,T)$ is finitely generated ($n\le2$) and $H^n(G,V)=H^n(G,T)\otimes\Q_\ell$;
$H^1(G,T)_{\tors}=\ker\iota=\delta H^0(G,W)\simeq H^0(G,W)/H^0(G,W)_{\mathrm{div}}$, and $\simeq H^0(G,W)$ if $H^0(G,V)=0$;
$0\to H^1(G,T)\otimes\Q_\ell/\Z_\ell\to H^1(G,W)\to H^2(G,T)_{\tors}\to0$ with the first term $=\operatorname{im}\pi=H^1(G,W)_{\mathrm{div}}$; the finite-level sequences.
\emph{(proved here, given \Cref{ch:cohomology})} \Cref{prop:galcoh:tvw}.
\item \textbf{$\ell$-adic Kummer theory.} $H^1(K,\Z_\ell(1))\simeq\widehat{K^\times}=\varprojlim K^\times/(K^\times)^{\ell^k}$, $a\mapsto\kappa_a$; $H^1(K,\Q_\ell/\Z_\ell(1))\simeq K^\times\otimes\Q_\ell/\Z_\ell$.
\emph{(proved here, given \Cref{ch:cohomology})} \Cref{thm:galcoh:kummer-ladic}.
\item \textbf{Local Kummer theory} ($K/\Q_p$ finite, $\ell\ne p$, residue field $\F_q$, $v=v_\ell(q-1)$): $0\to\mathcal O_K^\times\otimes\Z_\ell\to\widehat{K^\times}\xrightarrow{v_K}\Z_\ell\to0$, so
$H^1(K,\Z_\ell(1))\simeq\Z_\ell\times\Z/\ell^v$ with torsion the unit classes; $\#H^1(K,\mu_{\ell^k})=\ell^{k+\min(k,v)}$; $H^1(K,\Q_\ell/\Z_\ell(1))\simeq\Q_\ell/\Z_\ell$; $H^1(K,\Q_\ell(1))\simeq\Q_\ell$.
\emph{(proved here)} \Cref{prop:galcoh:kummer-local}.
\item \textbf{Local finiteness.} For $K/\Q_p$ finite and $M$ finite, $H^n(K,M)$ is finite and $0$ for $n\ge3$; hence $H^n(K,T)$ is finitely generated and \Cref{prop:galcoh:tvw}
applies to $G_K$. \emph{(stated; reference; sketched for $H^1$ with $p\nmid\#M$)} \Cref{thm:galcoh:local-finiteness}.
\item \textbf{Local Euler characteristic.} $\#H^0(K,M)\#H^2(K,M)/\#H^1(K,M)=|\#M|_K$; for $V$: $\sum(-1)^i\dim H^i(K,V)=-[K:\Q_p]\dim V$ if $\ell=p$, $0$ if $\ell\ne p$.
\emph{(stated; reference)} \Cref{thm:galcoh:local-euler}.
\item \textbf{Tate local duality.} $H^i(K,M)\times H^{2-i}(K,M^D)\to H^2(K,\mu)\simeq\Q/\Z$ is perfect; $H^i(K,W)\times H^{2-i}(K,T^D)\to\Q_\ell/\Z_\ell$, $H^i(K,T)\times H^{2-i}(K,W^D)\to\Q_\ell/\Z_\ell$ and
$H^i(K,V)\times H^{2-i}(K,V^D)\to\Q_\ell$ are perfect; $\langle\pi x,y\rangle=\langle x,\iota y\rangle\bmod\Z_\ell$. \emph{(stated; reference)} \Cref{thm:galcoh:tate-local-duality}.
\item \textbf{Consequences.} $H^2(K,M)\simeq H^0(K,M^D)^*$; for $p\nmid\#M$, $\#H^1(K,M)=\#H^0(K,M)\#H^0(K,M^D)$; for $\ell\ne p$, $\dim H^1(K,V)=\dim H^0(K,V)+\dim H^0(K,V^D)$ and
$H^2(K,T)\simeq H^0(K,W^D)^*$. \emph{(proved here, given \Cref{thm:galcoh:local-euler,thm:galcoh:tate-local-duality})} \Cref{cor:galcoh:local-duality-consequences}.
\item \textbf{Orders over $\Q_p$.} $\#H^i(\Q_p,\Z/\ell^k)=\ell^k,\ell^{k+m},\ell^m$ and $\#H^i(\Q_p,\mu_{\ell^k})=\ell^m,\ell^{k+m},\ell^k$ ($m=\min(k,v_\ell(p-1))$, $p\ne\ell$); $\mu_p$ over $\Q_p$ has
characteristic $p^{-1}$; $E[5]\subset E(\Q_{11})$ for \curve{11a1}. \emph{(proved here, given the stated theorems)} \Cref{ex:galcoh:local-orders}.
\item \textbf{Unramified classes} (definition; $K/\Q_p$ finite): $\Hur(K,M)=\operatorname{im}(\Inf\colon H^1(G_K/I_K,M^{I_K})\to H^1(K,M))=\ker(H^1(K,M)\to H^1(I_K,M))\simeq H^1(G_K/I_K,M^{I_K})$.
\Cref{def:galcoh:unramified}.
\item \textbf{Computation.} $\Hur(\Q_p,M)\simeq M^I/(\Frob-1)M^I$; $\#\Hur(\Q_p,M)=\#H^0(\Q_p,M)$ for finite $M$; $\Hur(\Q_p,V)=0$ iff $H^0(\Q_p,V)=0$ iff $\det(1-\Frob\mid V^I)\ne0$ iff
$\det(1-\Phi\mid V^I)\ne0$; for $p\ne\ell$, $\#\Hur(\Q_p,T)=\ell^{v_\ell\det(1-\Frob\mid V^I)}=\ell^{v_\ell\det(1-\Phi\mid V^I)}$ when finite. \emph{(proved here, given \Cref{ch:cohomology})}
\Cref{prop:galcoh:unramified-computed}.
\item \textbf{The inertia sequence.} $0\to\Hur(\Q_p,M)\to H^1(\Q_p,M)\to H^1(I,M)^{G/I}\to0$ is exact (conjugation action), for $M$ of type (D), and for types (C), (V) at $p\ne\ell$.
\emph{(proved here, given \Cref{ch:cohomology} and \Cref{thm:galcoh:local-finiteness}; types (C), (V) by passage to the limit)} \Cref{prop:galcoh:ur-sequence}.
\item \textbf{Exact orthogonality} ($p\ne\ell$). $\Hur(\Q_p,M)$ and $\Hur(\Q_p,M^D)$ are exact annihilators ($M$ finite of $\ell$-power order); so are $\Hur(\Q_p,V)$ and $\Hur(\Q_p,V^D)$.
False for $p=\ell$. \emph{(proved here, given \Cref{thm:galcoh:local-euler,thm:galcoh:tate-local-duality})} \Cref{thm:galcoh:ur-orthogonal}.
\item \textbf{Propagation.} For $L_V\subseteq H^1(K,V)$: $L_T=\iota^{-1}L_V\supseteq H^1(K,T)_{\tors}$, $L_W=\pi L_V$ divisible; if $L_V^\perp=L_{V^D}$ then $L_W^\perp=L_{T^D}$ and
$(H^1(K,W)/L_W)^*\simeq L_{T^D}$; for $T$ unramified, $p\ne\ell$, the propagation of $\Hur$ is $\Hur$. \emph{(proved here, given \Cref{thm:galcoh:tate-local-duality})}
\Cref{prop:galcoh:propagation}.
\item \textbf{Vanishing unramified cohomology.} If $p\ne\ell$ and $\Hur(\Q_p,V)=0$: $H^0(\Q_p,V)=0$, $\iota^{-1}(0)=H^1(\Q_p,T)_{\tors}\simeq H^0(\Q_p,W)$ is finite and contains $\Hur(\Q_p,T)$, and
$\pi(0)=0$. (The paper: ``$H^1_f(\Q_p,T^D)$ is the torsion of $H^1(\Q_p,T^D)$, a finite group''.) \emph{(proved here, given \Cref{thm:galcoh:local-finiteness})}
\Cref{cor:galcoh:ur-vanishing}.
\item \textbf{Examples.} $\Hur(\Q_p,\Z_\ell(1))=\Z_\ell(1)/(p-1)$; $H^1(\Q_p,\Z_\ell)=\Hur$; Tate curve: $\Hur(\Q_p,V_\ell E)=0$; $\ad V_\ell E$ at a multiplicative prime: $\det(1-\Phi\mid(V^D)^I)=1-p^{-2}$,
$\det(1-\Frob\mid(V^D)^I)=1-p^2$, $\#\Hur(\Q_p,T^D)=\ell^{v_\ell(p^2-1)}$; types $\mathrm{III}$, $\mathrm{III}^*$: $\ell^{v_\ell(q-\chi_{-4}(q))}$. \emph{(proved here, given \Cref{ch:galrep}; the last an illustration depending on \Cref{ch:lfunctions})}
\Cref{ex:galcoh:ur-examples}.
\item \textbf{Unit classes.} For $K/\Q_p$ finite, $p\nmid n$: the Kummer isomorphism maps $\mathcal O_K^\times$ onto $\Hur(K,\mu_n)$; the valuation sequence of $\widehat{K^\times}$ is the inertia
sequence for $\Z_\ell(1)$, with $\Res\kappa_x=v_K(x)t_\ell$ and $\Hom(I_K,\Z_\ell(1))^{G/I}=\Z_\ell t_\ell$. \emph{(proved here)} \Cref{prop:galcoh:unit-classes}.
\item \textbf{Value of an unramified class.} $H^1(G_{\Q_q}/I_q,\Z_\ell(1))\simeq\Z_\ell(1)/(q-1)$; for $u\in\Z_q^\times$, $\kappa_u$ is inflated as a cocycle and $\kappa_u(\Frob)$ is the class of $\bar u$ in
$\F_q^\times\otimes\Z_\ell$. \emph{(proved here)} \Cref{prop:galcoh:unramified-value}.
\item \textbf{The splitting given by $q$.} $H^1(\Q_q,\Z_\ell(1))=\Z_\ell[\kappa_q]\oplus\Hur(\Q_q,\Z_\ell(1))$ canonically; $\operatorname{pr}_{\ur}[\kappa_x]$ is the class of the unit part of $x$.
\emph{(proved here)} \Cref{prop:galcoh:canonical-splitting}.
\item \textbf{The corner of $\rho(\Frob)$.} For $E$ multiplicative at $q\ne\ell$, $\ell$ odd, with $\rho=\sm{\chi_\ell\varepsilon}{\varepsilon\kappa}0\varepsilon$: $\kappa=n\kappa_q+\kappa_u$,
$\operatorname{pr}_{\ur}[\kappa]$ is the class of $\bar u$; $\kappa|_I=nt_\ell$ (so $n'=n$); for the lift fixing the roots of $q$, $c'=\varepsilon\kappa_u(\Frob)$, of class $\pm$ the class of $\bar u$; in
general $c'$ is defined modulo $(n,q-1)\Z_\ell$; $\min(v_\ell(n),v_\ell(q-1),v_\ell(c'))=\min(v_\ell(n),v_\ell(q-1),v_\ell(c))$. \emph{(proved here, given \Cref{ch:local,ch:galrep})}
\Cref{thm:galcoh:tate-kummer-class}.
\item \textbf{The formula of Lemma lem:tamfactor} (a pointer). $\len\bigl((H^1(I,\ad T)_{\tors})^{G/I}\bigr)=v_\ell(n)+\min(v_\ell(n),v_\ell(q-1),v_\ell(c))$: the model computation of
\Cref{ch:cohomology} with the identification of $n'$, $c'$ above; the lemma is assembled in \Cref{ch:blochkato}. \emph{(proved here, given \Cref{ch:cohomology})} \Cref{cor:galcoh:tamfactor}.
\item \textbf{Tate parameters at $\ell=3$.} \curve{77c1}, \curve{791a1}, \curve{1577a1} at $q=7,7,19$: extra term $0,1,1$, bound $1,2,2=v_3(\degphi)$; \curve{11a1} at $\ell=5$: Eisenstein.
\emph{(computed)} \Cref{ex:galcoh:tate-kummer-examples}.
\item \textbf{Lemma lem:tamagawa-exp, cohomologically.} $\rho\bmod\ell^t$ is unramified at $q$ iff $nt_\ell\equiv0\pmod{\ell^t}$ iff $\ell^t\mid n$ iff $q_E$ is an $\ell^t$-th power in $\Q_q^{\ur}$.
\emph{(proved here)} \Cref{rem:galcoh:tamagawa-exp-cohomological}.
\item \textbf{$G_{\Q,S}$} (definition): $G_\Q/H_S$, $H_S$ generated by the inertia groups above $p\notin S$; localisation $\Res_v$ is independent of the embedding. \Cref{def:galcoh:gqs}.
\item \textbf{Unramified outside $S$.} $H^1(G_{\Q,S},M)=\{c\in H^1(\Q,M):\Res_pc\in\Hur\ \forall p\notin S\}$; for finite $M$ every class is unramified almost everywhere. \emph{(proved here)}
\Cref{prop:galcoh:unramified-almost-everywhere}.
\item \textbf{Finiteness for $G_{\Q,S}$.} $H^1(G_{\Q,S},M)$ and $H^2(G_{\Q,S},M)$ are finite for finite $M$; so \Cref{prop:galcoh:tvw} applies to $G_{\Q,S}$. \emph{($H^1$ sketched, given
Hermite-Minkowski; $H^2$ and $\operatorname{cd}$ stated; reference)} \Cref{thm:galcoh:gqs-finite}.
\item \textbf{The real place.} $2H^n(\R,M)=0$; for $\ell$ odd $H^n(\R,M)=0$ ($n\ge1$) and $\hat H^0(\R,M)=0$; so $H^1_f(\R,T^D)=0$; $H^1(\R,\mu_2)=\Z/2$. \emph{(proved here)}
\Cref{prop:galcoh:real-place}.
\item \textbf{Global Kummer theory.} $H^1(G_{\Q,S},\mu_n)\simeq\Z_S^\times/(\Z_S^\times)^n$, finite. \emph{(proved here)} \Cref{ex:galcoh:global-kummer}.
\item \textbf{Kummer sequence of $E$.} $0\to E(K)/nE(K)\xrightarrow{\delta}H^1(K,E[n])\to H^1(K,E)[n]\to0$, $\delta(P)=[\sigma\mapsto\sigma Q-Q]$. \emph{(proved here)} \Cref{prop:galcoh:ec-kummer}.
\item \textbf{Local Kummer image} ($p\nmid n$). $\#E(\Q_p)/n=\#E(\Q_p)[n]=\#\Hur(\Q_p,E[n])$; for good reduction the image is $\Hur$; it is its own annihilator (Tate); for $W=E[\ell^\infty]$ at a
multiplicative prime, $\#\Hur(\Q_p,W)$ is the $\ell$-part of $c_p$. \emph{(proved here; the self-duality stated)} \Cref{prop:galcoh:ec-local-image}.
\item \textbf{Selmer and Sha} (definition and proposition): $\Sel_n(E/\Q)$, $\Sha(E/\Q)$, and $0\to E(\Q)/n\to\Sel_n(E/\Q)\to\Sha(E/\Q)[n]\to0$. \emph{(proved here)}
\Cref{def:galcoh:selmer-sha,prop:galcoh:selmer-sha}.
\item \textbf{Finiteness; weak Mordell-Weil.} $\Sel_n(E/\Q)\subseteq H^1(G_{\Q,S},E[n])$ is finite; $E(\Q)/nE(\Q)$ is finite. \emph{(proved here, given \Cref{thm:galcoh:gqs-finite})}
\Cref{thm:galcoh:selmer-finite,cor:galcoh:weak-mw}.
\item \textbf{Selmer structures} (definition): local conditions $\mathcal L_v$, unramified almost everywhere; $H^1_{\mathcal L}(\Q,M)$; relaxed $\mathcal L^\Sigma$, strict $\mathcal L_\Sigma$, classical and
propagated structures. \Cref{def:galcoh:selmer-structure,ex:galcoh:selmer-structures}.
\item \textbf{Dual structures} (definition and proposition): $\mathcal L^*_v=\mathcal L_v^\perp$ on $M^D$, or on $T^D$ for structures on $W$; $\mathcal L^*$ is a Selmer structure, $\mathcal L^{**}=\mathcal L$,
$(\mathcal L^\Sigma)^*=(\mathcal L^*)_\Sigma$, and the dual of a propagated structure is propagated. \emph{(proved here, given \Cref{thm:galcoh:tate-local-duality})}
\Cref{def:galcoh:dual-structure,prop:galcoh:dual-structure}.
\item \textbf{$\ell$-adic Selmer groups.} $H^1_{\mathcal L}(\Q,W)$ is cofinitely generated; if $H^1_{L_V}(\Q,V)=0$ then $H^1_{\mathcal L}(\Q,T)=H^1(G_{\Q,S},T)_{\tors}\simeq H^0(\Q,W)/\mathrm{div}$;
$\operatorname{corank}H^1_{\mathcal L}(\Q,W)=\dim H^1_{L_V}(\Q,V)$. \emph{(proved here, given \Cref{thm:galcoh:gqs-finite}; the corank formula stated)} \Cref{prop:galcoh:selmer-l-adic}.
\item \textbf{Poitou-Tate.} The nine-term exact sequence for $M$ finite unramified outside $S$; the image of $H^1(G_S,M)$ in $\bigoplus_{v\in S}H^1(\Q_v,M)$ is the exact annihilator of
that of $H^1(G_S,M^D)$. \emph{(stated; reference)} \Cref{thm:galcoh:poitou-tate}.
\item \textbf{Reciprocity.} $\sum_v\operatorname{inv}_v(\Res_vx\cup\Res_vy)=0$ for $x\in H^1(\Q,M)$, $y\in H^1(\Q,M^D)$. \emph{(proved here, given \Cref{thm:galcoh:brauer})} \Cref{prop:galcoh:reciprocity}.
\item \textbf{Greenberg-Wiles.} $\#H^1_{\mathcal L}(\Q,M)/\#H^1_{\mathcal L^*}(\Q,M^D)=\frac{\#H^0(\Q,M)}{\#H^0(\Q,M^D)}\prod_v\frac{\#\mathcal L_v}{\#H^0(\Q_v,M)}$. \emph{(stated; reference; finiteness proved)}
\Cref{thm:galcoh:greenberg-wiles}.
\item \textbf{A check by hand.} For $M=\Z/\ell$ unramified everywhere both sides are $1$; relaxed at $p$, the formula counts the characters of order $\ell$ and conductor $p$.
\emph{(proved here, given the stated theorems)} \Cref{ex:galcoh:gw-check}.
\item \textbf{Comparison of orders.} For $\mathcal L\subseteq\mathcal L'$: $\#H^1_{\mathcal L'}/\#H^1_{\mathcal L}=\prod_v\#\mathcal L'_v/\#\mathcal L_v\cdot\#H^1_{\mathcal L'^*}/\#H^1_{\mathcal L^*}$. \emph{(proved here, given
\Cref{thm:galcoh:greenberg-wiles})} \Cref{cor:galcoh:gw-comparison}.
\item \textbf{The comparison sequence.} $0\to H^1_{\mathcal L}\to H^1_{\mathcal L'}\to\bigoplus_v\mathcal L'_v/\mathcal L_v\to(H^1_{\mathcal L^*})^*\to(H^1_{\mathcal L'^*})^*\to0$ is exact ($M$ finite). \emph{(complex property
proved here; exactness proved here, given \Cref{thm:galcoh:poitou-tate,thm:galcoh:greenberg-wiles})} \Cref{thm:galcoh:comparison}.
\item \textbf{The comparison sequence for $(W,T^D)$.} The same with $W$ and the dual structures on $T^D$. \emph{(sketched here; reference)} \Cref{thm:galcoh:comparison-l-adic}.
\item \textbf{Relaxing at $\Sigma$} (the abstract form of DFG04, Lemma 2.1). For propagated structures with $L_{V,v}^\perp=L_{V^D,v}$: the five-term sequence
$0\to H^1_{\mathcal L}(\Q,W)\to H^1_{\mathcal L^\Sigma}(\Q,W)\to\bigoplus_{p\in\Sigma}(\mathcal L^*_p)^*\to H^1_{\mathcal L^*}(\Q,T^D)^*\to H^1_{(\mathcal L^*)_\Sigma}(\Q,T^D)^*\to0$; if $H^1_{L_{V^D}}(\Q,V^D)=0$ and
$H^0(\Q_p,V^D)=0$ for one $p\in\Sigma$, then $0\to H^1_{\mathcal L}(\Q,W)\to H^1_{\mathcal L^\Sigma}(\Q,W)\to\bigoplus_{p\in\Sigma\cup\{\infty\}}(\mathcal L^*_p)^*\to H^0(\Q,W^D)^*\to0$.
\emph{(proved here, given \Cref{thm:galcoh:comparison-l-adic})} \Cref{thm:galcoh:relaxing-sigma}.
\item \textbf{The length identity.} $\len H^1_{\mathcal L^\Sigma}(\Q,W)=\len H^1_{\mathcal L}(\Q,W)+\sum_{p\in\Sigma}\len\mathcal L^*_p-\len H^0(\Q,W^D)$. \emph{(proved here, given
\Cref{thm:galcoh:relaxing-sigma})} \Cref{cor:galcoh:length-identity}.
\item \textbf{A toy case.} $T=\Z_\ell$, $\Sigma=\{p\}$: the sequence identifies the characters of $\ell$-power order unramified outside $p$ with $(\Z_p^\times\otimes\Z_\ell)^*$.
\emph{(proved here, given the stated theorems)} \Cref{ex:galcoh:relaxing-toy}.
\item \textbf{What \Cref{ch:blochkato} adds.} The Bloch-Kato conditions and their orthogonality at $\ell$, Theorem thm:dfg27, Lemma lem:fpr, and the value of the Euler factor.
\Cref{rem:galcoh:handover-ch18}.
\end{enumerate}
\end{sloppypar}

\section*{Where the paper uses this chapter}

\begin{itemize}
\item \emph{Lemma lem:tamfactor (the Tamagawa factor at a multiplicative prime), the end of its proof}: the group $H^1(G,\Z_\ell(1))=\varprojlim\Q_q^\times/(\Q_q^\times)^{\ell^k}$
(\Cref{thm:galcoh:kummer-ladic,prop:galcoh:kummer-local}); the splitting $\Z_\ell\kappa_q\oplus H^1_{\mathrm{ur}}$ and the class of the unit part
(\Cref{prop:galcoh:unit-classes,prop:galcoh:unramified-value,prop:galcoh:canonical-splitting}); the value $c'$ and its indeterminacy (\Cref{thm:galcoh:tate-kummer-class}); the
assembled formula (\Cref{cor:galcoh:tamfactor}, with \Cref{ch:cohomology,ch:blochkato}); its check with $M=T$ (\Cref{prop:galcoh:ec-local-image}(iv)).
\item \emph{Lemma lem:tamagawa-exp}: ``By Kummer theory over $\Q_q^{\mathrm{ur}}$'' (\Cref{prop:galcoh:kummer-class,rem:galcoh:tamagawa-exp-cohomological}).
\item \emph{The paragraph after Lemma lem:dfg21}: $H^1_{\mathrm{ur}}(\Q_p,V^D)=(V^D)^{I_p}/(\Frob_p-1)$ and its vanishing (\Cref{prop:galcoh:unramified-computed},
\Cref{warn:galcoh:frobenius}); ``$H^1_f(\Q_p,T^D)$ is the torsion of $H^1(\Q_p,T^D)$, a finite group'' (\Cref{prop:galcoh:tvw}, \Cref{cor:galcoh:ur-vanishing}); ``$H^1(\R,T^D)$ is
killed by $2$'' (\Cref{prop:galcoh:real-place}).
\item \emph{Lemma lem:dfg21 (DFG04, Lemma 2.1)}: Tate local duality and propagation (\Cref{thm:galcoh:tate-local-duality,prop:galcoh:propagation}), the exact orthogonality at
$p\ne\ell$ (\Cref{thm:galcoh:ur-orthogonal}), Selmer structures and their duals (\Cref{def:galcoh:selmer-structure,prop:galcoh:dual-structure,prop:galcoh:selmer-l-adic}), and
the comparison sequences (\Cref{thm:galcoh:comparison,thm:galcoh:comparison-l-adic,thm:galcoh:relaxing-sigma}).
\item \emph{The proof of Theorem thm:main and Remark rem:ns}: the length identity (\Cref{cor:galcoh:length-identity}), with $H^0(\Q,W^D)=0$, resp. of length $h_3(E)=1$.
\item \emph{The Selmer groups of Theorem thm:main and Section 4.2}: $H^1_f(\Q,\ad T\otimes\Q_\ell/\Z_\ell)$ and $H^1_\Sigma(\Q,W)$ as Selmer groups of a structure and its relaxation
(\Cref{ex:galcoh:selmer-structures}).
\item \emph{Section 1.2 ($\Tam_p$) and Lemma lem:fpr, through \Cref{ch:blochkato}}: the inertia sequence and the order of $\Hur(\Q_p,T^D)$
(\Cref{prop:galcoh:ur-sequence,prop:galcoh:unramified-computed}, \Cref{ex:galcoh:ur-examples}(d), (e)).
\item \emph{Section sec:tam and Remark rem:ko(1)}: the extra term at a multiplicative prime (\Cref{thm:galcoh:tate-kummer-class}, \Cref{ex:galcoh:tate-kummer-examples}).
\end{itemize}

\section{Galois cohomology of a field; Hilbert 90; Kummer theory}\label{sec:galcoh:field}

This section specialises the cohomology of profinite groups of \Cref{ch:cohomology} to absolute Galois groups. Its two theorems, Hilbert's
Theorem 90 in cohomological form and Kummer theory without the hypothesis $\mu_n\subset K$, are the reason the paper may write
$H^1(G,\Z_\ell(1))=\varprojlim\Q_q^\times/(\Q_q^\times)^{\ell^k}$ in the proof of Lemma lem:tamfactor and argue ``by Kummer theory'' in the proof of
Lemma lem:tamagawa-exp. The finite field computation (\Cref{ex:galcoh:kummer-finite-field}) is the identification
$\Z_\ell(1)/(q-1)\simeq\F_q^\times\otimes\Z_\ell$ that the same proof uses.

\begin{convention}[fields, Galois groups, coefficients]\label{conv:galcoh:fields}
$K$ is a field, $\overline K$ a separable closure and $G_K=\Gal(\overline K/K)$ with the Krull topology (\Cref{ch:galois}). For a $G_K$-module $M$ we write
$H^n(K,M):=H^n(G_K,M)$, continuous cochain cohomology as in \Cref{ch:cohomology}, for the three types of coefficients used there: (D) discrete modules
(every element has an open stabiliser); (C) finitely generated $\Z_\ell$-modules with the $\ell$-adic topology and a continuous action; (V) finite-dimensional
$\Q_\ell$-vector spaces with a continuous action. Whenever $\ell$-adic coefficients occur, $\operatorname{char}K\ne\ell$. For $K=\Q_p$ the inertia group
$I_p\trianglelefteq G_{\Q_p}$ has quotient $G_{\Q_p}/I_p=\Gal(\overline\F_p/\F_p)\simeq\hat\Z$; for $K=\R$, $G_\R=\{1,c\}$ with $c$ complex conjugation.
For $K\subseteq L\subseteq\overline K$ the restriction $\Res\colon H^n(K,M)\to H^n(L,M)$ is induced by $G_L\subseteq G_K$; for $L/K$ Galois, inflation-restriction
(\Cref{ch:cohomology}) gives the exact sequence $0\to H^1(\Gal(L/K),M^{G_L})\to H^1(K,M)\to H^1(L,M)^{\Gal(L/K)}$.
\end{convention}

\begin{convention}[Frobenius]\label{conv:galcoh:frobenius}
$\Frob=\Frob_p$ is an \emph{arithmetic} Frobenius: an element of $G_{\Q_p}$ inducing $x\mapsto x^p$ on $\overline\F_p$, well defined modulo $I_p$. It acts on $\mu_n$
($p\nmid n$) by $\zeta\mapsto\zeta^p$, so $\chi_\ell(\Frob_p)=p$ for $\ell\ne p$, where $\chi_\ell$ is the $\ell$-adic cyclotomic character. $\Phi=\Frob^{-1}$ is the
geometric Frobenius. Every cohomological statement of this chapter (conjugation actions, unramified cohomology, the matrices of Lemma lem:tamfactor) uses
$\Frob$; $\Phi$ appears only where a local $L$-factor is named, and then $L_p(V,s)=\det(1-\Phi\,p^{-s}\mid V^{I_p})^{-1}$ (\Cref{ch:local}). An eigenvalue
$\lambda$ of $\Frob$ is the eigenvalue $\lambda^{-1}$ of $\Phi$.
\end{convention}

\begin{convention}[connecting maps]\label{conv:galcoh:signs}
For an exact sequence $0\to A\to B\to C\to0$ of $G$-modules the connecting map $H^0(G,C)\to H^1(G,A)$ sends $c$ to the class of $g\mapsto gb-b$, where $b\in B$
lifts $c$ (\Cref{ch:cohomology}). For the sequence $1\to\mu_n\to\overline K^\times\to\overline K^\times\to1$, written multiplicatively, this is
$a\mapsto[\sigma\mapsto\sigma(\alpha)/\alpha]$ with $\alpha^n=a$: the convention of the Kummer characters of \Cref{ch:galois} and of the Kummer cocycles of
\Cref{ch:local}. No sign is introduced anywhere in this chapter.
\end{convention}

Recall from \Cref{ch:cohomology} that for a discrete $G_K$-module $M$ every continuous cochain factors through a finite quotient, so
\begin{equation}\label{eq:galcoh:direct-limit}
H^n(K,M)=\varinjlim_LH^n\bigl(\Gal(L/K),M^{G_L}\bigr),
\end{equation}
the limit over the finite Galois extensions $L/K$ inside $\overline K$, with inflation as transition maps. In particular $H^0(K,M)=M^{G_K}$, and for $M$ with trivial action
$H^1(K,M)=\Hom_{\mathrm{cont}}(G_K,M)$.

\begin{example}[Galois modules]\label{ex:galcoh:galois-modules}
The following are discrete $G_K$-modules: the additive group $\overline K$; the multiplicative group $\overline K^\times$; $\mu_n=\mu_n(\overline K)$ for $n$ prime to
$\operatorname{char}K$; $\Z/n$ with trivial action and its twists $\Z/n(r)=\mu_n^{\otimes r}$; for an elliptic curve $E/K$, the group $E(\overline K)$ and its
subgroup $E[n]$, $\operatorname{char}K\nmid n$ (\Cref{ch:ellcurves}); the $\ell$-primary modules $\mu_{\ell^\infty}=\bigcup_k\mu_{\ell^k}$ and $E[\ell^\infty]$. The
following are of type (C): $\Z_\ell(1)=\varprojlim_k\mu_{\ell^k}$ (transition maps $\zeta\mapsto\zeta^\ell$), on which $\sigma$ acts by multiplication by $\chi_\ell(\sigma)$;
the Tate module $T_\ell E$; $\ad T_\ell E$. For a finite field $K=\F_q$ the group $G_K$ is procyclic, $G_K\simeq\hat\Z$, topologically generated by
$\Frob_q\colon x\mapsto x^q$ (\Cref{ch:galois}).
\end{example}

\begin{theorem}[Hilbert's Theorem 90]\label{thm:galcoh:hilbert90}
Let $K$ be a field.
\begin{enumerate}[label=(\roman*)]
\item For every finite Galois extension $L/K$, $H^1(\Gal(L/K),L^\times)=0$ and $H^1(\Gal(L/K),L)=0$.
\item More generally, for every Galois extension $L/K$, finite or not, $H^1(\Gal(L/K),L^\times)=0$ and $H^1(\Gal(L/K),L)=0$ (continuous cochains, $L$ and
$L^\times$ discrete).
\item In particular $H^1(K,\overline K^\times)=0$ and $H^1(K,\overline K)=0$.
\item (Stated.) $H^n(K,\overline K)=0$ for all $n\ge1$.
\end{enumerate}
\end{theorem}

\begin{proof}
(i) This is Hilbert's Theorem 90 of \Cref{ch:galois}; we recall the argument because it is short. Let $G=\Gal(L/K)$ and let $(a_\sigma)$ be a $1$-cocycle with
values in $L^\times$, $a_{\sigma\tau}=a_\sigma\,\sigma(a_\tau)$. By Dedekind's independence of characters the map $c\mapsto\sum_{\tau\in G}a_\tau\tau(c)$ is not
identically zero; choose $c$ with $b=\sum_\tau a_\tau\tau(c)\ne0$. Then
\[
\sigma(b)=\sum_\tau\sigma(a_\tau)\,\sigma\tau(c)=\sum_\tau a_\sigma^{-1}a_{\sigma\tau}\,\sigma\tau(c)=a_\sigma^{-1}b,
\]
so $a_\sigma=b/\sigma(b)=\sigma(b^{-1})/b^{-1}$ is a coboundary. In the additive case, $a_{\sigma\tau}=a_\sigma+\sigma(a_\tau)$; as $L/K$ is separable the trace is
surjective, so there is $c$ with $\operatorname{Tr}_{L/K}(c)=1$, and $b=\sum_\tau a_\tau\tau(c)$ satisfies $\sigma(b)=\sum_\tau(a_{\sigma\tau}-a_\sigma)\sigma\tau(c)=b-a_\sigma$, so
$a_\sigma=b-\sigma(b)$ is a coboundary.

(ii) $L^\times$ is a discrete $\Gal(L/K)$-module, and for a finite Galois subextension $L'/K$ of $L/K$ the invariants of $\Gal(L/L')$ in $L^\times$ are $L'^\times$. By
\eqref{eq:galcoh:direct-limit} for the profinite group $\Gal(L/K)$, $H^1(\Gal(L/K),L^\times)=\varinjlim_{L'}H^1(\Gal(L'/K),L'^\times)$, a direct limit of zero groups by
(i). The additive case is the same. (iii) is (ii) with $L=\overline K$.

(iv) For finite $G$ the additive group $L$ is an induced $G$-module by the normal basis theorem, hence has vanishing cohomology in positive degrees;
the limit argument of (ii) applies. We do not use this; see \cite[Chapter X, Section 1]{SerreLocalFields} and \cite[Chapter II, Section 1]{SerreGaloisCoh}.
\end{proof}

\begin{remark}
Part (ii) applies to $L=\bar k(C)$, $K=k(C)$ for a curve $C$ over a perfect field $k$ (the extension $\bar k(C)/k(C)$ is Galois with group $G_k$): this is the
form in which \Cref{ch:curves} uses Hilbert 90 to show that a $G_k$-invariant principal divisor is the divisor of a function defined over $k$.
\end{remark}

\begin{theorem}[Kummer theory in cohomological form]\label{thm:galcoh:kummer}
Let $n\ge1$ be prime to $\operatorname{char}K$.
\begin{enumerate}[label=(\roman*)]
\item The sequence $1\to\mu_n\to\overline K^\times\xrightarrow{\ x\mapsto x^n\ }\overline K^\times\to1$ is an exact sequence of discrete $G_K$-modules, and its connecting map
\[
\delta=\delta_K\colon K^\times/(K^\times)^n\longrightarrow H^1(K,\mu_n),\qquad\delta(a)=\bigl[\sigma\mapsto\sigma(\alpha)/\alpha\bigr]\quad(\alpha\in\overline K,\ \alpha^n=a),
\]
is an isomorphism of groups.
\item $H^2(K,\mu_n)\simeq\operatorname{Br}(K)[n]$, the $n$-torsion of the \emph{Brauer group} $\operatorname{Br}(K):=H^2(K,\overline K^\times)$.
\item If $\mu_n\subset K$, then $H^1(K,\mu_n)=\Hom_{\mathrm{cont}}(G_K,\mu_n)$ and $\delta(a)$ is the Kummer character $\chi_a$ of \Cref{ch:galois}.
\item For a separable algebraic extension $L/K$, $\Res\circ\delta_K=\delta_L\circ(K^\times/K^{\times n}\to L^\times/L^{\times n})$. (Stated.) For $L/K$ finite,
$\Cor\circ\delta_L=\delta_K\circ\Nm_{L/K}$.
\end{enumerate}
\end{theorem}

\begin{proof}
(i) Exactness at the left and in the middle is clear. For $a\in\overline K^\times$ the polynomial $X^n-a$ has derivative $nX^{n-1}$, which is prime to $X^n-a$ because $n$
is invertible in $K$ and $a\ne0$; so $X^n-a$ is separable and has a root in the separable closure $\overline K$, and $x\mapsto x^n$ is surjective. The long exact sequence
(\Cref{ch:cohomology}) begins
\begin{multline*}
1\to\mu_n(K)\to K^\times\xrightarrow{\ n\ }K^\times\xrightarrow{\ \delta\ }H^1(K,\mu_n)\to H^1(K,\overline K^\times)\\
\xrightarrow{\ n\ }H^1(K,\overline K^\times)\to H^2(K,\mu_n)\to H^2(K,\overline K^\times)\xrightarrow{\ n\ }H^2(K,\overline K^\times).
\end{multline*}
By \Cref{thm:galcoh:hilbert90}(iii) the fifth term is zero, so $\delta$ induces an isomorphism $K^\times/(K^\times)^n\simeq H^1(K,\mu_n)$. By
\Cref{conv:galcoh:signs} the class $\delta(a)$ is represented by $\sigma\mapsto\sigma(\alpha)/\alpha$ for any lift $\alpha$ of $a$, that is, any $n$-th root of $a$ in
$\overline K$. (ii) The same sequence, with $H^1(K,\overline K^\times)=0$, identifies $H^2(K,\mu_n)$ with the kernel of multiplication by $n$ on $H^2(K,\overline K^\times)$.
(iii) If $\mu_n\subset K$ the action on $\mu_n$ is trivial, a cocycle is a continuous homomorphism, and coboundaries are zero; $\sigma\mapsto\sigma(\alpha)/\alpha$ is
the Kummer character $\chi_a$ of \Cref{ch:galois}. (iv) Choose $\overline L=\overline K$; for $\sigma\in G_L$ the restriction of the cocycle $\sigma\mapsto\sigma(\alpha)/\alpha$
is the cocycle that defines $\delta_L(a)$ with the same root $\alpha$. The corestriction formula is \cite[Chapter XIV, Section 1]{SerreLocalFields} and
\cite[Chapter VI, Section 2]{NSW}; we do not use it.
\end{proof}

\begin{remark}
The connecting map measures the failure of the choice of an $n$-th root to be Galois equivariant: a root $\alpha$ is defined only up to $\mu_n$, and
$\sigma(\alpha)/\alpha$ records how $\sigma$ moves it, much as a phase is defined only up to a constant and its variation is what is observed. Changing $\alpha$ to
$\zeta\alpha$ changes the cocycle by $\sigma\mapsto\sigma(\zeta)/\zeta$, a coboundary.
\end{remark}

\begin{definition}[Kummer map, Kummer class, Kummer cocycle]\label{def:galcoh:kummer-class}
The isomorphism $\delta$ of \Cref{thm:galcoh:kummer} is the \emph{Kummer map} of $K$ (for $n$); $\delta(a)\in H^1(K,\mu_n)$ is the \emph{Kummer class} of
$a\in K^\times$, and a cocycle $\sigma\mapsto\sigma(\alpha)/\alpha$ with $\alpha^n=a$ is a \emph{Kummer cocycle} of $a$. For $\ell\ne\operatorname{char}K$ and a compatible
system $(a^{1/\ell^k})_{k\ge0}$ of $\ell$-power roots of $a$ (that is, $(a^{1/\ell^{k+1}})^\ell=a^{1/\ell^k}$), the \emph{$\ell$-adic Kummer cocycle} is
\[
\kappa_a(\sigma)=\bigl(\sigma(a^{1/\ell^k})/a^{1/\ell^k}\bigr)_{k\ge0}\in\varprojlim_k\mu_{\ell^k}=\Z_\ell(1),
\]
the cocycle of \Cref{ch:galois,ch:local}.
\end{definition}

\begin{proposition}[properties of Kummer classes]\label{prop:galcoh:kummer-class}
Let $n$ be prime to $\operatorname{char}K$ and $a\in K^\times$.
\begin{enumerate}[label=(\roman*)]
\item For every separable algebraic extension $L/K$: $\Res_{G_L}\delta_K(a)=0$ if and only if $a\in(L^\times)^n$.
\item If $\mu_n\subset K$ and $\alpha^n=a$, the kernel of $\delta(a)=\chi_a$ is $G_{K(\alpha)}$, and $K(\alpha)/K$ is cyclic of degree equal to the order of $a$ in
$K^\times/(K^\times)^n$.
\item The reduction modulo $\ell^k$ of the $\ell$-adic Kummer cocycle $\kappa_a$, that is its image under $\Z_\ell(1)\to\mu_{\ell^k}$, is the Kummer cocycle of $a$ in
$H^1(K,\mu_{\ell^k})$ attached to the root $a^{1/\ell^k}$.
\end{enumerate}
\end{proposition}

\begin{proof}
(i) By \Cref{thm:galcoh:kummer}(iv) the restriction of $\delta_K(a)$ is $\delta_L(a)$, which is zero if and only if $a\in(L^\times)^n$ by \Cref{thm:galcoh:kummer}(i) for
$L$. (ii) $\chi_a(\sigma)=1$ if and only if $\sigma(\alpha)=\alpha$, so $\ker\chi_a=G_{K(\alpha)}$ and $\Gal(K(\alpha)/K)\simeq\chi_a(G_K)$, a subgroup of the cyclic group
$\mu_n$. Its order is the order of $\chi_a$ in $\Hom(G_K,\mu_n)$; since $\chi_a^m=\chi_{a^m}$ and $\delta$ is injective, $\chi_a^m=1$ if and only if $a^m\in(K^\times)^n$.
(iii) The projection $\Z_\ell(1)\to\mu_{\ell^k}$ takes the component of index $k$, which is $\sigma(a^{1/\ell^k})/a^{1/\ell^k}$.
\end{proof}

\begin{example}[Kummer theory over a finite field]\label{ex:galcoh:kummer-finite-field}
Let $K=\F_q$ with $q$ a power of a prime $p\ne\ell$, $v=v_\ell(q-1)$, and $\Frob_q$ the generator $x\mapsto x^q$ of $G_K\simeq\hat\Z$. Two computations of
$H^1(\F_q,\mu_{\ell^k})$ agree.
\begin{enumerate}[label=(\alph*)]
\item By Kummer theory it is $\F_q^\times/(\F_q^\times)^{\ell^k}$, cyclic of order $\ell^{\min(k,v)}$, $\F_q^\times$ being cyclic of order $q-1$ (\Cref{ch:rings}).
\item By the procyclic computation of \Cref{ch:cohomology} ($G_K\simeq\hat\Z$, finite coefficients), $c\mapsto c(\Frob_q)$ identifies it with
$\mu_{\ell^k}/(\Frob_q-1)\mu_{\ell^k}=\mu_{\ell^k}/\mu_{\ell^k}^{q-1}$, of order $\gcd(\ell^k,q-1)=\ell^{\min(k,v)}$.
\end{enumerate}
The composite isomorphism sends $u\in\F_q^\times$ to $\delta(u)(\Frob_q)=\Frob_q(\alpha)/\alpha=\alpha^{q-1}$, where $\alpha\in\overline\F_q$ and $\alpha^{\ell^k}=u$. For $k\le v$
this is $(\alpha^{\ell^k})^{(q-1)/\ell^k}=u^{(q-1)/\ell^k}$, the map of \Cref{ch:galois} (its Kummer proposition over finite fields). Passing to the limit over $k$
(\Cref{thm:galcoh:kummer-ladic} below, or directly the procyclic computation with $\Frob_q$ acting on $\Z_\ell(1)$ by $q$):
\[
H^1(\F_q,\Z_\ell(1))\;=\;\Z_\ell(1)/(q-1)\Z_\ell(1)\;\simeq\;\F_q^\times\otimes\Z_\ell=\F_q^\times/(\F_q^\times)^{\ell^v}\;\simeq\;\Z/\ell^v .
\]
The first equality is $c\mapsto c(\Frob_q)$; the second isomorphism is the Kummer isomorphism, canonical; the third depends on the choice of a generator. The
class of $u$ has valuation at least $k$ (for $k\le v$) if and only if $u\in(\F_q^\times)^{\ell^k}$, where, as in Lemma lem:tamfactor, a class $c\in\Z/\ell^v$ has
$v_\ell(c)=\min(v_\ell(\tilde c),v)$ for any lift $\tilde c\in\Z$; in particular $v_\ell(0)=v$.
\end{example}

\begin{theorem}[Brauer groups]\label{thm:galcoh:brauer}
\begin{enumerate}[label=(\roman*)]
\item $\operatorname{Br}(\F_q)=0$, $\operatorname{Br}(\C)=0$ and $\operatorname{Br}(\R)=\Z/2$.
\item Let $K$ be a finite extension of $\Q_p$. There is a canonical isomorphism $\operatorname{inv}_K\colon\operatorname{Br}(K)\to\Q/\Z$, the \emph{invariant}. For a finite
extension $L/K$, $\operatorname{inv}_L\circ\Res=[L:K]\cdot\operatorname{inv}_K$ and $\operatorname{inv}_K\circ\Cor=\operatorname{inv}_L$. With \Cref{thm:galcoh:kummer}(ii),
$H^2(K,\mu_n)\simeq\Z/n$ canonically, for every $n\ge1$.
\item (Albert, Brauer, Hasse, Noether.) The sequence
\[
0\to\operatorname{Br}(\Q)\to\bigoplus_v\operatorname{Br}(\Q_v)\xrightarrow{\ \sum_v\operatorname{inv}_v\ }\Q/\Z\to0
\]
is exact, the sum over all places $v$ of $\Q$, with $\operatorname{inv}_\infty\colon\operatorname{Br}(\R)\simeq\tfrac12\Z/\Z$. In particular an element of $\operatorname{Br}(\Q)$
has only finitely many nonzero local components, and the sum of its local invariants is $0$.
\end{enumerate}
\end{theorem}

\noindent We do not prove this; see \cite[Chapters X, XII and XIII]{SerreLocalFields} for (i) and (ii), Serre's and Tate's chapters in
\cite[Chapters VI and VII]{CasselsFrohlich} for (ii) and (iii), and \cite[Chapter II, Section 5]{SerreGaloisCoh}. Part (ii) is the cohomological form of local class
field theory (\Cref{ch:local}); it provides the target $H^2(K,\mu)\simeq\Q/\Z$ of the local pairing of \Cref{sec:galcoh:local}, and part (iii) is the reciprocity law behind
\Cref{prop:galcoh:reciprocity}.

\begin{paperlink}\label{pl:galcoh:kummer}
The proof of Lemma lem:tamfactor writes the Galois action on $T_\ell E$ through matrices ``with $\varepsilon$ unramified of order at most $2$ and $\kappa$ the Kummer
cocycle of $q_E$ in $H^1(G,\Z_\ell(1))=\varprojlim\Q_q^\times/(\Q_q^\times)^{\ell^k}$''. The equality is \Cref{thm:galcoh:kummer-ladic} below for $K=\Q_q$, which
rests on \Cref{thm:galcoh:kummer} at each finite level; the class of $\kappa$ is the Kummer class of $q_E$ in the sense of \Cref{def:galcoh:kummer-class}, and the
cocycle itself is the $\ell$-adic Kummer cocycle of the same definition. The proof of Lemma lem:tamagawa-exp argues ``By Kummer theory over
$\Q_q^{\mathrm{ur}}$, which contains $\mu_{\ell^t}$ as $q\ne\ell$''; that is \Cref{prop:galcoh:kummer-class}(i), (iii) with $L=\Q_q^{\mathrm{ur}}$, and
\Cref{rem:galcoh:tamagawa-exp-cohomological} gives the cohomological reading.
\end{paperlink}
\section{\texorpdfstring{$\ell$}{l}-adic coefficients: \texorpdfstring{$T$}{T}, \texorpdfstring{$V$}{V}, \texorpdfstring{$W$}{W} and Kummer duals}\label{sec:galcoh:ladic}

The paper works with three modules built from one lattice: $T^D$, $V^D$ and $W^D$ in Section 4.2 and Lemma lem:dfg21, and $W=\ad T\otimes\Q_\ell/\Z_\ell$ in
Theorem thm:main. This section defines the Kummer duals in general, proves the exact relation between the cohomology of $T$, $V$ and $W$ (in particular that the
torsion of $H^1(G,T)$ comes from $H^0(G,W)$, the fact behind ``$H^1_f(\Q_p,T^D)$ is the torsion of $H^1(\Q_p,T^D)$, a finite group'' after Lemma lem:dfg21), and
computes $H^1(K,\Z_\ell(1))$, the group in which the Kummer class of the Tate parameter lives in Lemma lem:tamfactor.

\begin{convention}[lattices and duals]\label{conv:galcoh:duals}
$T$ is a free $\Z_\ell$-module of finite rank with a continuous $\Z_\ell$-linear action of a profinite group $G$; $V=T\otimes_{\Z_\ell}\Q_\ell$ and
$W=T\otimes_{\Z_\ell}\Q_\ell/\Z_\ell=V/T$, so that $W[\ell^k]=\ell^{-k}T/T\simeq T/\ell^kT$ (\Cref{ch:modules}). The maps induced by $T\subset V\to W$ are written
\[
\iota\colon H^n(G,T)\to H^n(G,V),\qquad\pi\colon H^n(G,V)\to H^n(G,W).
\]
$A^*$ is the \emph{Pontryagin dual}: $\Hom(A,\Q/\Z)$ for a finite group, $\Hom_{\mathrm{cont}}(A,\Q_\ell/\Z_\ell)$ for a finitely generated or a discrete $\ell$-primary
$\Z_\ell$-module (\Cref{ch:modules}); this is the paper's ${}^*$ in Lemma lem:dfg21. $T^\vee=\Hom_{\Z_\ell}(T,\Z_\ell)$ and $V^\vee$ are linear duals. $M(1)=M\otimes_{\Z_\ell}\Z_\ell(1)$.
``Length'' means $\Z_\ell$-length, $\len=\len_{\Z_\ell}$, and $\#X=\ell^{\len X}$ for a finite $\Z_\ell$-module $X$.
\end{convention}

\begin{definition}[Kummer duals]\label{def:galcoh:kummer-duals}
Let $\mu$ be the group of roots of unity in $\overline K$ of order prime to $\operatorname{char}K$. For a finite
$G$-module $M$ of order prime to $\operatorname{char}K$, the \emph{Kummer dual} is $M^D=\Hom(M,\mu)$; if $M$ has $\ell$-power order, $M^D=\Hom(M,\mu_{\ell^\infty})=M^*(1)$. For
$T$ as in \Cref{conv:galcoh:duals}:
\[
T^D:=\Hom_{\Z_\ell}(T,\Z_\ell(1))=T^\vee(1),\qquad V^D:=T^D\otimes\Q_\ell=\Hom_{\Q_\ell}(V,\Q_\ell(1)),\qquad W^D:=V^D/T^D .
\]
In all cases $G$ acts on homomorphisms by $(g\phi)(x)=g\,\phi(g^{-1}x)$.
\end{definition}

These are the paper's $T^D$, $V^D$, $W^D$ of Section 4.2, where $T$ is the adjoint lattice $A=\ad T_\ell E$; the identification $T^D\simeq A(1)$ through the trace form
is made in \Cref{ch:blochkato}.

\begin{proposition}[Kummer duality for lattices]\label{prop:galcoh:kummer-duals}
\begin{enumerate}[label=(\roman*)]
\item The evaluation pairings
\begin{gather*}
T\times T^D\to\Z_\ell(1),\qquad V\times V^D\to\Q_\ell(1),\qquad M\times M^D\to\mu,\\
W\times T^D\to\Q_\ell/\Z_\ell(1),\qquad T\times W^D\to\Q_\ell/\Z_\ell(1)
\end{gather*}
are $G$-equivariant and perfect: each module is identified with the continuous homomorphisms from the other into the target.
\item $W^D\simeq\Hom(T,\Q_\ell/\Z_\ell(1))=T^*(1)$, and for each $k$, $W[\ell^k]^D=(T/\ell^kT)^D\simeq T^D/\ell^kT^D\simeq W^D[\ell^k]$.
\item $(T^D)^D=T$ canonically; $W^D$ is the Kummer dual of $T$ and $T^D$ that of $W$ in the sense of (i).
\item $\Z_\ell^D=\Z_\ell(1)$ and $\Z_\ell(1)^D=\Z_\ell$; for $T=\Z_\ell$, $W^D=\mu_{\ell^\infty}$, and the Kummer dual of $W=\Q_\ell/\Z_\ell$ in the sense of (iii) is $\Z_\ell(1)$. For an elliptic curve $E/K$ and $n$ prime to
$\operatorname{char}K$, the Weil pairing gives $E[n]^D\simeq E[n]$ and $(T_\ell E)^D\simeq T_\ell E$.
\end{enumerate}
\end{proposition}

\begin{proof}
(i) Equivariance: $\langle gx,g\phi\rangle=(g\phi)(gx)=g\,\phi(g^{-1}gx)=g\langle x,\phi\rangle$. Perfectness of $T\times T^D\to\Z_\ell(1)$ is linear algebra over $\Z_\ell$ for
free modules of finite rank ($\Z_\ell(1)$ is free of rank one), and tensoring with $\Q_\ell$ gives the second pairing. For the third, a choice of basis $e_1,\dots,e_r$ of $T$
and a generator of $\Z_\ell(1)$ identifies $T^D$ with $\Z_\ell^r$ by the dual basis and $W$ with $(\Q_\ell/\Z_\ell)^r$, and the pairing becomes
$((x_i),(y_i))\mapsto\sum x_iy_i\in\Q_\ell/\Z_\ell$; that this identifies each side with the continuous dual of the other is the duality between $\Z_\ell$ and
$\Q_\ell/\Z_\ell$ of \Cref{ch:modules}. The fourth pairing is the third for $T^D$ in place of $T$, by (iii). For finite $M$ the pairing $M\times\Hom(M,\mu)\to\mu$ is
perfect because $\mu$ contains a cyclic group of order $\exp(M)$ and $\#\Hom(M,\mu)=\#M$.
(ii) For $T$ free, $\Hom(T,\Z_\ell(1))\otimes\Q_\ell/\Z_\ell=\Hom(T,\Z_\ell(1)\otimes\Q_\ell/\Z_\ell)$, and $V^D/T^D=T^D\otimes\Q_\ell/\Z_\ell$. Similarly
$\Hom(T/\ell^kT,\mu_{\ell^k})=\Hom(T,\Z_\ell(1))\otimes\Z/\ell^k$, and $W^D[\ell^k]=\ell^{-k}T^D/T^D\simeq T^D/\ell^kT^D$.
(iii) The map $T\to\Hom(\Hom(T,\Z_\ell(1)),\Z_\ell(1))$, $x\mapsto(\phi\mapsto\phi(x))$, is an isomorphism of free modules of the same rank (check on a basis) and is
equivariant by (i). The rest is a restatement of (i).
(iv) The first statements are immediate from the definitions. The Weil pairing $e_n\colon E[n]\times E[n]\to\mu_n$ is bilinear, alternating, perfect and
Galois equivariant (\Cref{ch:ellcurves}), so $P\mapsto e_n(P,\cdot)$ is an isomorphism $E[n]\to E[n]^D$; the pairings for $n=\ell^k$ are compatible and give
$T_\ell E\simeq\Hom(T_\ell E,\Z_\ell(1))$.
\end{proof}

\begin{proposition}[the cohomology of $T$, $V$, $W$]\label{prop:galcoh:tvw}
Let $G$ be a profinite group and $T$ as in \Cref{conv:galcoh:duals}, and suppose that $H^1(G,T/\ell T)$ and $H^2(G,T/\ell T)$ are finite. (This holds for $G=G_{\F_q}$,
$G_\R$, $G_{\Q_p}$ and $G_{\Q,S}$: \Cref{ch:cohomology}, \Cref{thm:galcoh:local-finiteness} and \Cref{thm:galcoh:gqs-finite}.)
\begin{enumerate}[label=(\roman*)]
\item $H^n(G,T)$ is a finitely generated $\Z_\ell$-module for $n\le2$, and $H^n(G,V)=H^n(G,T)\otimes_{\Z_\ell}\Q_\ell$, so that $\iota$ is $x\mapsto x\otimes1$.
\item (Torsion of $H^1(G,T)$.) The sequence $H^0(G,V)\to H^0(G,W)\xrightarrow{\ \delta\ }H^1(G,T)\xrightarrow{\ \iota\ }H^1(G,V)$ is exact, $\ker\iota=H^1(G,T)_{\tors}$, and the
image of $H^0(G,V)$ in $H^0(G,W)$ is $H^0(G,T)\otimes\Q_\ell/\Z_\ell$, which is the maximal divisible subgroup $H^0(G,W)_{\mathrm{div}}$. Hence
\[
H^1(G,T)_{\tors}\;\simeq\;H^0(G,W)/H^0(G,W)_{\mathrm{div}},
\]
a finite group. If $H^0(G,V)=0$, then $H^0(G,W)$ is finite and $\delta\colon H^0(G,W)\to H^1(G,T)_{\tors}$ is an isomorphism.
\item (Divisible part of $H^1(G,W)$.) The sequence $0\to H^1(G,T)\otimes\Q_\ell/\Z_\ell\to H^1(G,W)\to H^2(G,T)_{\tors}\to0$ is exact; the first term is the image of
$\pi$ and equals $H^1(G,W)_{\mathrm{div}}$. So $H^1(G,W)$ is cofinitely generated, of corank $\rank_{\Z_\ell}H^1(G,T)=\dim_{\Q_\ell}H^1(G,V)$.
\item (Finite levels.) For $n\ge1$ and $k\ge1$ the sequences
\begin{gather*}
0\to H^{n-1}(G,T)/\ell^k\to H^{n-1}(G,T/\ell^kT)\to H^n(G,T)[\ell^k]\to0,\\
0\to H^{n-1}(G,W)/\ell^k\to H^n(G,W[\ell^k])\to H^n(G,W)[\ell^k]\to0
\end{gather*}
are exact.
\end{enumerate}
\end{proposition}

\begin{proof}
(iv) The sequences $0\to T\xrightarrow{\ell^k}T\to T/\ell^kT\to0$ and $0\to W[\ell^k]\to W\xrightarrow{\ell^k}W\to0$ (exact because $W$ is divisible) admit continuous sections
(\Cref{ch:cohomology}), so they give long exact sequences, and the stated short exact sequences are the usual extraction of a long exact sequence at the
two maps given by multiplication by $\ell^k$.

(i) From the first sequence of (iv) with $k=1$ and induction on $k$ (the filtration $\ell^jT/\ell^kT$ has graded pieces $\simeq T/\ell T$), all $H^n(G,T/\ell^kT)$, $n\le2$, are finite.
By \Cref{ch:cohomology} (the theorem of Tate on limits) $H^n(G,T)=\varprojlim_kH^n(G,T/\ell^kT)$ for $n\le2$, a compact pro-$\ell$ $\Z_\ell$-module, and by (iv) its
quotient by $\ell$ embeds in the finite group $H^n(G,T/\ell T)$. A compact pro-$\ell$ $\Z_\ell$-module $X$ with $X/\ell X$ finite is finitely generated, by the topological form of
Nakayama's lemma (\Cref{ch:modules}; \cite[Appendix B]{Rubin}).
The equality $H^n(G,V)=H^n(G,T)\otimes\Q_\ell$ is \Cref{ch:cohomology}: a continuous cochain with values in $V$ has compact image, hence takes values in $\ell^{-a}T$ for
some $a$.

(ii) The sequence $0\to T\to V\to W\to0$ has a continuous section (\Cref{ch:cohomology}), whence the exact sequence
$H^0(G,T)\to H^0(G,V)\to H^0(G,W)\xrightarrow{\delta}H^1(G,T)\xrightarrow{\iota}H^1(G,V)$. By (i), $\iota$ is $X\to X\otimes\Q_\ell$ for the finitely generated $\Z_\ell$-module
$X=H^1(G,T)$, whose kernel is $X_{\tors}$; so $\delta(H^0(G,W))=H^1(G,T)_{\tors}$, a finite group. Next, $T^G=T\cap V^G$, and $V^G=T^G\otimes\Q_\ell$ (if $x\in V^G$
then $\ell^ax\in T\cap V^G$ for some $a$). So the image of $V^G$ in $W=V/T$ is $V^G/(V^G\cap T)=V^G/T^G\simeq T^G\otimes\Q_\ell/\Z_\ell$, a divisible group; call it $D$.
Then $D\subseteq(W^G)_{\mathrm{div}}$, and $(W^G)_{\mathrm{div}}/D$ is divisible (a quotient of a divisible group) and finite (a subgroup of $W^G/D\simeq H^1(G,T)_{\tors}$),
hence zero, since a finite divisible group is trivial. This proves $D=(W^G)_{\mathrm{div}}$ and the displayed isomorphism. If $V^G=0$, then $D=0$ and $W^G\simeq H^1(G,T)_{\tors}$
is finite.

(iii) The continuation of the same long exact sequence reads $H^1(G,T)\xrightarrow{\iota}H^1(G,V)\xrightarrow{\pi}H^1(G,W)\to H^2(G,T)\to H^2(G,V)$. The image of $\pi$
is $\coker\iota=(X\otimes\Q_\ell)/\iota(X)\simeq X\otimes\Q_\ell/\Z_\ell$ (for $X$ finitely generated, $X\otimes\Q_\ell/\Z_\ell=(X/X_{\tors})\otimes\Q_\ell/\Z_\ell$), and the cokernel of $\pi$ is
$\ker(H^2(G,T)\to H^2(G,T)\otimes\Q_\ell)=H^2(G,T)_{\tors}$, finite by (i). So $H^1(G,W)$ is an extension of a finite group by $(\Q_\ell/\Z_\ell)^r$, $r=\rank X$, which is
cofinitely generated of corank $r$; the image of $\pi$ is divisible of finite index, and the argument at the end of (ii) shows that it is the maximal divisible
subgroup.
\end{proof}

\begin{warning}\label{warn:galcoh:torsion}
$H^1(G,V)$ is a $\Q_\ell$-vector space and sees no torsion: two lattices with the same $V$ can have different $H^1(G,T)_{\tors}$, and it is $H^0(G,W)$, not $H^0(G,V)$,
that controls it. The torsion is $H^0(G,W)$ modulo its divisible part, not $H^0(G,W)$ itself; the two agree when $H^0(G,V)=0$, which is the case in every application
in the paper (\Cref{cor:galcoh:ur-vanishing,thm:galcoh:relaxing-sigma}).
\end{warning}

\begin{example}[checks of \Cref{prop:galcoh:tvw}]\label{ex:galcoh:tvw-examples}
\begin{enumerate}[label=(\alph*)]
\item $G=G_{\F_q}$, $q$ prime to $\ell$, $T=\Z_\ell(1)$: $\Frob_q$ acts on $V=\Q_\ell(1)$ by $q\ne1$, so $H^0(G,V)=0$; $H^0(G,W)=\mu_{\ell^\infty}(\F_q)$ has order $\ell^{v_\ell(q-1)}$; and
$H^1(G,T)=\Z_\ell(1)/(q-1)$ (\Cref{ex:galcoh:kummer-finite-field}) is all torsion, isomorphic to $H^0(G,W)$, as (ii) predicts.
\item $G=G_\R$, $\ell$ odd, any $T$: $H^n(\R,T)=0$ for $n\ge1$ (\Cref{prop:galcoh:real-place}), and $H^0(\R,W)=W^{c=1}=T^{c=1}\otimes\Q_\ell/\Z_\ell$ is divisible, so
$H^1(\R,T)_{\tors}=0$ is consistent with (ii).
\item $G=G_{\Q_p}$, $p\ne\ell$, $T=\Z_\ell(1)$: by \Cref{prop:galcoh:kummer-local}, $H^1(\Q_p,\Z_\ell(1))_{\tors}$ is the image of $\Z_p^\times\otimes\Z_\ell$, of order
$\ell^{v_\ell(p-1)}=\#\mu_{\ell^\infty}(\Q_p)=\#H^0(\Q_p,\Q_\ell/\Z_\ell(1))$.
\end{enumerate}
\end{example}

\begin{theorem}[$\ell$-adic Kummer theory]\label{thm:galcoh:kummer-ladic}
Let $K$ be a field of characteristic different from $\ell$, and $\widehat{K^\times}:=\varprojlim_kK^\times/(K^\times)^{\ell^k}$ (transition maps the projections).
\begin{enumerate}[label=(\roman*)]
\item The Kummer maps induce an isomorphism $\widehat{K^\times}\simeq H^1(K,\Z_\ell(1))$; the class of $a\in K^\times$ is $\kappa_a:=(\delta_k(a))_k$, represented by the
$\ell$-adic Kummer cocycle of \Cref{def:galcoh:kummer-class} for any compatible system of roots.
\item The Kummer maps induce an isomorphism $K^\times\otimes\Q_\ell/\Z_\ell\simeq H^1(K,\Q_\ell/\Z_\ell(1))$.
\item $\kappa_{ab}=\kappa_a+\kappa_b$; $\kappa_a=0$ if and only if $a\in\bigcap_k(K^\times)^{\ell^k}$; for a separable algebraic extension $L/K$ the restriction of $\kappa_a$ to
$G_L$ is $\kappa_a$ computed over $L$.
\end{enumerate}
\end{theorem}

\begin{proof}
(i) The groups $H^0(K,\mu_{\ell^k})$ are finite, so by \Cref{ch:cohomology} (Tate's theorem on limits, the case $n=1$)
$H^1(K,\Z_\ell(1))=\varprojlim_kH^1(K,\mu_{\ell^k})$, with transition maps induced by $\mu_{\ell^{k+1}}\to\mu_{\ell^k}$, $\zeta\mapsto\zeta^\ell$. On the Kummer side: if
$\alpha^{\ell^{k+1}}=a$, the image of the cocycle $\sigma\mapsto\sigma(\alpha)/\alpha$ is $\sigma\mapsto\sigma(\alpha^\ell)/\alpha^\ell$, the Kummer cocycle of $a$ at level $k$ for the root
$\beta=\alpha^\ell$. So the transition maps correspond to the projections $K^\times/(K^\times)^{\ell^{k+1}}\to K^\times/(K^\times)^{\ell^k}$, and \Cref{thm:galcoh:kummer}(i) at
each level gives (i). That $(\delta_k(a))_k$ is represented by $\kappa_a$ is \Cref{prop:galcoh:kummer-class}(iii).
(ii) Cohomology of discrete modules commutes with direct limits (by \eqref{eq:galcoh:direct-limit}, a cochain of a profinite group with values in a direct limit
factors through a finite quotient and takes finitely many values), so $H^1(K,\mu_{\ell^\infty})=\varinjlim_kH^1(K,\mu_{\ell^k})$ along the inclusions
$\mu_{\ell^k}\subset\mu_{\ell^{k+1}}$. If $\alpha^{\ell^k}=a$, then $\alpha^{\ell^{k+1}}=a^\ell$ and the cocycle $\sigma\mapsto\sigma(\alpha)/\alpha$ is also the Kummer cocycle of $a^\ell$
at level $k+1$; so the transition maps are $a\mapsto a^\ell$, and $\varinjlim(K^\times/(K^\times)^{\ell^k},a\mapsto a^\ell)=\varinjlim K^\times\otimes\Z/\ell^k=K^\times\otimes\Q_\ell/\Z_\ell$.
(iii) The first two statements follow from (i), since $\widehat{K^\times}$ is the $\ell$-adic completion and $a$ maps to $0$ if and only if it lies in every $(K^\times)^{\ell^k}$; the last
is \Cref{thm:galcoh:kummer}(iv) at each level.
\end{proof}

\begin{proposition}[local Kummer theory, $\ell\ne p$]\label{prop:galcoh:kummer-local}
Let $K$ be a finite extension of $\Q_p$, $\ell\ne p$, with ring of integers $\mathcal O_K$, uniformiser $\varpi$, residue field $\F_q$ and normalised valuation $v_K$; put
$v=v_\ell(q-1)$.
\begin{enumerate}[label=(\roman*)]
\item The sequence $0\to\mathcal O_K^\times\otimes\Z_\ell\to\widehat{K^\times}\xrightarrow{\ v_K\ }\Z_\ell\to0$ is exact and split (by $\varpi$), and
$\mathcal O_K^\times\otimes\Z_\ell\simeq\F_q^\times\otimes\Z_\ell\simeq\Z/\ell^v$. Hence $H^1(K,\Z_\ell(1))\simeq\Z_\ell\times\Z/\ell^v$, with torsion subgroup the classes of units.
\item $\#H^1(K,\mu_{\ell^k})=\ell^k\cdot\#\mu_{\ell^k}(K)=\ell^{k+\min(k,v)}$.
\item $H^1(K,\Q_\ell/\Z_\ell(1))\simeq K^\times\otimes\Q_\ell/\Z_\ell\simeq\Q_\ell/\Z_\ell$, through $v_K$.
\item $H^1(K,\Q_\ell(1))\simeq\widehat{K^\times}\otimes\Q_\ell\simeq\Q_\ell$, spanned by the class of $\varpi$.
\end{enumerate}
\end{proposition}

\begin{proof}
(i) By \Cref{ch:local}, $K^\times=\varpi^{\Z}\times\mu_{q-1}\times U^{(1)}$ with $U^{(1)}=1+\varpi\mathcal O_K$ a pro-$p$ group, on which raising to an $\ell^k$-th power is
bijective. So $K^\times/(K^\times)^{\ell^k}\simeq\Z/\ell^k\times\mu_{q-1}/\mu_{q-1}^{\ell^k}$, the first factor given by $v_K$, and the second has order $\ell^{\min(k,v)}$; the limit
is $\Z_\ell\times\Z/\ell^v$, and the second factor is $\mathcal O_K^\times\otimes\Z_\ell\simeq\F_q^\times\otimes\Z_\ell$. Apply \Cref{thm:galcoh:kummer-ladic}(i).
(ii) is the computation at level $k$, with \Cref{thm:galcoh:kummer}(i). (iii) In $K^\times\otimes\Q_\ell/\Z_\ell$ the finite group $\mu_{q-1}$ and the uniquely $\ell$-divisible
group $U^{(1)}$ contribute nothing, and $\varpi^{\Z}\otimes\Q_\ell/\Z_\ell=\Q_\ell/\Z_\ell$. (iv) By \Cref{ch:cohomology}, $H^1(K,\Q_\ell(1))=H^1(K,\Z_\ell(1))\otimes\Q_\ell$, and the
torsion dies.
\end{proof}

\noindent Numerical instances (with $v_3(6)=1$, $v_3(18)=2$, $v_5(10)=1$, $v_3(4)=0$): $H^1(\Q_7,\Z_3(1))\simeq\Z_3\times\Z/3$; $H^1(\Q_{19},\Z_3(1))\simeq\Z_3\times\Z/9$;
$H^1(\Q_{11},\Z_5(1))\simeq\Z_5\times\Z/5$; $H^1(\Q_5,\Z_3(1))\simeq\Z_3$, every unit of $\Z_5$ being a cube.
\section[Local Galois cohomology]{Local Galois cohomology: finiteness, Euler characteristic, Tate local duality}\label{sec:galcoh:local}

The three theorems of this section are the deep input of the chapter, and they are stated with references. Finiteness is behind the words ``a finite group''
in the paragraph after Lemma lem:dfg21; Tate local duality is what makes the dual terms $H^1_f(\Q_p,T^D)^*$ of Lemma lem:dfg21 the right objects (through
\Cref{prop:galcoh:propagation}); and the Euler characteristic, combined with duality, gives the counting arguments of \Cref{sec:galcoh:unramified}. The
examples at the end make the statements concrete over $\Q_p$.

Throughout the section $K$ is a finite extension of $\Q_p$ (the paper needs only $K=\Q_p$), $M$ is a finite discrete $G_K$-module, and $M^D=\Hom(M,\mu)$ its Kummer dual
(\Cref{def:galcoh:kummer-duals}).

\begin{theorem}[finiteness; cohomological dimension]\label{thm:galcoh:local-finiteness}
\begin{enumerate}[label=(\roman*)]
\item $H^n(K,M)$ is finite for every $n$, and $H^n(K,M)=0$ for $n\ge3$.
\item Consequently, for $T$ as in \Cref{conv:galcoh:duals} with a continuous action of $G_K$, the groups $H^n(K,T)$ are finitely generated $\Z_\ell$-modules, the
$H^n(K,V)$ are finite-dimensional, the $H^n(K,W)$ are cofinitely generated, and \Cref{prop:galcoh:tvw} applies with $G=G_K$.
\end{enumerate}
\end{theorem}

\noindent We do not prove (i); see \cite[Chapter II, Section 5]{SerreGaloisCoh}, \cite[Chapter VII, Section 2]{NSW} and \cite[Chapter I, Section 2]{MilneADT}.
(ii) The hypothesis of \Cref{prop:galcoh:tvw} is (i) for $M=T/\ell T$, so for $n\le2$ the statements about $T$, $V$ and $H^1(K,W)$ are \Cref{prop:galcoh:tvw}(i), (iii).
For $n\ge3$: by (i) and induction on $k$, $H^n(K,T/\ell^kT)=0$, so $H^n(K,T)=\varprojlim_kH^n(K,T/\ell^kT)=0$ by Tate's limit theorem (\Cref{ch:cohomology};
its hypothesis, finiteness of the $H^{n-1}(K,T/\ell^kT)$, holds by (i)), hence $H^n(K,V)=0$; and $H^n(K,W)=\varinjlim_kH^n(K,W[\ell^k])=0$ since cohomology of
discrete modules commutes with direct limits. For $H^2(K,W)$: the sequence of \Cref{prop:galcoh:tvw}(iv) with $n=2$ shows $H^2(K,W)[\ell]$ is a quotient of the finite
group $H^2(K,W[\ell])$, and $H^2(K,W)$ is $\ell$-primary torsion, so it is cofinitely generated (\Cref{ch:modules}). Alternatively, $H^2(K,W)\simeq H^0(K,T^D)^*$
by \Cref{thm:galcoh:tate-local-duality}.

\begin{proof}[Sketch of proof of (i) for $H^1$ and $p\nmid\#M$]
This is the case used at the primes of $\Sigma$; at $p=\ell$ the theorem is used as stated. The action of $G_K$ on $M$ factors through $\Gal(L/K)$ for a finite Galois extension $L/K$, and inflation-restriction gives
$0\to H^1(\Gal(L/K),M)\to H^1(K,M)\to\Hom_{\mathrm{cont}}(G_L,M)$. The first group is finite (finite group, finite module). A continuous homomorphism $G_L\to M$ kills
the wild inertia group $P_L$, which is pro-$p$ while $M$ has order prime to $p$; so it factors through $\Gal(L^{\tame}/L)$, which is topologically generated by two
elements, a lift of Frobenius and a generator of tame inertia (\Cref{ch:local}). Such a homomorphism is determined by two elements of $M$, so there are finitely
many.
\end{proof}

\begin{theorem}[Tate's local Euler characteristic formula]\label{thm:galcoh:local-euler}
\[
\frac{\#H^0(K,M)\cdot\#H^2(K,M)}{\#H^1(K,M)}=\bigl|\#M\bigr|_K=p^{-[K:\Q_p]\,v_p(\#M)} .
\]
In particular the left side is $1$ when $p\nmid\#M$. For $V$ of type (V) over $\Q_\ell$: $\dim H^0(K,V)-\dim H^1(K,V)+\dim H^2(K,V)$ equals $-[K:\Q_p]\dim V$ if $\ell=p$,
and $0$ if $\ell\ne p$.
\end{theorem}

\noindent We do not prove this; see \cite[Chapter II, Section 5]{SerreGaloisCoh}, \cite[Chapter I, Section 2]{MilneADT} and \cite[Chapter VII, Section 3]{NSW}. The
$\ell$-adic form follows from the finite one applied to $T/\ell^kT$, using \Cref{prop:galcoh:tvw}(iv) to compare ranks.

\begin{theorem}[Tate local duality]\label{thm:galcoh:tate-local-duality}
\begin{enumerate}[label=(\roman*)]
\item For $0\le i\le2$ the cup product (\Cref{ch:cohomology}) followed by the evaluation $M\otimes M^D\to\mu$ and the invariant (\Cref{thm:galcoh:brauer}(ii),
through $H^2(K,\mu)=\varinjlim_nH^2(K,\mu_n)\simeq\Q/\Z$),
\[
H^i(K,M)\times H^{2-i}(K,M^D)\longrightarrow H^2(K,\mu)\xrightarrow{\ \operatorname{inv}_K\ }\Q/\Z,
\]
is a perfect pairing of finite groups. It is natural in $M$, compatible with restriction and corestriction, and compatible with connecting maps up to the sign
$(-1)^i$.
\item ($\ell$-adic forms.) For $T$ as in \Cref{conv:galcoh:duals} with a continuous action of $G_K$ and $0\le i\le2$, the cup products with the pairings of
\Cref{prop:galcoh:kummer-duals}(i) give perfect pairings
\begin{gather*}
H^i(K,W)\times H^{2-i}(K,T^D)\to\Q_\ell/\Z_\ell,\qquad H^i(K,T)\times H^{2-i}(K,W^D)\to\Q_\ell/\Z_\ell,\\
H^i(K,V)\times H^{2-i}(K,V^D)\to\Q_\ell ;
\end{gather*}
the first two identify the discrete group with the Pontryagin dual of the finitely generated one, the third is a perfect pairing of finite-dimensional
$\Q_\ell$-vector spaces.
\item (Compatibility.) For $x\in H^1(K,V)$ and $y\in H^1(K,T^D)$, $\langle\pi x,y\rangle_{W,T^D}=\langle x,\iota y\rangle_{V,V^D}\bmod\Z_\ell$; symmetrically with the roles of
$(T,V,W)$ and $(T^D,V^D,W^D)$ exchanged.
\end{enumerate}
\end{theorem}

\noindent We do not prove this; see \cite[Chapter II, Section 5]{SerreGaloisCoh}, \cite{TateDuality}, \cite[Chapter I, Section 2]{MilneADT} and
\cite[Chapter VII, Section 2]{NSW} for (i), and \cite[Section 1.4]{Rubin} for (ii) and (iii). Part (ii) is obtained from (i) applied to $T/\ell^kT$ and
$(T/\ell^kT)^D=T^D/\ell^kT^D$ (\Cref{prop:galcoh:kummer-duals}(ii)) by passing to the limit. Part (iii) is the naturality of the cup product for the
compatible pairings $V\times T^D\to\Q_\ell(1)$ and $W\times T^D\to\Q_\ell/\Z_\ell(1)$, together with the fact that $H^2(K,\Q_\ell(1))=\Q_\ell\to H^2(K,\Q_\ell/\Z_\ell(1))=\Q_\ell/\Z_\ell$ is the
projection. No sign matters for any order or length in this book, and we do not track signs.

\begin{corollary}[consequences of local duality]\label{cor:galcoh:local-duality-consequences}
\begin{enumerate}[label=(\roman*)]
\item $H^2(K,M)\simeq H^0(K,M^D)^*$. In particular $H^2(K,\mu_{\ell^k})\simeq\Z/\ell^k$ and $H^2(K,\Z/\ell^k)\simeq\mu_{\ell^k}(K)^*$.
\item If $p\nmid\#M$, then $\#H^1(K,M)=\#H^0(K,M)\cdot\#H^0(K,M^D)$. If $\ell\ne p$, then $\dim H^1(K,V)=\dim H^0(K,V)+\dim H^0(K,V^D)$.
\item If $\ell\ne p$, then $H^2(K,T)\simeq H^0(K,W^D)^*$; it is finite if $H^0(K,V^D)=0$.
\end{enumerate}
\end{corollary}

\begin{proof}
(i) is \Cref{thm:galcoh:tate-local-duality}(i) with $i=2$; for $M=\mu_{\ell^k}$, $M^D=\Z/\ell^k$ with trivial action, and for $M=\Z/\ell^k$, $M^D=\mu_{\ell^k}$. (ii)
\Cref{thm:galcoh:local-euler} gives $\#H^1=\#H^0\cdot\#H^2$, and (i) gives $\#H^2(K,M)=\#H^0(K,M^D)$. The dimension formula is the same argument with the $\ell$-adic
forms of the two theorems. (iii) is \Cref{thm:galcoh:tate-local-duality}(ii) with $i=2$; if $H^0(K,V^D)=0$, then $H^0(K,W^D)$ is finite by \Cref{prop:galcoh:tvw}(ii)
applied to $T^D$.
\end{proof}

\noindent For instance, if $p\ne\ell$: $\dim H^1(\Q_p,\Q_\ell)=1+0=1$ and $\dim H^1(\Q_p,\Q_\ell(1))=0+1=1$ (compare \Cref{prop:galcoh:kummer-local}(iv)); and for an
elliptic curve $E/\Q_p$, $E[\ell]^D\simeq E[\ell]$ (\Cref{prop:galcoh:kummer-duals}(iv)) gives $\#H^1(\Q_p,E[\ell])=\#E(\Q_p)[\ell]^2$.

\begin{example}[orders of local cohomology groups]\label{ex:galcoh:local-orders}
\begin{enumerate}[label=(\alph*)]
\item Let $p\ne\ell$, $v=v_\ell(p-1)$, $k\ge1$ and $m=\min(k,v)$, so that $\#\mu_{\ell^k}(\Q_p)=\ell^m$. Then for $i=0,1,2$:
\[
\#H^i(\Q_p,\Z/\ell^k)=\ell^k,\ \ell^{k+m},\ \ell^m;\qquad\#H^i(\Q_p,\mu_{\ell^k})=\ell^m,\ \ell^{k+m},\ \ell^k .
\]
For $\mu_{\ell^k}$ this needs no duality: $H^0$ is $\mu_{\ell^k}(\Q_p)$, $H^1$ is \Cref{prop:galcoh:kummer-local}(ii), and $H^2\simeq\operatorname{Br}(\Q_p)[\ell^k]=\Z/\ell^k$ by
\Cref{thm:galcoh:kummer}(ii) and \Cref{thm:galcoh:brauer}(ii); the Euler characteristic $\ell^m\ell^k/\ell^{k+m}=1$ is visible. For $\Z/\ell^k$ the values come from
\Cref{cor:galcoh:local-duality-consequences}(i), (ii). For $(p,\ell,k)=(7,3,1)$: $3,9,3$ and $3,9,3$; for $(19,3,2)$: $9,81,9$ and $9,81,9$; for $(5,3,2)$, where $m=0$:
$9,9,1$ and $1,9,9$.
\item $p$ odd, $M=\mu_p$ over $\Q_p$: $H^0=\mu_p(\Q_p)=1$ (the roots of unity in $\Q_p$ have order dividing $p-1$); $H^1=\Q_p^\times/(\Q_p^\times)^p$ has order $p^2$,
since $\Q_p^\times=p^{\Z}\times\mu_{p-1}\times(1+p\Z_p)$ and $1+p\Z_p\simeq\Z_p$; $H^2=\operatorname{Br}(\Q_p)[p]\simeq\Z/p$. The Euler characteristic is $1\cdot p/p^2=p^{-1}=|p|_p$, as
\Cref{thm:galcoh:local-euler} predicts for $\#M=p$.
\item The curve \curve{11a1} over $\Q_{11}$, $\ell=5$. The reduction at $11$ is split multiplicative of type $\mathrm I_5$ ($a_{11}=1$), with Tate parameter
$q_E=11^5u$, $u\in\Z_{11}^\times$, $\bar u=10=-1$ (computed with PARI, \Cref{ex:galcoh:tate-kummer-examples}). As $-1=(-1)^5$, $\bar u$ is a fifth power in $\F_{11}^\times$,
so $u$ is a fifth power in $\Z_{11}^\times$ (Hensel), and $q_E=(11u^{1/5})^5$ is a fifth power in $\Q_{11}$; and $\mu_5\subset\Q_{11}$ as $11\equiv1\pmod5$. The
$5$-torsion of the Tate curve is generated by $\mu_5$ and a fifth root of $q_E$ (\Cref{ch:reduction}), so $E[5]\subset E(\Q_{11})$: all twelve roots of the
$5$-division polynomial lie in $\Q_{11}$, and $4x^3+b_2x^2+2b_4x+b_6$ is a square in $\Q_{11}$ at each of them (checked with \texttt{polrootspadic}). Hence
$\#H^1(\Q_{11},E[5])=\#E(\Q_{11})[5]^2=5^4$ by \Cref{cor:galcoh:local-duality-consequences}(ii). This is an Eisenstein example ($E(\Q)_{\tors}=\Z/5$), used for illustration
only.
\end{enumerate}
\end{example}
\section{Unramified cohomology}\label{sec:galcoh:unramified}

Unramified classes are the local conditions at every prime $p\ne\ell$ in the Selmer groups of the paper. This section computes them
($\Hur(\Q_p,M)=M^{I_p}/(\Frob_p-1)M^{I_p}$), proves that for $p\ne\ell$ the unramified subgroups of $M$ and of its Kummer dual are exact annihilators, explains how local
conditions on $V$ propagate to $T$ and $W$, and draws the consequence the paper states after Lemma lem:dfg21: at a prime $p\in\Sigma$ the propagated condition
on $T^D$ is the torsion of $H^1(\Q_p,T^D)$, a finite group, isomorphic to $H^0(\Q_p,W^D)$. The torsion form of the exact sequence of \Cref{prop:galcoh:ur-sequence}
is the starting point of the Tamagawa factors of \Cref{ch:blochkato}.

In this section $p$ is a prime, $G=G_{\Q_p}$, $I=I_p$, $\Frob$ an arithmetic Frobenius, and $M$ a $G$-module which is either discrete and torsion (type (D)), or of
type (C), or of type (V). The parts that need $p\ne\ell$ say so.

\begin{definition}[unramified classes]\label{def:galcoh:unramified}
Let $K$ be a finite extension of $\Q_p$, with inertia group $I_K\trianglelefteq G_K$, and $M$ a $G_K$-module of type (D), (C) or (V). The group of \emph{unramified classes} is
\[
\Hur(K,M):=\operatorname{im}\bigl(\Inf\colon H^1(G_K/I_K,M^{I_K})\to H^1(K,M)\bigr)=\ker\bigl(\Res\colon H^1(K,M)\to H^1(I_K,M)\bigr);
\]
for $K=\Q_p$, $G_K=G$ and $I_K=I$, this is $\Hur(\Q_p,M)$. The results below are stated for $\Q_p$, the case the paper uses; their proofs apply verbatim to $K$, with
$\Frob$ an arithmetic Frobenius of $K$ (acting by $x\mapsto x^{q}$ on the residue field $\F_q$ of $K$), $G_K/I_K\simeq\hat\Z$, and $|\cdot|_K$ in the Euler characteristic.
The two descriptions agree, and inflation is injective, by inflation-restriction (\Cref{ch:cohomology}, valid for the three types of coefficients); so
$\Hur(K,M)\simeq H^1(G_K/I_K,M^{I_K})$. The module $M$ is \emph{unramified} if $I_K$ acts trivially on it. The definition is made for every $p$, including $p=\ell$, where it is
\emph{not} the right local condition (\Cref{warn:galcoh:p-equals-l}; \Cref{ch:blochkato}).
\end{definition}

\begin{proposition}[computation of unramified cohomology]\label{prop:galcoh:unramified-computed}
\begin{enumerate}[label=(\roman*)]
\item $c\mapsto c(\Frob)$ induces an isomorphism $\Hur(\Q_p,M)\simeq M^I/(\Frob-1)M^I$.
\item If $M$ is finite, $\#\Hur(\Q_p,M)=\#H^0(\Q_p,M)$.
\item For $V$ of type (V), $\dim\Hur(\Q_p,V)=\dim H^0(\Q_p,V)$, and the following are equivalent: $\Hur(\Q_p,V)=0$; $H^0(\Q_p,V)=0$; $1$ is not an eigenvalue of
$\Frob$ on $V^I$; $\det(1-\Frob\mid V^I)\ne0$; $\det(1-\Phi\mid V^I)\ne0$.
\item Let $p\ne\ell$ and $T$ be a lattice with $T\otimes\Q_\ell=V$. Then $T^I=T\cap V^I$ is a lattice in $V^I$, saturated in $T$, and $\Hur(\Q_p,T)=T^I/(\Frob-1)T^I$. If the
conditions of (iii) hold, $\Hur(\Q_p,T)$ is finite of order
\[
\#\Hur(\Q_p,T)=\ell^{\,v_\ell(\det(1-\Frob\mid V^I))}=\ell^{\,v_\ell(\det(1-\Phi\mid V^I))}.
\]
\end{enumerate}
\end{proposition}

\begin{proof}
(i) $G/I\simeq\hat\Z$ is procyclic with topological generator the image of $\Frob$, and $M^I$ is a $G/I$-module of the same type as $M$ ($M^I$ is closed in $M$ for types
(C) and (V), being an intersection of kernels of continuous maps). By the procyclic computations of \Cref{ch:cohomology} (the discrete torsion case, and the compact
case, which applies because $\hat\Z$ surjects onto $\Z_\ell$), $c\mapsto c(\Frob)$ is an isomorphism $H^1(G/I,M^I)\simeq M^I/(\Frob-1)M^I$. Compose with the inverse of
inflation. (ii) For finite $M^I$ the endomorphism $\Frob-1$ has kernel and cokernel of the same order, and its kernel is $(M^I)^{\Frob=1}=M^G$. (iii) The same with
dimensions; the kernel of $\Frob-1$ on $V^I$ is $V^G$. The equivalences of the last four conditions are linear algebra, and on $V^I$, of dimension $d$,
\[
\det(1-\Phi)=\det\bigl(\Frob^{-1}(\Frob-1)\bigr)=(-1)^d\det(\Phi)\det(1-\Frob),
\]
so one determinant vanishes if and only if the other does. (iv) $T^I=T\cap V^I$ is clear, and it is saturated in $T$ (if $\ell x\in T^I$ with $x\in T$, then $x$ is $I$-fixed, $T$
being torsion-free); every $x\in V^I$ has $\ell^ax\in T\cap V^I$, so $T^I$ spans $V^I$. Part (i) for $M=T$ gives $\Hur(\Q_p,T)=T^I/(\Frob-1)T^I$. If $\det(1-\Frob\mid V^I)\ne0$,
$\Frob-1$ is an injective endomorphism of the free $\Z_\ell$-module $T^I$, and its cokernel has order $\ell^{v_\ell(\det(\Frob-1))}$ by the elementary divisor theorem
(\Cref{ch:modules}). Finally $\det(\Phi\mid T^I)$ is a unit of $\Z_\ell$ ($\Phi$ is an automorphism of the lattice $T^I$), so the two determinants have the same valuation by
the displayed identity.
\end{proof}

\begin{warning}[the Frobenius convention]\label{warn:galcoh:frobenius}
The \emph{vanishing} of $\Hur(\Q_p,V)$ and the \emph{order} of $\Hur(\Q_p,T)$ do not depend on whether one writes $\Frob$ or $\Phi$ (\Cref{prop:galcoh:unramified-computed}(iii),
(iv)). The \emph{value} of the determinant does. Only $\det(1-\Phi\mid(V^D)^{I_p})$ equals $L_p(A_f,1)^{-1}$ (\Cref{conv:galcoh:frobenius}; \Cref{ex:galcoh:ur-examples}(d)
shows the two values at a multiplicative prime). The revised paper writes the geometric Frobenius $\Frob_p^{-1}$ in this determinant, after Lemma lem:dfg21.
\end{warning}

\begin{proposition}[unramified classes and the action on $H^1(I,M)$]\label{prop:galcoh:ur-sequence}
Let $M$ be of type (D), or of type (C) or (V) with $p\ne\ell$. The sequence
\[
0\to\Hur(\Q_p,M)\to H^1(\Q_p,M)\xrightarrow{\ \Res\ }H^1(I,M)^{G/I}\to0
\]
is exact, where $G/I$ acts on $H^1(I,M)$ by the conjugation action $(\sigma\cdot c)(i)=\sigma\,c(\sigma^{-1}i\sigma)$ of \Cref{ch:cohomology}.
\end{proposition}

\begin{proof}
Exactness on the left and in the middle is \Cref{def:galcoh:unramified}; only the surjectivity of $\Res$ onto the invariants needs proof. The restriction lands in the
invariants because conjugation by an element of $G$ acts trivially on $H^1(G,M)$ (\Cref{ch:cohomology}).

\emph{Type (D).} The five-term exact sequence of inflation-restriction for $I\trianglelefteq G$ and discrete coefficients (\Cref{ch:cohomology}) reads
$0\to H^1(G/I,M^I)\to H^1(G,M)\to H^1(I,M)^{G/I}\to H^2(G/I,M^I)$, and the last group vanishes: $G/I\simeq\hat\Z$, and $H^2(\hat\Z,-)=0$ for discrete torsion modules
(\Cref{ch:cohomology}, the procyclic computations).

\emph{Type (C), $p\ne\ell$.} Put $M_k=M/\ell^kM$, a finite discrete module. By the case (D), each
$\Res_k\colon H^1(G,M_k)\to A_k:=H^1(I,M_k)^{G/I}$ is surjective. The groups $H^1(G,M_k)$ are finite (\Cref{thm:galcoh:local-finiteness}), and so are the $H^1(I,M_k)$: the
wild inertia group $P$ is pro-$p$ and $M_k$ has $\ell$-power order, so $H^1(P,M_k)=0$, inflation-restriction gives $H^1(I,M_k)=H^1(I/P,M_k^P)$, and $I/P$ is procyclic
(\Cref{ch:local}), so this is a quotient of $M_k^P$ (\Cref{ch:cohomology}). By Tate's limit theorem in degree $1$ (\Cref{ch:cohomology}, valid for every profinite group),
$H^1(G,M)=\varprojlim_kH^1(G,M_k)$ and $H^1(I,M)=\varprojlim_kH^1(I,M_k)$, compatibly with restriction and with the action of $G/I$; and invariants commute with inverse
limits, so $H^1(I,M)^{G/I}=\varprojlim_kA_k$. Let $(y_k)_k\in\varprojlim_kA_k$. The fibres $\Res_k^{-1}(y_k)$ are nonempty finite sets, and reduction maps
$\Res_{k+1}^{-1}(y_{k+1})$ into $\Res_k^{-1}(y_k)$; an inverse limit of nonempty finite sets is nonempty (compactness), so there is a compatible family $(x_k)_k$ with
$\Res_kx_k=y_k$, that is, an $x\in H^1(G,M)$ with $\Res x=(y_k)_k$.

\emph{Type (V), $p\ne\ell$.} Let $T\subset M$ be a $G$-stable lattice (\Cref{ch:cohomology}), so that $H^1(G,M)=H^1(G,T)\otimes\Q_\ell$ and $H^1(I,M)=H^1(I,T)\otimes\Q_\ell$
(\Cref{ch:cohomology}). The group $X=H^1(I,T)=\varprojlim_kH^1(I,T/\ell^kT)$ is compact, and $X/\ell X$ embeds in the finite group $H^1(I,T/\ell T)$ (\Cref{prop:galcoh:tvw}(iv),
which holds for any profinite group); so $X$ is finitely generated by topological Nakayama (\Cref{ch:modules}) and its torsion is killed by some $N\ge1$. Let
$y\in H^1(I,M)^{G/I}$ and write $y=\ell^{-a}\bar x$ with $x\in X$ and $\bar x$ its image. For $g\in G/I$, $gx-x$ maps to $0$ in $H^1(I,M)$, so it is torsion and
$g(Nx)=Nx$. By the case (C), $Nx=\Res z$ with $z\in H^1(G,T)$, and then $y=\Res\bigl((N\ell^a)^{-1}\iota z\bigr)$.
\end{proof}

\noindent \Cref{ch:cohomology} extends the five-term sequence to coefficients of types (C) and (V) when the quotient is $\hat\Z$, which gives the proposition at $p=\ell$ as
well; the paper needs it only at $p\ne\ell$, where the argument above is self-contained.

\begin{remark}\label{rem:galcoh:tamagawa-preview}
For $p\ne\ell$ and $M$ of type (C), $H^1(I,M)$ is computed through the tame quotient $I\to\Z_\ell(1)$ (\Cref{ch:cohomology}, the tame inertia example). Taking torsion
subgroups in \Cref{prop:galcoh:ur-sequence} for $M=T^D$ is the first step of the formula of Fontaine and Perrin-Riou, and the ideal
$\Fitt_{\Z_\ell}\bigl((H^1(I_p,M)_{\tors})^{G/I}\bigr)$ (the Fitting ideal \emph{of the invariants}) is the Tamagawa factor $\Tam_p(M)$ of Section 1.2 of the paper
(\Cref{ch:blochkato}).
\end{remark}

\begin{theorem}[exact orthogonality of unramified classes]\label{thm:galcoh:ur-orthogonal}
Let $p\ne\ell$.
\begin{enumerate}[label=(\roman*)]
\item For a finite $G$-module $M$ of $\ell$-power order, $\Hur(\Q_p,M)$ and $\Hur(\Q_p,M^D)$ are exact annihilators of each other under the pairing of
\Cref{thm:galcoh:tate-local-duality}(i).
\item For $V$ of type (V), $\Hur(\Q_p,V)$ and $\Hur(\Q_p,V^D)$ are exact annihilators under the pairing of \Cref{thm:galcoh:tate-local-duality}(ii).
\end{enumerate}
\end{theorem}

\begin{proof}
(i) \emph{Orthogonality.} Let $x=\Inf(a)$ and $y=\Inf(b)$, where $a\in H^1(G/I,M^I)$ and $b\in H^1(G/I,(M^D)^I)$. The cup product commutes with inflation
(\Cref{ch:cohomology}), so $x\cup y$ is the inflation of $a\cup b$, and its image under $M\otimes M^D\to\mu_{\ell^\infty}$ is the inflation of a class in
$H^2(G/I,\mu_{\ell^\infty}^I)$. That group is zero ($G/I\simeq\hat\Z$ and $\mu_{\ell^\infty}^I=\mu_{\ell^\infty}$ is discrete torsion), so $\langle x,y\rangle=0$.
\emph{Counting.} By \Cref{prop:galcoh:unramified-computed}(ii), \Cref{cor:galcoh:local-duality-consequences}(i) and \Cref{thm:galcoh:local-euler} (with $p\nmid\#M$),
\begin{align*}
\#\Hur(\Q_p,M)\cdot\#\Hur(\Q_p,M^D)&=\#H^0(\Q_p,M)\cdot\#H^0(\Q_p,M^D)\\
&=\#H^0(\Q_p,M)\cdot\#H^2(\Q_p,M)=\#H^1(\Q_p,M).
\end{align*}
As the pairing is perfect, the annihilator of $\Hur(\Q_p,M)$ has order $\#H^1(\Q_p,M)/\#\Hur(\Q_p,M)=\#\Hur(\Q_p,M^D)$; it contains $\Hur(\Q_p,M^D)$, so it equals it, and
symmetrically.
(ii) The orthogonality argument is the same, with $H^2(\hat\Z,\Q_\ell(1))=0$. By \Cref{prop:galcoh:unramified-computed}(iii) and
\Cref{cor:galcoh:local-duality-consequences}(ii), $\dim\Hur(\Q_p,V)+\dim\Hur(\Q_p,V^D)=\dim H^0(\Q_p,V)+\dim H^0(\Q_p,V^D)=\dim H^1(\Q_p,V)$, and a perfect pairing of
finite-dimensional spaces finishes the proof.
\end{proof}

\begin{warning}[the case $p=\ell$]\label{warn:galcoh:p-equals-l}
Over $\Q_p$, $p$ odd: $\#\Hur(\Q_p,\Z/p)=\#H^0(\Q_p,\Z/p)=p$ and $\#\Hur(\Q_p,\mu_p)=\#H^0(\Q_p,\mu_p)=1$ (\Cref{prop:galcoh:unramified-computed}(ii), which holds for every
$p$), while $\#H^1(\Q_p,\Z/p)=p\cdot\#H^2(\Q_p,\Z/p)\cdot p=p^2$ by \Cref{thm:galcoh:local-euler} with $\#H^2(\Q_p,\Z/p)=\#\mu_p(\Q_p)=1$. The annihilator of $\Hur(\Q_p,\Z/p)$ has
order $p$, not $1$: the unramified subgroups are not exact annihilators when $p=\ell$. (By local class field theory the annihilator of $\Hur(\Q_p,\Z/p)$ is the image of
$\Z_p^\times$ in $\Q_p^\times/(\Q_p^\times)^p=H^1(\Q_p,\mu_p)$; we do not use this.) At $\ell$ the right local condition is that of Bloch and Kato (\Cref{ch:blochkato}),
and the book never writes $\Hf=\Hur$ at $p=\ell$.
\end{warning}

The next proposition explains how a local condition on $V$ determines conditions on $T$ and on $W$, and why the dual of the quotient of $H^1(\Q_p,W)$ by the
propagated condition is the propagated condition on $T^D$; this is the mechanism that turns the paper's definition $H^1_f(\Q_p,T^D)=\iota^{-1}H^1_f(\Q_p,V^D)$ into the
dual terms of Lemma lem:dfg21.

\begin{proposition}[propagation of local conditions]\label{prop:galcoh:propagation}
Let $K$ be a finite extension of $\Q_p$ ($p=\ell$ allowed) and $T$ as in \Cref{conv:galcoh:duals}, with a continuous action of $G_K$. For a $\Q_\ell$-subspace
$L_V\subseteq H^1(K,V)$ put
\[
L_T:=\iota^{-1}(L_V)\subseteq H^1(K,T),\qquad L_W:=\pi(L_V)\subseteq H^1(K,W),
\]
the conditions \emph{propagated} from $L_V$; likewise $L_{T^D}$, $L_{W^D}$ from a subspace $L_{V^D}\subseteq H^1(K,V^D)$.
\begin{enumerate}[label=(\roman*)]
\item $L_T\supseteq H^1(K,T)_{\tors}$, and $L_W$ is divisible, so $L_W\subseteq H^1(K,W)_{\mathrm{div}}$.
\item If $L_V$ and $L_{V^D}$ are exact annihilators under the pairing $H^1(K,V)\times H^1(K,V^D)\to\Q_\ell$, then $L_W$ and $L_{T^D}$ are exact annihilators under
$H^1(K,W)\times H^1(K,T^D)\to\Q_\ell/\Z_\ell$; hence $\bigl(H^1(K,W)/L_W\bigr)^*\simeq L_{T^D}$. Symmetrically $L_T$ and $L_{W^D}$ are exact annihilators.
\item If $p\ne\ell$, $T$ is unramified and $L_V=\Hur(K,V)$, then $L_T=\Hur(K,T)$ and $L_W=\Hur(K,W)$.
\end{enumerate}
\end{proposition}

\begin{proof}
(i) $\ker\iota=H^1(K,T)_{\tors}$ by \Cref{prop:galcoh:tvw}(ii), and $\pi$ is $\Z_\ell$-linear on a $\Q_\ell$-vector space, so its image is divisible.
(ii) Let $y\in L_{T^D}$ and $x\in L_V$. By \Cref{thm:galcoh:tate-local-duality}(iii), $\langle\pi x,y\rangle=\langle x,\iota y\rangle\bmod\Z_\ell=0$, since $\iota y\in L_{V^D}$ is orthogonal
to $L_V$. So $L_{T^D}\subseteq L_W^\perp$. Conversely let $y\in H^1(K,T^D)$ be orthogonal to $L_W$. For $x\in L_V$ and $k\ge0$, $\ell^{-k}x\in L_V$, so
$\ell^{-k}\langle x,\iota y\rangle\in\Z_\ell$ for every $k$, whence $\langle x,\iota y\rangle=0$. Thus $\iota y\in L_V^\perp=L_{V^D}$ and $y\in L_{T^D}$. This proves
$L_W^\perp=L_{T^D}$. The pairing identifies the discrete group $H^1(K,W)$ with the Pontryagin dual of $H^1(K,T^D)$, and for a subgroup $A$ of a discrete group,
$(A^\perp)^\perp=A$ (Pontryagin duality, \Cref{ch:modules}); so $L_{T^D}^\perp=L_W$ as well, and $(H^1(K,W)/L_W)^*=L_W^\perp=L_{T^D}$. The symmetric statement is the same
argument with the second pairing of \Cref{thm:galcoh:tate-local-duality}(ii).
(iii) As $I$ acts trivially, $H^1(I,T)=\Hom_{\mathrm{cont}}(I,T)$ and $H^1(I,V)=\Hom_{\mathrm{cont}}(I,V)$, and $\Hom(I,T)\to\Hom(I,V)$ is injective, $T$ being torsion-free. So
for $x\in H^1(K,T)$, $\Res_I(\iota x)=0$ if and only if $\Res_Ix=0$: $L_T=\Hur(K,T)$. Inflation commutes with $\pi$, and $\pi\colon H^1(G/I,V)=V/(\Frob-1)V\to
H^1(G/I,W)=W/(\Frob-1)W$ is surjective because $V\to W$ is; so $\pi(\Hur(K,V))=\Hur(K,W)$.
\end{proof}

\begin{warning}
For $p\ne\ell$ and $T$ \emph{ramified} at $p$, the propagated condition $\iota^{-1}(\Hur(\Q_p,V))$ can be strictly larger than $\Hur(\Q_p,T)$: it contains all of
$H^1(\Q_p,T)_{\tors}$. The quotient is what the Tamagawa factor of \Cref{ch:blochkato} measures. The two coincide when $T$ is unramified at $p$
(\Cref{prop:galcoh:propagation}(iii)); they must not be identified in general.
\end{warning}

\begin{corollary}[vanishing unramified cohomology]\label{cor:galcoh:ur-vanishing}
Let $p\ne\ell$, let $T$ be a lattice with a continuous action of $G_{\Q_p}$, and suppose $\Hur(\Q_p,V)=0$. Then $H^0(\Q_p,V)=0$, and for the conditions propagated from
$L_V=\Hur(\Q_p,V)=0$:
\[
L_T=\iota^{-1}(0)=H^1(\Q_p,T)_{\tors}\simeq H^0(\Q_p,W),\qquad L_W=0 .
\]
$L_T$ is a finite group, and it contains the finite group $\Hur(\Q_p,T)$.
\end{corollary}

\begin{proof}
$H^0(\Q_p,V)=0$ by \Cref{prop:galcoh:unramified-computed}(iii). Then $L_T=\ker\iota=H^1(\Q_p,T)_{\tors}$, which is $\delta(H^0(\Q_p,W))\simeq H^0(\Q_p,W)$ and finite by
\Cref{prop:galcoh:tvw}(ii) and \Cref{thm:galcoh:local-finiteness}. $\Hur(\Q_p,T)$ is finite by \Cref{prop:galcoh:unramified-computed}(iv), hence torsion, hence contained in
$L_T$. $L_W=\pi(0)=0$.
\end{proof}

\begin{paperlink}\label{pl:galcoh:after-dfg21}
The paragraph after Lemma lem:dfg21 in the revised paper reads: ``For $p\in\Sigma$ the group $H^1_f(\Q_p,V^D)=H^1_{\mathrm{ur}}(\Q_p,V^D)=(V^D)^{I_p}/(\Frob_p-1)$
vanishes, because $\det(1-\Frob_p^{-1}\mid(V^D)^{I_p})=L_p(A_f,1)^{-1}\ne0$; for the same reason $H^0(\Q_p,V^D)=0$, \dots'', and, in the next sentence, ``So
$H^1_f(\Q_p,T^D)$ is the torsion of $H^1(\Q_p,T^D)$, a finite group.'' In the book's terms: the equality $H^1_f=\Hur$ at $p\ne\ell$ is \Cref{ch:blochkato}; the
equality with $(V^D)^{I_p}/(\Frob_p-1)$ is \Cref{prop:galcoh:unramified-computed}(i); the vanishing and the vanishing of $H^0(\Q_p,V^D)$ are
\Cref{prop:galcoh:unramified-computed}(iii), applied to $V^D$, with the nonvanishing of $L_p(A_f,1)^{-1}$ supplied by \Cref{ch:lfunctions}; and the last sentence is
\Cref{cor:galcoh:ur-vanishing} applied to $T^D$, since $H^1_f(\Q_p,T^D)$ is defined in Lemma lem:dfg21 as the propagated condition $\iota^{-1}H^1_f(\Q_p,V^D)$.
\emph{Commentary.} The book proves slightly more than the paper states: $H^1_f(\Q_p,T^D)\simeq H^0(\Q_p,W^D)$, which is how \Cref{ch:blochkato} computes its length
(Lemma lem:fpr). The determinant in the revised paper is written with $\Frob_p^{-1}$, the geometric Frobenius, as \Cref{warn:galcoh:frobenius} requires (the
original version wrote $\Frob_p$, which gives $1-p^2$ instead of $1-p^{-2}$ at a multiplicative prime); the vanishing, and the finiteness of $H^1_f(\Q_p,T^D)$,
hold in either convention.
\end{paperlink}

\begin{example}[unramified cohomology in the paper's situations]\label{ex:galcoh:ur-examples}
Let $p\ne\ell$.
\begin{enumerate}[label=(\alph*)]
\item $T=\Z_\ell(1)$: $\Hur(\Q_p,\Z_\ell(1))=\Z_\ell(1)/(p-1)\Z_\ell(1)$, cyclic of order $\ell^{v_\ell(p-1)}$; under \Cref{thm:galcoh:kummer-ladic} it is the image of
$\Z_p^\times$ (\Cref{prop:galcoh:unit-classes}).
\item $T=\Z_\ell$ with trivial action: every continuous homomorphism $\psi\colon G_{\Q_p}\to\Z_\ell$ is unramified. Indeed $\psi|_I$ factors through the tame character
(\Cref{ch:local}), and $\psi(\Frob\,t\,\Frob^{-1})=\psi(t^p)$ gives $\psi(t)=p\psi(t)$, so $\psi(t)=0$. Hence $H^1(\Q_p,\Z_\ell)=\Hur(\Q_p,\Z_\ell)=\Z_\ell$, the dual
$\Hur(\Q_p,\Q_\ell/\Z_\ell(1))=\mu_{\ell^\infty}/\mu_{\ell^\infty}^{p-1}=0$ (a divisible group), and \Cref{prop:galcoh:propagation}(ii), (iii) are consistent.
\item The Tate curve: if $E$ has multiplicative reduction at $p$, $T=T_\ell E$ has a basis $e_1,e_2$ with $I$ acting through $e_2\mapsto e_2+n'e_1$, $n'\ne0$, and
$\Frob e_1=\varepsilon pe_1$, $\varepsilon=\pm1$ (\Cref{ch:galrep}). So $V^I=\Q_\ell e_1$ with $\Frob$-eigenvalue $\varepsilon p\ne1$; $\Hur(\Q_p,V_\ell E)=0$ and
$H^1(\Q_p,T_\ell E)_{\tors}\simeq H^0(\Q_p,E[\ell^\infty])=E(\Q_p)[\ell^\infty]$.
\item $V=\ad V_\ell E$ at a multiplicative prime: $V^I=\Q_\ell e$, $e=e_1\otimes e_2^\vee$, with $\Frob e=pe$ (\Cref{ch:cohomology}, the model computation). For
$V^D=V(1)$ the eigenvalue of $\Frob$ on $(V^D)^I$ is $p^2$ and that of $\Phi$ is $p^{-2}$. So $\det(1-\Phi\mid(V^D)^I)=1-p^{-2}$, the paper's $L_p(A_f,1)^{-1}$ at a
multiplicative prime, while $\det(1-\Frob\mid(V^D)^I)=1-p^2$: different numbers, with the same $\ell$-adic valuation $v_\ell(p^2-1)$ and the same nonvanishing.
By \Cref{prop:galcoh:unramified-computed}(iv), $\#\Hur(\Q_p,T^D)=\ell^{v_\ell(p^2-1)}$.
\item (An illustration that depends on \Cref{ch:lfunctions}.) Types $\mathrm{III}$, $\mathrm{III}^*$ at $q\ge5$: by Lemma lem:euler (proved in \Cref{ch:lfunctions}) $(\ad V)^{I_q}$ is a line on which $\Frob$ acts by
$\chi_{-4}(q)=\pm1$. Then $(V^D)^{I_q}$ is a line with $\Frob$-eigenvalue $q\chi_{-4}(q)$, and $\#\Hur(\Q_q,T^D)=\ell^{v_\ell(1-q\chi_{-4}(q))}=\ell^{v_\ell(q-\chi_{-4}(q))}$, which is
$\ell^{e_\ell(E,q)}$ in the notation of Definition def:terms.
\item Good reduction at $p$: $V_\ell E$ is unramified and the eigenvalues of $\Frob$ have absolute value $\sqrt p$ (Hasse), so $1$ is not one of them and
$\Hur(\Q_p,V_\ell E)=0$.
\end{enumerate}
\end{example}
\section{The Kummer class of the Tate parameter}\label{sec:galcoh:tate}

This section proves the last sentences of the paper's proof of Lemma lem:tamfactor: the splitting of $H^1(\Q_q,\Z_\ell(1))$ given by the uniformiser $q$, the
identification of the unramified component of the Kummer class of the Tate parameter with the class of the reduction of its unit part, and the exact
indeterminacy of the corner entry $c'$ of $\rho(\Frob)$. The field-theoretic facts (Hensel's lemma, the tame character $t_\ell$, the decomposition
$\kappa_x=v_K(x)\kappa_\varpi+\kappa_u$, the value of $\kappa_u$ at Frobenius, the Frobenius lift fixing roots of the uniformiser) are proved in \Cref{ch:local}, and
the group cohomology of $\ad T$ is \Cref{ch:cohomology} (the model computation). What is added here is the cohomological reading and the identification of $c'$; the assembly of Lemma lem:tamfactor is \Cref{ch:blochkato}.

Throughout, $q\ne\ell$ are primes, $G=G_{\Q_q}$, $I=I_q$, $v=v_\ell(q-1)$, and $t_\ell\colon I\to\Z_\ell(1)$ is the tame character of \Cref{ch:local}: for any compatible
system of $\ell$-power roots of $q$, $t_\ell(\sigma)=\kappa_q(\sigma)$ for $\sigma\in I$, independently of the system.

\begin{proposition}[unramified Kummer classes are unit classes]\label{prop:galcoh:unit-classes}
Let $K$ be a finite extension of $\Q_p$, $p\ne\ell$, with residue field $k_K$, ring of integers $\mathcal O_K$ and uniformiser $\varpi$.
\begin{enumerate}[label=(\roman*)]
\item For $n$ prime to $p$, the Kummer isomorphism $K^\times/(K^\times)^n\simeq H^1(K,\mu_n)$ maps the image of $\mathcal O_K^\times$ onto $\Hur(K,\mu_n)$. In the limit,
$\mathcal O_K^\times\otimes\Z_\ell\subseteq\widehat{K^\times}$ corresponds to $\Hur(K,\Z_\ell(1))$.
\item The valuation sequence of \Cref{prop:galcoh:kummer-local}(i) corresponds to the sequence of \Cref{prop:galcoh:ur-sequence} for $M=\Z_\ell(1)$:
\[
\begin{array}{ccccccccc}
0&\to&\mathcal O_K^\times\otimes\Z_\ell&\to&\widehat{K^\times}&\xrightarrow{\ v_K\ }&\Z_\ell&\to&0\\
&&\downarrow&&\downarrow&&\downarrow&&\\
0&\to&\Hur(K,\Z_\ell(1))&\to&H^1(K,\Z_\ell(1))&\xrightarrow{\ \Res\ }&\Hom(I_K,\Z_\ell(1))^{G/I}&\to&0
\end{array}
\]
with vertical isomorphisms the Kummer map and $a\mapsto a\,t_\ell$; in particular $\Res\kappa_x=v_K(x)\,t_\ell$ and $\Hom_{\mathrm{cont}}(I_K,\Z_\ell(1))^{G/I}=\Z_\ell\,t_\ell$.
\end{enumerate}
\end{proposition}

\begin{proof}
(i) A unit $u$ has an $n$-th root in $K^{\ur}$: $X^n-\bar u$ is separable over $k_K$, so it has a simple root in a finite extension of $k_K$, which lifts by Hensel's lemma
to a root in the corresponding finite unramified extension of $K$ (\Cref{ch:local}). So $I_K$ fixes an $n$-th root of $u$, and the Kummer cocycle of $u$ built from it vanishes on $I_K$:
$\delta(u)\in\Hur(K,\mu_n)$. Conversely let $x=\varpi^au$ with $\Res_{I_K}\delta(x)=0$. By \Cref{prop:galcoh:kummer-class}(i) for $L=K^{\ur}$ (whose absolute Galois group is
$I_K$), $x$ is an $n$-th power in $K^{\ur}$; as $u$ is one, so is $\varpi^a$, and since $v_K$ extends to $K^{\ur}$ with the same value group, $n\mid a$. So $x\in\mathcal O_K^\times(K^\times)^n$.
The limit statement follows by \Cref{thm:galcoh:kummer-ladic}. (ii) The formula $\Res\kappa_x=v_K(x)t_\ell$ is \Cref{ch:local} (the theorem on the Kummer cocycle on
inertia), so the right square commutes. Every continuous homomorphism $I_K\to\Z_\ell(1)$ factors through $t_\ell$, which is surjective, so
$\Hom_{\mathrm{cont}}(I_K,\Z_\ell(1))=\Z_\ell t_\ell$; and $t_\ell$ is invariant under the conjugation action, since $t_\ell(\sigma^{-1}i\sigma)=\chi_\ell(\sigma)^{-1}t_\ell(i)$
(the tame relation, \Cref{ch:local}) gives $(\sigma\cdot t_\ell)(i)=\chi_\ell(\sigma)t_\ell(\sigma^{-1}i\sigma)=t_\ell(i)$. The left square commutes by (i), and the outer vertical
maps are isomorphisms, hence so is the middle one (as we knew from \Cref{thm:galcoh:kummer-ladic}).
\end{proof}

\begin{proposition}[the value of an unramified Kummer class]\label{prop:galcoh:unramified-value}
Let $K=\Q_q$. Then $\Z_\ell(1)^I=\Z_\ell(1)$, and $c\mapsto c(\Frob)$ gives
\[
H^1(G/I,\Z_\ell(1))\;\simeq\;\Z_\ell(1)/(q-1)\Z_\ell(1)\;=\;\mu_{\ell^v}(\F_q)\qquad(\text{cyclic of order }\ell^v).
\]
Let $u\in\Z_q^\times$ and let $\kappa_u$ be a Kummer cocycle built from a compatible system of $\ell$-power roots of $u$ in $\Q_q^{\ur}$. Then $\kappa_u$ is inflated from $G/I$
as a cocycle (not only as a class), so $\kappa_u(\Frob)$ depends only on the image of $\Frob$ in $G/I$, and its class in $\Z_\ell(1)/(q-1)$ is $\bar u^{(q-1)/\ell^v}$, that is, the
class of the reduction $\bar u$ in $\F_q^\times\otimes\Z_\ell=\F_q^\times/(\F_q^\times)^{\ell^v}$. Hence the composite
\[
\Z_q^\times\otimes\Z_\ell\xrightarrow{\ \sim\ }\Hur(\Q_q,\Z_\ell(1))=H^1(G/I,\Z_\ell(1))\xrightarrow{\ \sim\ }\Z_\ell(1)/(q-1)
\]
is the isomorphism ``reduce modulo $q$ and take the class in $\F_q^\times\otimes\Z_\ell$''.
\end{proposition}

\begin{proof}
$I$ acts trivially on $\Z_\ell(1)$ because $\Q_q^{\ur}$ contains all $\ell$-power roots of unity, and $\Frob$ acts by $q$; the isomorphism is the procyclic computation of
\Cref{ch:cohomology} (its example with $\Frob$ acting by an integer $q$ prime to $\ell$). For $\sigma\in I$ and $g\in G$, $\kappa_u(g\sigma)=\kappa_u(g)+\chi_\ell(g)\kappa_u(\sigma)=\kappa_u(g)$,
because the chosen roots of $u$ lie in $\Q_q^{\ur}$ and are fixed by $\sigma$; so $\kappa_u$ factors through $G/I$. The class of $\kappa_u(\Frob)$ is computed in
\Cref{ch:local} (the theorem on the Kummer cocycle on inertia and on Frobenius, part (iii)): it is the Teichm\"uller lift of $\bar u^{(q-1)/\ell^v}$. Finally
$\Z_q^\times\otimes\Z_\ell=\F_q^\times\otimes\Z_\ell$ because $1+q\Z_q$ is pro-$q$, and $y\mapsto y^{(q-1)/\ell^v}$ maps $\F_q^\times$ onto $\mu_{\ell^v}(\F_q)$ with kernel
$(\F_q^\times)^{\ell^v}$.
\end{proof}

After a choice of a generator of $\Z_\ell(1)$, $\Z_\ell(1)/(q-1)\simeq\Z/\ell^v$, and the valuation of a class is read as in \Cref{ex:galcoh:kummer-finite-field}: $v_\ell(c)=\min(v_\ell(\tilde c),v)$,
so $v_\ell(0)=v$; it does not depend on the generator.

\begin{proposition}[the canonical splitting attached to the uniformiser $q$]\label{prop:galcoh:canonical-splitting}
For $K=\Q_q$ the prime $q$ is a canonical uniformiser and $\Q_q^\times=q^{\Z}\times\Z_q^\times$ canonically. Hence
\[
H^1(\Q_q,\Z_\ell(1))=\Z_\ell\,[\kappa_q]\;\oplus\;\Hur(\Q_q,\Z_\ell(1))
\]
canonically, and there is a canonical retraction
\[
\operatorname{pr}_{\ur}\colon H^1(\Q_q,\Z_\ell(1))\to\Hur(\Q_q,\Z_\ell(1))=H^1(G/I,\Z_\ell(1)^I)
\]
with
$\operatorname{pr}_{\ur}[\kappa_x]=[\kappa_{xq^{-v_q(x)}}]$: it sends the class of $x$ to the class of its unit part.
\end{proposition}

\begin{proof}
Tensor the decomposition $\Q_q^\times=q^{\Z}\times\Z_q^\times$ with $\Z_\ell$ after $\ell$-adic completion: $\widehat{\Q_q^\times}=\Z_\ell\,\hat q\oplus(\Z_q^\times\otimes\Z_\ell)$. Apply the
Kummer isomorphism (\Cref{thm:galcoh:kummer-ladic}) and \Cref{prop:galcoh:unit-classes}(i).
\end{proof}

\begin{remark}
Two cautions. First, the retraction is not functorial: inflation goes from $H^1(G/I,-)$ to $H^1(G,-)$, not the other way, and over a general finite extension $K$ of
$\Q_q$ the analogous splitting depends on the choice of uniformiser (replacing $\varpi$ by $\varpi w$, $w$ a unit, changes $\operatorname{pr}_{\ur}[\kappa_\varpi]$ from $0$ to
$-[\kappa_w]$, since $\varpi=\varpi'w^{-1}$; compare \Cref{ex:galcoh:ex-uniformiser}). Over $\Q_q$ the uniformiser $q$ is canonical, and the paper's phrase ``the image of $\kappa$ in $H^1(G/I,\Z_\ell(1)^I)$'' has the precise meaning
$\operatorname{pr}_{\ur}[\kappa]$. Second, the value $\kappa(\Frob)$ of a cocycle at a Frobenius lift is a different object from the class $\operatorname{pr}_{\ur}[\kappa]$;
\Cref{thm:galcoh:tate-kummer-class}(iii), (iv) compare them.
\end{remark}

\begin{theorem}[the corner of $\rho(\Frob)$ at a multiplicative prime]\label{thm:galcoh:tate-kummer-class}
Let $E/\Q$ have multiplicative reduction at $q$, let $\ell\ne q$ be odd, $q_E\in\Q_q^\times$ the Tate parameter, $n=v_q(q_E)=n_q$, $u=q_Eq^{-n}\in\Z_q^\times$ its unit part,
and $\varepsilon\colon G\to\{\pm1\}$ the unramified character with $\varepsilon(\Frob)=a_q(E)$. Fix compatible systems of $\ell$-power roots of $q$ and of $u$ (the latter in
$\Q_q^{\ur}$), give $q_E$ the product system, and let $\kappa=\kappa_{q_E}$, $\kappa_q$, $\kappa_u$ be the Kummer cocycles. Fix a generator $\zeta$ of $\Z_\ell(1)$ and write
$\Z_\ell(1)=\Z_\ell$. By \Cref{ch:reduction,ch:galrep}, $T=T_\ell E$ has a basis $e_1,e_2$ in which
\[
\rho(g)=\begin{pmatrix}\chi_\ell(g)\varepsilon(g)&\varepsilon(g)\kappa(g)\\0&\varepsilon(g)\end{pmatrix}\qquad(g\in G),
\]
the corner being $\varepsilon\kappa$. Then:
\begin{enumerate}[label=(\roman*)]
\item $\kappa=n\kappa_q+\kappa_u$ as cocycles; $[\kappa]$ is the Kummer class of $q_E$, and $\operatorname{pr}_{\ur}[\kappa]=[\kappa_u]$ is the class of $\bar u$ in
$\F_q^\times\otimes\Z_\ell=\Z_\ell(1)/(q-1)$.
\item $\kappa|_I=n\,t_\ell$. For $t\in I$ with $t_\ell(t)=\zeta$, $te_1=e_1$ and $te_2=e_2+ne_1$: the paper's $n'$ may be taken equal to $n$, and in any case $v_\ell(n')=v_\ell(n)$.
\item Let $\Frob_0$ be a Frobenius lift fixing the chosen $\ell$-power roots of $q$. Then $\kappa_q(\Frob_0)=0$ and
\[
\rho(\Frob_0)=\begin{pmatrix}\varepsilon q&c'\\0&\varepsilon\end{pmatrix},\qquad\varepsilon=a_q(E),\quad c'=\varepsilon\,\kappa_u(\Frob_0),
\]
and the class of $c'$ in $\Z_\ell(1)/(q-1)$ is $\varepsilon$ times the class of $\bar u$.
\item For an arbitrary lift $\Frob=\Frob_0\sigma$ ($\sigma\in I$) and an arbitrary cocycle in the class of $\kappa$ (another choice of roots of $q_E$), the corner $c'$ of $\rho(\Frob)$
changes by an element of the ideal $(n,q-1)\Z_\ell$: $\kappa(\Frob_0\sigma)=\kappa(\Frob_0)+q\,n\,t_\ell(\sigma)$, and a coboundary changes the value at $\Frob$ by $(q-1)z$,
$z\in\Z_\ell$. A change of $\zeta$ multiplies $c'$ by a unit.
\item Consequently, for every choice,
\[
\min\bigl(v_\ell(n),v_\ell(q-1),v_\ell(c')\bigr)=v_\ell\bigl((n,q-1,c')\Z_\ell\bigr)=\min\bigl(v_\ell(n),v_\ell(q-1),v_\ell(c)\bigr),
\]
where $c\in\Z/\ell^v$ is the class of $\bar u$ (for any generator) and $v_\ell(0)=v$.
\end{enumerate}
\end{theorem}

\begin{proof}
(i) The cocycle identity $\kappa=n\kappa_q+\kappa_u$ for the product system is \Cref{ch:local} (part (i) of its theorem on the Kummer cocycle). That $[\kappa]$ is the Kummer
class of $q_E$ is \Cref{thm:galcoh:kummer-ladic}(i). By \Cref{prop:galcoh:canonical-splitting}, $\operatorname{pr}_{\ur}[\kappa]=[\kappa_u]$, and by
\Cref{prop:galcoh:unramified-value} its value is the class of $\bar u$.
(ii) $\kappa_u$ vanishes on $I$ (its roots lie in $\Q_q^{\ur}$) and $\kappa_q|_I=t_\ell$, so $\kappa|_I=nt_\ell$; this is independent of the roots, since coboundaries
$(\chi_\ell-1)z$ vanish on $I$. As $\varepsilon|_I=1$ and $\chi_\ell|_I=1$, $\rho(t)=\sm1{n}01$ in the basis $e_1,e_2$.
(iii) $\Frob_0$ fixes the chosen roots of $q$, so $\kappa_q(\Frob_0)=0$ and $\kappa(\Frob_0)=\kappa_u(\Frob_0)$; also $\chi_\ell(\Frob_0)=q$ and $\varepsilon(\Frob_0)=a_q(E)$.
The class of $\kappa_u(\Frob_0)$ modulo $(q-1)$ is the class of $\bar u$ by \Cref{prop:galcoh:unramified-value}. Such a lift exists by \Cref{ch:local} (its corollary on
the Frobenius lift fixing roots).
(iv) By the cocycle identity and (ii), $\kappa(\Frob_0\sigma)=\kappa(\Frob_0)+\chi_\ell(\Frob_0)\kappa(\sigma)=\kappa(\Frob_0)+q\,n\,t_\ell(\sigma)$, and $\varepsilon(\Frob_0\sigma)=\varepsilon(\Frob_0)$; so
$c'$ changes by $\varepsilon qnt_\ell(\sigma)\in n\Z_\ell$. Another system of roots of $q_E$ changes $\kappa$ by the coboundary $g\mapsto(\chi_\ell(g)-1)z$ (\Cref{ch:local}),
whose value at any Frobenius lift is $(q-1)z$; it corresponds to the change of basis $e_2\mapsto e_2+ze_1$. Replacing $\zeta$ by $\zeta^a$, $a\in\Z_\ell^\times$, multiplies all
coordinates of $\kappa$ by $a^{-1}$.
(v) For elements $x_1,\dots,x_r$ of $\Z_\ell$, the ideal they generate is $\ell^m\Z_\ell$ with $m=\min_iv_\ell(x_i)$ (with $v_\ell(0)=\infty$). By (iv) the ideal $(n,q-1,c')$
does not depend on any choice, and by (iii) it equals $(n,q-1,\tilde c)$ for any $\tilde c\in\Z_\ell$ lifting the class of $\bar u$ in $\Z_\ell/(q-1)=\Z/\ell^v$ (the sign
$\varepsilon$ is a unit). Finally $\min(v_\ell(n),v,v_\ell(\tilde c))=\min(v_\ell(n),v,\min(v_\ell(\tilde c),v))=\min(v_\ell(n),v,v_\ell(c))$.
\end{proof}

\begin{corollary}[the formula of Lemma lem:tamfactor]\label{cor:galcoh:tamfactor}
In the situation of \Cref{thm:galcoh:tate-kummer-class}, with $M=\ad T$, the model computation of \Cref{ch:cohomology} gives
\[
\len\Bigl(\bigl(H^1(I,\ad T)_{\tors}\bigr)^{G/I}\Bigr)=v_\ell(n)+\min\bigl(v_\ell(n),v_\ell(q-1),v_\ell(c)\bigr).
\]
\end{corollary}

\begin{proof}
\Cref{ch:cohomology} proves the length $v_\ell(n')+\min(v_\ell(n'),v_\ell(q-1),v_\ell(c'))$ for matrices of this shape with any $n'$ and $c'$; this chapter contributes only the
identification of $n'$ and $c'$: $v_\ell(n')=v_\ell(n)$ by \Cref{thm:galcoh:tate-kummer-class}(ii), and the minimum by part (v).
\end{proof}

\noindent Lemma lem:tamfactor itself (the Fitting ideal of this module is $\Tam_q(\ad T)$, the assembly of the proof, and the additive potentially multiplicative case) is
\Cref{ch:blochkato}; the corollary is the pointer from the two computations to that assembly.

\begin{paperlink}\label{pl:galcoh:tamfactor}
The end of the proof of Lemma lem:tamfactor in the revised paper reads: ``Finally, $q$ being a uniformiser of $\Q_q$, $H^1(\Q_q,\Z_\ell(1))$ splits as
$\Z_\ell\kappa_q\oplus H^1_{\mathrm{ur}}$, and $\kappa=n\kappa_q+\kappa_u$, where $\kappa_u$ is unramified, inflated from the class of $u$ in
$H^1(G/I,\Z_\ell(1))=\Z_\ell(1)/(q-1)\simeq\F_q^\times\otimes\Z_\ell$. The value $c'=\varepsilon\kappa(\Frob)$ is a cocycle value: it is $\varepsilon$ times the class of $u$ for the
Frobenius lift fixing a compatible system of $\ell$-power roots of $q$, and in general it is well defined only modulo $\ell^{\min(v_\ell(n),v_\ell(q-1))}$ (changing the lift
changes it by a multiple of $n'$, changing the cocycle by a multiple of $q-1$).'' The splitting is \Cref{prop:galcoh:canonical-splitting}; the decomposition and the
inflation are \Cref{thm:galcoh:tate-kummer-class}(i) and \Cref{prop:galcoh:unramified-value}; the value at the special lift is part (iii), and the indeterminacy, the
ideal $(n,q-1)\Z_\ell=\ell^{\min(v_\ell(n),v_\ell(q-1))}\Z_\ell$, is part (iv). ``The class of $u$'' means the class of its reduction $\bar u$, as the statement of the lemma
says. The original version of the paper wrote $\kappa$ instead of $\varepsilon\kappa$ in the corner and ``well defined modulo $q-1$''; the revised wording is the one proved
here. The check ``the same computation for $M=T$ gives $(\Z/\ell^{v_\ell(n)})\bar e_1$ with $\Frob$ acting by $\varepsilon$, that is, the $\ell$-part of the Tamagawa number
$c_q(E)$'' is \Cref{prop:galcoh:ec-local-image}(iv) below, in its cohomological form.
\end{paperlink}

\begin{example}[Tate parameters and the extra term at $\ell=3$]\label{ex:galcoh:tate-kummer-examples}
The table lists optimal curves with a multiplicative prime $q\equiv1\pmod3$, $3\mid n$, and $3\nmid N$; $c$ is the class of $\bar u$ in $\F_q^\times/(\F_q^\times)^{3^{v_3(q-1)}}$,
``extra'' is $\min(v_3(n),v_3(q-1),v_3(c))$, and ``bound'' is $\sum_qt_3(E,q)+\text{extra}$, the right side of (eq:identity) without the Selmer term (these curves have no additive
prime, so no term $e_3$). The last row is at $\ell=5$. The Tate parameters were computed with PARI as \texttt{ellinit(E,O(q\^{}30)).tate[3]}, the Kodaira symbols with
\texttt{elllocalred}, the degrees with \texttt{ellmoddegree} and checked against Cremona's \texttt{degphi} file, and the mod $\ell$ images read from Cremona's
\texttt{galrep} file.
\begin{center}
\small\begin{tabular}{lccccccc}
curve, $q$&$a_q$&$n$&$\bar u$&$\bar u$ a cube?&extra&bound&$v_3(\degphi)$\\\hline
\curve{77c1}, $7$&$-1$&$3$&$5$&no&$0$&$1$&$1$ ($\degphi=6$)\\
\curve{791a1}, $7$&$-1$&$3$&$6$&yes&$1$&$2$&$2$ ($\degphi=144$)\\
\curve{1577a1}, $19$&$+1$&$3$&$8$&yes, not a ninth power&$1$&$2$&$2$ ($\degphi=720$)\\
\curve{11a1}, $11$ ($\ell=5$)&$+1$&$5$&$10$&a fifth power&$1$&(Eisenstein)&$v_5(1)=0$
\end{tabular}
\end{center}
The cubes in $\F_7^\times$ are $\{1,6\}$; the cubes in $\F_{19}^\times$ are $\{1,7,8,11,12,18\}$ and the ninth powers $\{1,18\}$; the fifth powers in $\F_{11}^\times$ are $\{1,10\}$.
The other bad primes are $11$ for \curve{77c1}, $113$ for \curve{791a1} and $83$ for \curve{1577a1}, all of type $\mathrm I_2$, with $t_3=0$. In the first three rows the
mod $3$ image is all of $\GL_2(\F_3)$ (the \texttt{galrep} file lists no exceptional image at $3$), so the hypothesis of Theorem thm:main holds, and the bound equals
$v_3(\degphi)$: the extra term of Remark rem:ko(1) is present exactly when $\bar u$ is a cube, and the Selmer term is $0$. For \curve{11a1} at $\ell=5$ the local term
$t_5+\text{extra}=2$ exceeds $v_5(\degphi)=0$: here $E[5]$ is reducible ($E(\Q)_{\tors}=\Z/5$, image \texttt{5Cs.1.1}), and the example shows why the theorem needs $\bar\rho$
irreducible.
\begin{lstlisting}[language=GP]
E = ellinit([1,0,1,-31,117]);           \\ 791a1
qE = ellinit(E, O(7^30)).tate[3];
n = valuation(qE, 7); ub = lift(Mod(truncate(qE/7^n), 7));
[n, ub, ispower(Mod(ub,7), 3), ellap(E,7), ellmoddegree(E)]
\\ [3, 6, 1, -1, 144]
\end{lstlisting}
\end{example}

\begin{remark}[Lemma lem:tamagawa-exp in cohomological form]\label{rem:galcoh:tamagawa-exp-cohomological}
Let $E$ be multiplicative at $q\ne\ell$ and $t\ge1$. By \Cref{thm:galcoh:tate-kummer-class}(ii) the action of $I$ on $T/\ell^tT$ is $\sm1{nt_\ell}01$ modulo $\ell^t$, for split and
nonsplit reduction alike ($\varepsilon|_I=1$). So $\rho\bmod\ell^t$ is unramified at $q$ if and only if $nt_\ell\equiv0\pmod{\ell^t}$ on $I$, if and only if $\ell^t\mid n$, $t_\ell$ being
surjective. Equivalently, by \Cref{prop:galcoh:kummer-class}(i) over $\Q_q^{\ur}\supset\mu_{\ell^t}$: $q_E\in\bigl((\Q_q^{\ur})^\times\bigr)^{\ell^t}$ if and only if $\ell^t\mid n$. The lemma
itself is proved in \Cref{ch:galrep}; this remark only gives the cohomological reading of the paper's phrase ``By Kummer theory over $\Q_q^{\mathrm{ur}}$''.
\end{remark}
\section{Global Galois cohomology: \texorpdfstring{$G_\Q$}{G\_Q}, \texorpdfstring{$G_{\Q,S}$}{G\_Q,S} and the real place}\label{sec:galcoh:global}

The Selmer groups of Theorem thm:main are subgroups of the cohomology of $G_\Q$ cut out by local conditions. This section introduces the group $G_{\Q,S}$ through
which the cohomology of a module unramified outside $S$ is computed, proves that classes unramified outside $S$ are exactly the classes of $G_{\Q,S}$, states the
finiteness theorem behind every finiteness statement of the paper, and disposes of the real place, which for $\ell$ odd contributes nothing (the paper's
``At $\infty$, $H^1(\R,T^D)$ is killed by $2$, so $H^1_f(\R,T^D)=0$ for $\ell$ odd'').

\begin{definition}[places, $\Q_S$, $G_{\Q,S}$]\label{def:galcoh:gqs}
$p$ runs over the primes and $v$ over the places of $\Q$, including $\infty$, with $\Q_\infty=\R$. For each $v$ fix an embedding $\overline\Q\to\overline\Q_v$; it identifies $G_{\Q_v}$
with a decomposition group in $G_\Q$ (\Cref{ch:ant,ch:local}), and gives the \emph{localisation} $\Res_v\colon H^n(\Q,M)\to H^n(\Q_v,M)$. A different embedding replaces
$G_{\Q_v}$ by a conjugate, and the conjugation action of $G_\Q$ on its own cohomology is trivial (\Cref{ch:cohomology}), so $\Res_v$ does not depend on the
embedding. $S$ is a finite set of places containing $\infty$; $\Q_S\subseteq\overline\Q$ is the maximal extension of $\Q$ unramified at every prime $p\notin S$; and
$G_{\Q,S}=\Gal(\Q_S/\Q)=G_\Q/H_S$, where $H_S$ is the closed normal subgroup generated by the inertia groups $I_{\mathfrak p}$ of all primes $\mathfrak p$ of $\overline\Q$ above
primes $p\notin S$. A $G_\Q$-module $M$ (finite, of type (C) or (V), or discrete $\ell$-primary such as $W$) is \emph{unramified outside $S$} if $H_S$ acts trivially; it is then a
$G_{\Q,S}$-module. In the paper $\Sigma$ is a finite set of primes (there, the primes dividing $N$) with $\ell\notin\Sigma$.
\end{definition}

\noindent Examples: $\mu_n$ is unramified outside $\{\infty\}\cup\{p\mid n\}$; $E[n]$, $T_\ell E$ and $\ad T_\ell E$ are unramified outside $\{\infty,\ell\}\cup\{p\mid N\}$, $N$ the
conductor of $E$, by the criterion of N\'eron, Ogg and Shafarevich (\Cref{ch:galrep}).

\begin{proposition}[classes unramified outside $S$]\label{prop:galcoh:unramified-almost-everywhere}
Let $M$ be unramified outside $S$.
\begin{enumerate}[label=(\roman*)]
\item Inflation identifies $H^1(G_{\Q,S},M)$ with $\{c\in H^1(\Q,M):\Res_pc\in\Hur(\Q_p,M)\text{ for all }p\notin S\}$.
\item If $M$ is finite, every $c\in H^1(\Q,M)$ is unramified at all but finitely many primes.
\end{enumerate}
\end{proposition}

\begin{proof}
(i) Inflation-restriction for $H_S\trianglelefteq G_\Q$ (\Cref{ch:cohomology}), with $M^{H_S}=M$, gives the exact sequence
$0\to H^1(G_{\Q,S},M)\to H^1(\Q,M)\to H^1(H_S,M)=\Hom_{\mathrm{cont}}(H_S,M)$. So $c$ is inflated from $G_{\Q,S}$ if and only if the homomorphism $c|_{H_S}$ is zero. The set of
inertia groups $I_{\mathfrak p}$, $\mathfrak p$ above $p\notin S$, is stable under conjugation ($\sigma I_{\mathfrak p}\sigma^{-1}=I_{\sigma\mathfrak p}$), so the closed subgroup it generates is
normal and equals $H_S$; and the kernel of a continuous homomorphism is a closed subgroup. Hence $c|_{H_S}=0$ if and only if $c|_{I_{\mathfrak p}}=0$ for every such $\mathfrak p$.
For $i\in I_{\mathfrak p}$, which acts trivially on $M$, and $\sigma\in G_\Q$, the cocycle identity gives $c(\sigma i\sigma^{-1})=c(\sigma)+\sigma c(i)+\sigma i\,c(\sigma^{-1})=\sigma c(i)$,
because $\sigma i\,c(\sigma^{-1})=\sigma c(\sigma^{-1})=-c(\sigma)$. So $c$ vanishes on $I_{\sigma\mathfrak p}$ if and only if it vanishes on $I_{\mathfrak p}$, and it suffices to take the
prime $\mathfrak p$ of the fixed embedding, where the condition is $\Res_pc\in\Hur(\Q_p,M)$.
(ii) Let $L$ be the finite Galois extension cut out by the action on $M$. The restriction of $c$ to $G_L$ is a continuous homomorphism to the finite group $M$; its kernel is
an open subgroup of $G_L$, with fixed field $L'$ finite over $\Q$; let $L''$ be the Galois closure of $L'$. If $p$ is unramified in $L''$, every $I_{\mathfrak p}$ above $p$
lies in $G_{L''}\subseteq G_{L'}$, on which $c$ vanishes. Only finitely many primes ramify in $L''$.
\end{proof}

\begin{theorem}[finiteness for $G_{\Q,S}$]\label{thm:galcoh:gqs-finite}
Let $M$ be finite and unramified outside the finite set $S$.
\begin{enumerate}[label=(\roman*)]
\item $H^1(G_{\Q,S},M)$ is finite.
\item (Stated.) $H^2(G_{\Q,S},M)$ is finite.
\item (Stated.) If $\ell$ is odd, $M$ has $\ell$-power order and $\ell\in S$, then $H^n(G_{\Q,S},M)=0$ for $n\ge3$.
\end{enumerate}
Consequently, for a lattice $T$ unramified outside $S$, \Cref{prop:galcoh:tvw} applies with $G=G_{\Q,S}$: $H^1(G_{\Q,S},T)$ is finitely generated, $H^1(G_{\Q,S},W)$ is
cofinitely generated, and $H^1(G_{\Q,S},T)_{\tors}\simeq H^0(\Q,W)/H^0(\Q,W)_{\mathrm{div}}$.
\end{theorem}

\begin{proof}[Sketch of proof of (i)]
Let $L$ be the field cut out by $M$; it lies in $\Q_S$, and $H=\Gal(\Q_S/L)$ is open and normal in $G_{\Q,S}$ with trivial action on $M$. Inflation-restriction gives
$0\to H^1(\Gal(L/\Q),M)\to H^1(G_{\Q,S},M)\to\Hom_{\mathrm{cont}}(H,M)$. The first group is finite. A continuous homomorphism $H\to M$ has kernel $\Gal(\Q_S/L')$ with $L'/L$
Galois of degree dividing $\#M$ and unramified outside $S$, and by the theorem of Hermite and Minkowski (there are finitely many number fields of bounded degree
unramified outside a given finite set of primes; stated, \cite[Chapter VIII, Section 1]{SilvermanAEC}) there are finitely many such $L'$, each with finitely many
homomorphisms. For (ii) and (iii) see \cite[Chapter VIII, Section 3]{NSW} and \cite[Chapter I, Section 4]{MilneADT}. The consequence is \Cref{prop:galcoh:tvw}, whose
hypothesis is (i) and (ii) for $M=T/\ell T$; note $H^0(G_{\Q,S},W)=H^0(\Q,W)$ since $H_S$ acts trivially.
\end{proof}

\begin{proposition}[the real place]\label{prop:galcoh:real-place}
Let $G_\R=\{1,c\}$ and $M$ a $G_\R$-module. Then $2H^n(\R,M)=0$ for $n\ge1$. If $\ell$ is odd and $M$ is discrete $\ell$-primary, of type (C) or of type (V), then
$H^n(\R,M)=0$ for $n\ge1$, the Tate group $\hat H^0(\R,M)=M^{c=1}/(1+c)M$ is zero, and $M=M^+\oplus M^-$ with $M^\pm=\{m:cm=\pm m\}$. In particular, for $\ell$ odd every
local condition at $\infty$ is $0=H^1(\R,M)$, and $H^1_f(\R,T^D)\subseteq H^1(\R,T^D)=0$.
\end{proposition}

\begin{proof}
A finite group of order $2$ kills its cohomology in positive degrees (\Cref{ch:cohomology}; for continuous cochains with coefficients of type (C) as well). If $2$ is
invertible on $M$, multiplication by $2$ is an isomorphism of $H^n(\R,M)$ that is also zero, so $H^n(\R,M)=0$. If $cm=m$ then $m=(1+c)(m/2)$, so $\hat H^0=0$; and
$m=\frac{m+cm}2+\frac{m-cm}2$ gives the decomposition.
\end{proof}

\noindent For $\ell=2$ this fails: by \Cref{thm:galcoh:kummer}(i), $H^1(\R,\mu_2)=\R^\times/\R^{\times2}=\{\pm1\}$.

\begin{paperlink}\label{pl:galcoh:real}
The last sentence of the paragraph after Lemma lem:dfg21, ``At $\infty$, $H^1(\R,T^D)$ is killed by $2$, so $H^1_f(\R,T^D)=0$ for $\ell$ odd'', is
\Cref{prop:galcoh:real-place}: the first half holds for every module, the second needs $\ell$ odd, and \Cref{ex:galcoh:real-two} shows that it fails for $\ell=2$. The
summand $p=\infty$ of the sequence of Lemma lem:dfg21 is therefore $0$.
\end{paperlink}

\begin{example}[global Kummer theory]\label{ex:galcoh:global-kummer}
By \Cref{thm:galcoh:kummer}, $H^1(\Q,\mu_n)=\Q^\times/(\Q^\times)^n$. For $p\nmid n$ the localisation of $\delta(x)$ is $\delta_p(x)$ (\Cref{thm:galcoh:kummer}(iv)), which is
unramified if and only if $x\in\Z_p^\times(\Q_p^\times)^n$, that is, $n\mid v_p(x)$ (\Cref{prop:galcoh:unit-classes}(i)). So for $S\supseteq\{\infty\}\cup\{p\mid n\}$, by
\Cref{prop:galcoh:unramified-almost-everywhere}(i),
\[
H^1(G_{\Q,S},\mu_n)=\{x\in\Q^\times:n\mid v_p(x)\ \forall p\notin S\}/(\Q^\times)^n\;\simeq\;\Z_S^\times/(\Z_S^\times)^n,\qquad\Z_S^\times=\pm\prod_{p\in S\setminus\{\infty\}}p^{\Z}.
\]
(Unique factorisation in $\Z$: an $x$ with $n\mid v_p(x)$ for $p\notin S$ is a unit of $\Z_S=\Z[1/p:p\in S]$ times an $n$-th power, and a unit of $\Z_S$ that is an $n$-th power in
$\Q$ is an $n$-th power of a unit.) The group is finite, of order $n^{\#S-1}$ times $2$ or $1$ according as $n$ is even or odd: the model case of
\Cref{thm:galcoh:gqs-finite}.
\end{example}

\begin{example}\label{ex:galcoh:real-two}
$H^1(\R,\Z_2(1))=\Z/2$: complex conjugation acts on $\Z_2(1)$ by $-1$, and for the group of order $2$, $H^1=\ker(1+c)/(c-1)M=M/2M$.
\end{example}

\section{Elliptic curves: the Kummer map, Selmer groups and Sha}\label{sec:galcoh:selmer-ec}

The classical Selmer group of an elliptic curve is the model of every Selmer group in the paper, and its finiteness is the weak Mordell-Weil theorem that
\Cref{ch:ellcurves} deferred to this chapter. The local computation \Cref{prop:galcoh:ec-local-image}(iv) is the classical shadow of the Tamagawa factors:
for $W=E[\ell^\infty]$ at a multiplicative prime $p\ne\ell$, the unramified cohomology has the order of the $\ell$-part of the Tamagawa number $c_p$, which is the check
the paper performs at the end of the proof of Lemma lem:tamfactor.

\begin{proposition}[the Kummer sequence of an elliptic curve]\label{prop:galcoh:ec-kummer}
Let $E/K$ be an elliptic curve and $n$ prime to $\operatorname{char}K$. The sequence $0\to E[n]\to E(\overline K)\xrightarrow{[n]}E(\overline K)\to0$ is an exact sequence of
discrete $G_K$-modules, and its long exact sequence gives the \emph{Kummer map}
\[
\delta\colon E(K)/nE(K)\hookrightarrow H^1(K,E[n]),\qquad\delta(P)=[\sigma\mapsto\sigma Q-Q]\ \text{ with }nQ=P,
\]
and the exact sequence
\[
0\to E(K)/nE(K)\to H^1(K,E[n])\to H^1(K,E)[n]\to0,\qquad H^1(K,E):=H^1(K,E(\overline K)).
\]
\end{proposition}

\begin{proof}
$[n]$ is a nonconstant isogeny, hence surjective on $\overline K$-points, with kernel $E[n]$ (\Cref{ch:ellcurves}); and every point of $E(\overline K)$ has an open
stabiliser. The long exact sequence $E(K)\xrightarrow{n}E(K)\to H^1(K,E[n])\to H^1(K,E)\xrightarrow{n}H^1(K,E)$ gives the rest, and the connecting map is that of
\Cref{conv:galcoh:signs}.
\end{proof}

\noindent Replacing $E$ by the multiplicative group gives Kummer theory, and Hilbert 90 says that the analogue $H^1(K,\overline K^\times)$ of $H^1(K,E)$ vanishes; for
elliptic curves $H^1(K,E)$ need not vanish, and this is where $\Sha$ comes from.

\begin{proposition}[the local Kummer image]\label{prop:galcoh:ec-local-image}
Let $E/\Q_p$ be an elliptic curve and $n$ prime to $p$.
\begin{enumerate}[label=(\roman*)]
\item $\#E(\Q_p)/nE(\Q_p)=\#E(\Q_p)[n]=\#H^0(\Q_p,E[n])=\#\Hur(\Q_p,E[n])$.
\item If $E$ has good reduction at $p$, then $\delta(E(\Q_p)/nE(\Q_p))=\Hur(\Q_p,E[n])$.
\item (Tate; stated.) For every reduction type, $\delta(E(\Q_p)/nE(\Q_p))$ is its own exact annihilator under the pairing of \Cref{thm:galcoh:tate-local-duality}(i) with
$E[n]^D\simeq E[n]$.
\item Let $\ell\ne p$ be odd and $W=E[\ell^\infty]$. The Kummer image $E(\Q_p)\otimes\Q_\ell/\Z_\ell$ in $H^1(\Q_p,W)$ is $0$, while $\Hur(\Q_p,W)=W^I/(\Frob-1)W^I$; if $E$ has
multiplicative reduction of type $\mathrm I_m$ at $p$, then $\#\Hur(\Q_p,W)=\ell^{v_\ell(m)}$ if the reduction is split and $1$ if it is nonsplit: the $\ell$-part of the
Tamagawa number $c_p$.
\end{enumerate}
\end{proposition}

\begin{proof}
(i) Let $E_1(\Q_p)$ be the kernel of reduction; it is the group of the formal group on $p\Z_p$, on which $[n]$ is an isomorphism because $n$ is a unit (\Cref{ch:reduction}),
and $F=E(\Q_p)/E_1(\Q_p)$ is finite. The snake lemma for $[n]$ on $0\to E_1(\Q_p)\to E(\Q_p)\to F\to0$ gives $E(\Q_p)[n]\simeq F[n]$ and $E(\Q_p)/n\simeq F/n$, and
$\#F[n]=\#F/nF$ for a finite group. The last equality is \Cref{prop:galcoh:unramified-computed}(ii).
(ii) Let $P\in E(\Q_p)$. The reduction $\tilde E(\overline\F_p)$ is $n$-divisible, so there is a point of $\tilde E(\overline\F_p)$ with $n$ times it equal to $\tilde P$; it lies in
$\tilde E(\F_{p^r})$ for some $r$ and lifts, by Hensel's lemma for the smooth model, to $Q_0\in E(K')$, $K'$ the unramified extension of degree $r$. Then $P-nQ_0$ lies in
$E_1(K')$, on which $[n]$ is bijective, so $P=nQ$ with $Q\in E(K')\subseteq E(\Q_p^{\ur})$. Inertia fixes $Q$, so $\delta(P)$ vanishes on $I_p$: $\delta(E(\Q_p)/n)\subseteq\Hur(\Q_p,E[n])$,
and the orders agree by (i). (iii) See \cite[Chapter I, Section 3]{MilneADT}.
(iv) $E(\Q_p)$ has the finite-index subgroup $E_1(\Q_p)$, a pro-$p$ group, which is uniquely $\ell$-divisible; so $E(\Q_p)\otimes\Q_\ell/\Z_\ell=0$. The formula for $\Hur$ is
\Cref{prop:galcoh:unramified-computed}(i). For type $\mathrm I_m$, $T_\ell E$ has a basis with $\rho(t)=\sm1{m'}01$, $v_\ell(m')=v_\ell(m)$, and $\rho(\Frob)=\sm{\varepsilon p}{*}0\varepsilon$
(\Cref{thm:galcoh:tate-kummer-class}, \Cref{ch:galrep}). In $W=(\Q_\ell/\Z_\ell)e_1\oplus(\Q_\ell/\Z_\ell)e_2$ an element $xe_1+ye_2$ is $I$-invariant iff $m'y=0$, so $W^I$ is an extension of
$\Z/\ell^{v_\ell(m)}$ (the $y$-coordinate) by $(\Q_\ell/\Z_\ell)e_1\simeq\mu_{\ell^\infty}$. On the subgroup $(\Q_\ell/\Z_\ell)e_1$, $\Frob-1$ is multiplication by $\varepsilon p-1\ne0$, which is
surjective on a divisible group; on the quotient it is multiplication by $\varepsilon-1$. So $W^I/(\Frob-1)W^I\simeq(\Z/\ell^{v_\ell(m)})/(\varepsilon-1)$, which is
$\Z/\ell^{v_\ell(m)}$ if $\varepsilon=1$ and $0$ if $\varepsilon=-1$, $2$ being a unit. The Tamagawa number is $c_p=m$ in the split case and $1$ or $2$ in the nonsplit case
(\Cref{ch:reduction}).
\end{proof}

\noindent Part (iv) agrees with the check at the end of the proof of Lemma lem:tamfactor, which computes the invariants of the torsion of $H^1(I,T)$ instead
(\Cref{ch:cohomology}, the model computation for $M=T$): for $M=T_\ell E$, the propagated condition $\iota^{-1}\Hur(\Q_p,V)$ is all of $H^1(\Q_p,T)_{\tors}$ (\Cref{ex:galcoh:ur-examples}(c)),
and the gap between it and $\Hur(\Q_p,T)$ is measured by the same Tamagawa number.

\begin{definition}[the $n$-Selmer group and Sha]\label{def:galcoh:selmer-sha}
For $E/\Q$ and $n\ge1$,
\begin{gather*}
\Sel_n(E/\Q)=\bigl\{c\in H^1(\Q,E[n]):\Res_vc\in\delta_v\bigl(E(\Q_v)/nE(\Q_v)\bigr)\text{ for all }v\bigr\},\\
\Sha(E/\Q)=\ker\Bigl(H^1(\Q,E)\to\prod_vH^1(\Q_v,E)\Bigr),
\end{gather*}
where $\delta_v$ is the Kummer map of $E/\Q_v$.
\end{definition}

\begin{proposition}\label{prop:galcoh:selmer-sha}
The sequence $0\to E(\Q)/nE(\Q)\xrightarrow{\ \delta\ }\Sel_n(E/\Q)\to\Sha(E/\Q)[n]\to0$ is exact.
\end{proposition}

\begin{proof}
Compare the sequence of \Cref{prop:galcoh:ec-kummer} for $\Q$ with those for the $\Q_v$, through the localisation maps. A class $c\in H^1(\Q,E[n])$ lies in $\Sel_n$
if and only if its image in each $H^1(\Q_v,E)[n]$ is zero, that is, if and only if its image $\bar c\in H^1(\Q,E)[n]$ lies in $\Sha$. So $\Sel_n$ is the preimage of
$\Sha[n]$ under the surjection $H^1(\Q,E[n])\to H^1(\Q,E)[n]$, whose kernel $\delta(E(\Q)/n)$ lies in $\Sel_n$ (the localisation of a global Kummer class is a local
Kummer class).
\end{proof}

\begin{theorem}[finiteness of the Selmer group; weak Mordell-Weil]\label{thm:galcoh:selmer-finite}
Let $E/\Q$ have conductor $N_E$, $n\ge2$, and $S=\{\infty\}\cup\{p\mid nN_E\}$. Then $\Sel_n(E/\Q)\subseteq H^1(G_{\Q,S},E[n])$, and $\Sel_n(E/\Q)$ is finite.
\end{theorem}

\begin{proof}
For $p\notin S$, $E$ has good reduction and $p\nmid n$, so the local condition at $p$ is $\Hur(\Q_p,E[n])$ (\Cref{prop:galcoh:ec-local-image}(ii)); $E[n]$ is unramified outside $S$,
and \Cref{prop:galcoh:unramified-almost-everywhere}(i) puts $\Sel_n$ inside $H^1(G_{\Q,S},E[n])$, which is finite by \Cref{thm:galcoh:gqs-finite}(i).
\end{proof}

\begin{corollary}[weak Mordell-Weil]\label{cor:galcoh:weak-mw}
$E(\Q)/nE(\Q)$ is finite for every $n\ge1$.
\end{corollary}

\begin{proof}
It injects into $\Sel_n(E/\Q)$ by \Cref{prop:galcoh:selmer-sha}.
\end{proof}

\noindent With the descent by the canonical height (\Cref{ch:ellcurves}; \cite[Chapter VIII]{SilvermanAEC}) this gives the theorem of Mordell and Weil that $E(\Q)$
is finitely generated.

\begin{example}[\curve{11a1} and $n=5$]\label{ex:galcoh:selmer-11a1}
For $E=\curve{11a1}$, Cremona's \texttt{allbsd} file gives rank $0$, $E(\Q)\simeq\Z/5$, Tamagawa product $5$ and analytic order of $\Sha$ equal to $1$. The same value
is printed in \cite[Table 4]{Cremona97} (row 11A, column $S$), where for rank $0$, $S=L(E,1)\,\#E(\Q)_{\mathrm{tors}}^2/(\Omega\prod_pc_p)$, with $\Omega$ Cremona's real period ($L(E,1)/\Omega=1/5$ in that row), is defined as the number that the Birch and
Swinnerton-Dyer conjecture predicts to be $\#\Sha$; for rank $0$ it is exact, but it is a statement about $L(E,1)$, not a proof about $\Sha$, and the table does not claim
otherwise. That $\Sha$ is finite for a curve of analytic rank $0$ is a theorem of Kolyvagin \cite{Kolyvagin90} (with the modularity of $E$ and Gross and Zagier), but neither the table nor that finiteness gives $\Sha[5]=0$; here $\bar\rho_{E,5}$ is
reducible ($E$ has a rational point of order $5$), and a proof needs a separate argument, for instance a descent by the $5$-isogenies of the class, which we do not give. Granting $\Sha[5]=0$,
$\Sel_5(E/\Q)\simeq E(\Q)/5E(\Q)\simeq\Z/5$. Over $\Q_{11}$, $E[5]\subset E(\Q_{11})$ (\Cref{ex:galcoh:local-orders}(c)), so the local condition at $11$ is a subgroup of order
$\#E(\Q_{11})[5]=25$ (\Cref{prop:galcoh:ec-local-image}(i)) of $H^1(\Q_{11},E[5])$, which has order $625$ (\Cref{ex:galcoh:local-orders}(c)). It is not the unramified subgroup, although that has order $25$
too. Indeed $E$ is a Tate curve over $\Q_{11}$ ($a_{11}=1$), $E(\Q_{11})\simeq\Q_{11}^\times/q_E^{\Z}$, and as $q_E$ is a fifth power,
$E(\Q_{11})/5\simeq\Q_{11}^\times/(\Q_{11}^\times)^5$; the Kummer map of $E$ restricted to this is the Kummer map of $\mathbb G_m$ followed by $\mu_5\subset E[5]$. So the local
condition is $H^1(\Q_{11},\mu_5)=\Hom(G_{\Q_{11}},\mu_5)\subset\Hom(G_{\Q_{11}},E[5])$ (the action on $E[5]$ is trivial), which contains the ramified class $\delta(11)$. Its
intersection with the unramified subgroup has order $5$, in agreement with \Cref{prop:galcoh:ec-local-image}(iv) and $c_{11}=5$.
\end{example}
\section{Selmer structures and dual Selmer structures}\label{sec:galcoh:structures}

The paper's Selmer groups $H^1_f(\Q,W)$ and $H^1_\Sigma(\Q,W)$ are defined by local conditions: the Bloch-Kato condition at $\ell$, the unramified condition at the
other primes, and, for $H^1_\Sigma$, no condition at the primes of $\Sigma$. This section gives the general language (Selmer structures, after Mazur and Rubin), the
dual structure on the Kummer dual defined by local duality, and the two facts about $\ell$-adic Selmer groups that the proof of Lemma lem:dfg21 needs. The
Bloch-Kato structure itself is defined in \Cref{ch:blochkato}.

From now on $\ell$ is odd, so that every local condition at $\infty$ is $0=H^1(\R,M)$ (\Cref{prop:galcoh:real-place}).

\begin{definition}[Selmer structures]\label{def:galcoh:selmer-structure}
Let $M$ be a $G_\Q$-module unramified outside a finite set $S$ of places (containing $\infty$ and $\ell$), of one of the following kinds: finite of $\ell$-power order; $W$,
discrete and cofinitely generated; or $T$, of type (C). A \emph{Selmer structure} $\mathcal L$ on $M$ is a family of subgroups $\mathcal L_v\subseteq H^1(\Q_v,M)$, one for each place
$v$, closed when $M$ is of type (C), such that $\mathcal L_v=\Hur(\Q_v,M)$ for all $v$ outside a finite set $S_{\mathcal L}\supseteq S$. Its \emph{Selmer group} is
\begin{align*}
H^1_{\mathcal L}(\Q,M)&=\bigl\{c\in H^1(\Q,M):\Res_vc\in\mathcal L_v\text{ for all }v\bigr\}\\
&=\ker\Bigl(H^1(G_{\Q,S_{\mathcal L}},M)\to\bigoplus_{v\in S_{\mathcal L}}H^1(\Q_v,M)/\mathcal L_v\Bigr),
\end{align*}
the equality by \Cref{prop:galcoh:unramified-almost-everywhere}(i). We write $\mathcal L\subseteq\mathcal L'$ if $\mathcal L_v\subseteq\mathcal L'_v$ for all $v$.
\end{definition}

\begin{example}[Selmer structures]\label{ex:galcoh:selmer-structures}
\begin{enumerate}[label=(\alph*)]
\item The \emph{unramified} structure, $\mathcal L_p=\Hur(\Q_p,M)$ at every prime.
\item For a finite set $\Sigma$ of primes and a structure $\mathcal L$: the structure \emph{relaxed} at $\Sigma$, $\mathcal L^\Sigma_p=H^1(\Q_p,M)$ for $p\in\Sigma$ and $\mathcal L^\Sigma_v=\mathcal L_v$
otherwise; and the structure \emph{strict} at $\Sigma$, $\mathcal L_{\Sigma,p}=0$ for $p\in\Sigma$ and $\mathcal L_{\Sigma,v}=\mathcal L_v$ otherwise.
\item The classical structure $\mathcal L_v=\delta_v(E(\Q_v)/nE(\Q_v))$ on $E[n]$, $n=\ell^k$ ($\ell$ odd, as fixed above), a Selmer structure by \Cref{prop:galcoh:ec-local-image}(ii), with Selmer group
$\Sel_n(E/\Q)$.
\item \emph{Propagated structures.} Given a lattice $T$ unramified outside $S$ and $\Q_\ell$-subspaces $L_{V,v}\subseteq H^1(\Q_v,V)$ with $L_{V,v}=\Hur(\Q_v,V)$ for $v\notin S$, the
families $(\iota^{-1}L_{V,v})_v$ on $T$ and $(\pi L_{V,v})_v$ on $W$ (\Cref{prop:galcoh:propagation}) are Selmer structures: at $p\notin S$, $p\ne\ell$, they are the unramified
conditions by \Cref{prop:galcoh:propagation}(iii).
\end{enumerate}
The Bloch-Kato structure of \Cref{ch:blochkato} is of type (d), and the paper's $H^1_f(\Q,W)$ and $H^1_\Sigma(\Q,W)$ are the Selmer groups of the Bloch-Kato structure on $W$
and of its relaxation at $\Sigma$, the latter described in Section 4.2 of the paper as the group ``with no condition at the primes of $\Sigma$, the Bloch-Kato condition at
$\ell$ and the unramified condition elsewhere''.
\end{example}

\begin{definition}[dual Selmer structures]\label{def:galcoh:dual-structure}
(a) Let $M$ be finite of $\ell$-power order and $\mathcal L$ a Selmer structure on $M$. The \emph{dual structure} $\mathcal L^*$ on $M^D$ is $\mathcal L^*_v=\mathcal L_v^\perp$, the
annihilator under the local pairing of \Cref{thm:galcoh:tate-local-duality}(i) (at $v=\infty$ both sides are $0$).
(b) Let $T$ be a lattice unramified outside $S$ and $\mathcal L$ a Selmer structure on $W$. The \emph{dual structure} $\mathcal L^*$ on $T^D$ is $\mathcal L^*_v=\mathcal L_v^\perp\subseteq H^1(\Q_v,T^D)$,
the annihilator under the first pairing of \Cref{thm:galcoh:tate-local-duality}(ii); symmetrically for structures on $T$, with duals on $W^D$.
\end{definition}

\begin{proposition}[properties of dual structures]\label{prop:galcoh:dual-structure}
\begin{enumerate}[label=(\roman*)]
\item The dual of a Selmer structure is a Selmer structure: at every prime $p\notin S_{\mathcal L}$ (so $p\ne\ell$), $\Hur(\Q_p,M)^\perp=\Hur(\Q_p,M^D)$, and likewise for
$(W,T^D)$.
\item $\mathcal L^{**}=\mathcal L$; $\mathcal L\subseteq\mathcal L'$ implies $\mathcal L'^*\subseteq\mathcal L^*$; and $(\mathcal L^\Sigma)^*=(\mathcal L^*)_\Sigma$: relaxed and strict are dual.
\item If $\mathcal L$ on $W$ is propagated from $(L_{V,v})_v$ and $L_{V,v}^\perp=L_{V^D,v}$ for every $v$, then $\mathcal L^*$ is the structure on $T^D$ propagated from
$(L_{V^D,v})_v$.
\end{enumerate}
\end{proposition}

\begin{proof}
(i) For finite $M$ this is \Cref{thm:galcoh:ur-orthogonal}(i). For $(W,T^D)$ and $p\notin S$, $T$ and $T^D$ are unramified at $p$, so by
\Cref{prop:galcoh:propagation}(iii) $\Hur(\Q_p,W)$ and $\Hur(\Q_p,T^D)$ are propagated from $\Hur(\Q_p,V)$ and $\Hur(\Q_p,V^D)$, which are exact annihilators by
\Cref{thm:galcoh:ur-orthogonal}(ii); apply \Cref{prop:galcoh:propagation}(ii).
(ii) For a perfect pairing between a finite (or discrete) group and its dual, $(A^\perp)^\perp=A$ and $A\subseteq A'$ implies $A'^\perp\subseteq A^\perp$; and $H^1(\Q_p,M)^\perp=0$.
(iii) is \Cref{prop:galcoh:propagation}(ii) at each place.
\end{proof}

\begin{proposition}[$\ell$-adic Selmer groups]\label{prop:galcoh:selmer-l-adic}
Let $T$ be unramified outside $S$ and $\mathcal L$ the structures on $T$ and $W$ propagated from $(L_{V,v})_v$, and let
$H^1_{L_V}(\Q,V)=\{c\in H^1(G_{\Q,S},V):\Res_vc\in L_{V,v}\ \forall v\}$.
\begin{enumerate}[label=(\roman*)]
\item $H^1_{\mathcal L}(\Q,W)$ is cofinitely generated and $H^1_{\mathcal L}(\Q,T)$ is finitely generated.
\item If $H^1_{L_V}(\Q,V)=0$, then $H^1_{\mathcal L}(\Q,T)=H^1(G_{\Q,S},T)_{\tors}\simeq H^0(\Q,W)/H^0(\Q,W)_{\mathrm{div}}$, and $H^1_{\mathcal L}(\Q,T)\simeq H^0(\Q,W)$ if moreover
$H^0(\Q,V)=0$.
\item (Stated.) $\operatorname{corank}_{\Z_\ell}H^1_{\mathcal L}(\Q,W)=\dim_{\Q_\ell}H^1_{L_V}(\Q,V)$; in particular $H^1_{\mathcal L}(\Q,W)$ is finite if and only if $H^1_{L_V}(\Q,V)=0$.
\end{enumerate}
\end{proposition}

\begin{proof}
(i) $H^1_{\mathcal L}(\Q,W)\subseteq H^1(G_{\Q,S},W)$, which is cofinitely generated by \Cref{thm:galcoh:gqs-finite}, and a subgroup of a cofinitely generated group is
cofinitely generated (its Pontryagin dual is a quotient of a finitely generated $\Z_\ell$-module). $H^1_{\mathcal L}(\Q,T)$ is a submodule of the finitely generated
$H^1(G_{\Q,S},T)$, and $\Z_\ell$ is noetherian.
(ii) Let $c\in H^1_{\mathcal L}(\Q,T)$; then $c\in H^1(G_{\Q,S},T)$ (the conditions outside $S$ are unramified, \Cref{ex:galcoh:selmer-structures}(d) and
\Cref{prop:galcoh:unramified-almost-everywhere}(i)), and $\Res_v(\iota c)=\iota(\Res_vc)\in L_{V,v}$ for all $v$, so $\iota c\in H^1_{L_V}(\Q,V)=0$, and $c$ is torsion by
\Cref{prop:galcoh:tvw}(ii). Conversely, if $c\in H^1(G_{\Q,S},T)$ is torsion, so is each $\Res_vc$, which therefore lies in $\ker\iota\subseteq\iota^{-1}L_{V,v}$
(\Cref{prop:galcoh:propagation}(i)). The isomorphisms are \Cref{prop:galcoh:tvw}(ii) for $G=G_{\Q,S}$ (\Cref{thm:galcoh:gqs-finite}).
(iii) The image of $H^1_{\mathcal L}(\Q,T)\otimes\Q_\ell/\Z_\ell$ in $H^1_{\mathcal L}(\Q,W)$ has finite index, and $H^1_{\mathcal L}(\Q,T)\otimes\Q_\ell=H^1_{L_V}(\Q,V)$; we do not give the details,
see \cite[Section 1.5]{Rubin}.
\end{proof}

\begin{paperlink}\label{pl:galcoh:structures}
In Theorem thm:main the group $H^1_f(\Q,\ad T\otimes\Q_\ell/\Z_\ell)$ is the Selmer group of the Bloch-Kato structure on $W$, and in Section 4.2 of the paper
$H^1_\Sigma(\Q,W)\supseteq H^1_f(\Q,W)$ is the Selmer group of its relaxation at $\Sigma$ (\Cref{ex:galcoh:selmer-structures}(b), (d)). Lemma lem:dfg21 puts
$H^1_f(\Q_p,T^D)=\iota^{-1}H^1_f(\Q_p,V^D)$, the dual structure on $T^D$ in the sense of \Cref{prop:galcoh:dual-structure}(iii), and its hypothesis $H^1_f(\Q,V^D)=0$ is
the hypothesis of \Cref{prop:galcoh:selmer-l-adic}(ii) for $T^D$; that proposition then says (as $H^0(\Q,V^D)=0$) that the Selmer group $H^1_f(\Q,T^D)$ is $H^0(\Q,W^D)$, the group whose dual is
the last term of the lemma.
\end{paperlink}
\section[Global duality]{Global duality: Poitou-Tate and the formula of Greenberg and Wiles}\label{sec:galcoh:global-duality}

Global duality relates a Selmer group to the Selmer group of the dual structure on the Kummer dual. In the paper it enters only through Lemma lem:dfg21, the
exact sequence of Diamond, Flach and Guo, which is a comparison of two Selmer structures; this section states the Poitou-Tate sequence and the formula of Greenberg
and Wiles, proves the reciprocity law that makes the comparison maps well defined, and checks the formula by hand in the simplest case, where it recovers a
piece of class field theory over $\Q$. Throughout, $\ell$ is odd, $M$ is finite of $\ell$-power order, and $M^D$ is its Kummer dual.

\begin{theorem}[the Poitou-Tate exact sequence]\label{thm:galcoh:poitou-tate}
Let $M$ be unramified outside $S$, with $\infty,\ell\in S$, and $G_S=G_{\Q,S}$. There is an exact sequence
\[
\begin{array}{l}
0\to H^0(G_S,M)\to\bigoplus_{v\in S}\hat H^0(\Q_v,M)\to H^2(G_S,M^D)^*\\[2pt]
\qquad\to H^1(G_S,M)\to\bigoplus_{v\in S}H^1(\Q_v,M)\to H^1(G_S,M^D)^*\\[2pt]
\qquad\to H^2(G_S,M)\to\bigoplus_{v\in S}H^2(\Q_v,M)\to H^0(G_S,M^D)^*\to0,
\end{array}
\]
where $\hat H^0(\Q_p,M)=H^0(\Q_p,M)$ for primes $p$, $\hat H^0(\R,M)=0$ ($\ell$ odd), the maps $\bigoplus_vH^i(\Q_v,M)\to H^{2-i}(G_S,M^D)^*$ send $(c_v)_v$ to
$y\mapsto\sum_v\operatorname{inv}_v(c_v\cup\Res_vy)$, and the other maps are localisations or their duals. In particular the image of $H^1(G_S,M)$ in $\bigoplus_{v\in S}H^1(\Q_v,M)$ is the
exact annihilator of the image of $H^1(G_S,M^D)$ in $\bigoplus_{v\in S}H^1(\Q_v,M^D)$ under the sum of the local pairings.
\end{theorem}

\noindent We do not prove this; see \cite[Chapter I, Section 4]{MilneADT}, \cite[Chapter VIII, Section 6]{NSW}, \cite{TateDuality} and \cite[Section 1.7]{Rubin}. The
last sentence is exactness at $\bigoplus_{v\in S}H^1(\Q_v,M)$: the kernel of the map to $H^1(G_S,M^D)^*$ is by definition the annihilator of the image of $H^1(G_S,M^D)$, and
the annihilator of that annihilator is the image itself, the local pairings being perfect.

\begin{proposition}[reciprocity]\label{prop:galcoh:reciprocity}
For $x\in H^1(\Q,M)$ and $y\in H^1(\Q,M^D)$, the local cup products $\Res_vx\cup\Res_vy$ are zero for all but finitely many $v$, and
\[
\sum_v\operatorname{inv}_v\bigl(\Res_vx\cup\Res_vy\bigr)=0 .
\]
\end{proposition}

\begin{proof}
By \Cref{prop:galcoh:unramified-almost-everywhere}(ii), for all but finitely many primes $p$ (say $p\ne\ell$, with $M$, $x$ and $y$ unramified at $p$) both $\Res_px$ and $\Res_py$
are inflated from $G_{\Q_p}/I_p$, so their cup product, pushed to $\mu_{\ell^\infty}$, is inflated from $H^2(\hat\Z,\mu_{\ell^\infty})=0$, as in the proof of
\Cref{thm:galcoh:ur-orthogonal}. The cup product $x\cup y$, pushed to $\mu_{\ell^\infty}$, is a class in $H^2(\Q,\mu_{\ell^k})$ for some $k$, that is, by
\Cref{thm:galcoh:kummer}(ii), an element $\beta$ of $\operatorname{Br}(\Q)[\ell^k]$. Cup products commute with restriction, and the isomorphism of \Cref{thm:galcoh:kummer}(ii) commutes
with restriction (it comes from a connecting map), so the local components of $\beta$ are the classes $\Res_vx\cup\Res_vy$. The sum of the invariants of the local
components of an element of $\operatorname{Br}(\Q)$ is $0$ by \Cref{thm:galcoh:brauer}(iii).
\end{proof}

\begin{theorem}[the formula of Greenberg and Wiles]\label{thm:galcoh:greenberg-wiles}
Let $\mathcal L$ be a Selmer structure on $M$ and $\mathcal L^*$ its dual on $M^D$. Then $H^1_{\mathcal L}(\Q,M)$ and $H^1_{\mathcal L^*}(\Q,M^D)$ are finite, and
\[
\frac{\#H^1_{\mathcal L}(\Q,M)}{\#H^1_{\mathcal L^*}(\Q,M^D)}=\frac{\#H^0(\Q,M)}{\#H^0(\Q,M^D)}\prod_v\frac{\#\mathcal L_v}{\#H^0(\Q_v,M)},
\]
the product over all places including $\infty$; all but finitely many factors are $1$.
\end{theorem}

\noindent The finiteness follows from \Cref{def:galcoh:selmer-structure} and \Cref{thm:galcoh:gqs-finite}(i), and almost all factors are $1$ by
\Cref{prop:galcoh:unramified-computed}(ii). We do not prove the formula; it is due to Greenberg and to Wiles, see \cite[Section 1]{Wiles95} and
\cite[Section 2]{DDT}. The proof combines \Cref{thm:galcoh:poitou-tate} with Tate's global Euler characteristic formula
$\#H^0(G_S,M)\,\#H^2(G_S,M)/\#H^1(G_S,M)=\#H^0(\R,M)/\abs{\#M}_\infty$ (stated; \cite[Chapter I, Section 5]{MilneADT}), which supplies the factor $1/\#H^0(\R,M)$ at
infinity.

\begin{example}[the formula by hand]\label{ex:galcoh:gw-check}
Let $M=\Z/\ell$, so $M^D=\mu_\ell$.
\begin{enumerate}[label=(\alph*)]
\item Let $\mathcal L$ be unramified at every prime, $\ell$ included, and $0$ at $\infty$. Then $H^1_{\mathcal L}(\Q,\Z/\ell)$ consists of the homomorphisms $G_\Q\to\Z/\ell$ unramified
everywhere, which are $0$ because $\Q$ has no nontrivial unramified extension (Minkowski; stated). The dual structure: $\mathcal L^*_p=\Hur(\Q_p,\mu_\ell)$ for $p\ne\ell$
(\Cref{thm:galcoh:ur-orthogonal}), and at $\ell$ the annihilator of $\Hur(\Q_\ell,\Z/\ell)$, which is the image of $\Z_\ell^\times$ in $\Q_\ell^\times/(\Q_\ell^\times)^\ell$
(\Cref{warn:galcoh:p-equals-l}; this uses local class field theory). By \Cref{ex:galcoh:global-kummer} and these conditions,
$H^1_{\mathcal L^*}(\Q,\mu_\ell)=\{x\in\Q^\times:\ell\mid v_p(x)\text{ for all }p\}/(\Q^\times)^\ell=\{\pm1\}(\Q^\times)^\ell/(\Q^\times)^\ell=0$, as $-1=(-1)^\ell$. The left side is $1$. On
the right, $\#H^0(\Q,\Z/\ell)/\#H^0(\Q,\mu_\ell)=\ell$; the factor at each prime is $\#\Hur/\#H^0=1$ (\Cref{prop:galcoh:unramified-computed}(ii), valid at $\ell$ too); and the
factor at $\infty$ is $\#\mathcal L_\infty/\#H^0(\R,\Z/\ell)=1/\ell$. The product is $1$.
\item Relax $\mathcal L$ at one prime $p\ne\ell$ and put $m=1$ if $p\equiv1\pmod\ell$, $m=0$ otherwise. The dual structure becomes strict at $p$: $x$ with $\ell\mid v_q(x)$
for all primes $q$ and $x\in(\Q_p^\times)^\ell$; then $x=\pm y^\ell$ and the dual Selmer group is still $0$. On the right only the factor at $p$ changes, to
$\#H^1(\Q_p,\Z/\ell)/\#H^0(\Q_p,\Z/\ell)=\ell^{1+m}/\ell=\ell^m$ (\Cref{ex:galcoh:local-orders}(a)). So the formula says: the homomorphisms $G_\Q\to\Z/\ell$ unramified outside
$p$ (and at $\ell$, and at $\infty$) number $\ell^m$. This is what class field theory gives: they factor through $\Gal(\Q(\zeta_p)/\Q)\simeq(\Z/p)^\times$ by the theorem
of Kronecker and Weber, and $\Hom((\Z/p)^\times,\Z/\ell)$ has order $\ell^m$. The factor at $\infty$ is essential: without it the two sides would differ by $\ell$.
\end{enumerate}
\end{example}

\begin{corollary}[comparison of two structures]\label{cor:galcoh:gw-comparison}
For Selmer structures $\mathcal L\subseteq\mathcal L'$ on $M$,
\[
\frac{\#H^1_{\mathcal L'}(\Q,M)}{\#H^1_{\mathcal L}(\Q,M)}=\prod_v\frac{\#\mathcal L'_v}{\#\mathcal L_v}\cdot\frac{\#H^1_{\mathcal L'^*}(\Q,M^D)}{\#H^1_{\mathcal L^*}(\Q,M^D)} .
\]
\end{corollary}

\begin{proof}
Divide the formula of \Cref{thm:galcoh:greenberg-wiles} for $\mathcal L'$ by the one for $\mathcal L$; the global terms $\#H^0$ and the local terms $\#H^0(\Q_v,M)$ cancel.
\end{proof}

\noindent This is the numerical shadow of the exact sequence of \Cref{thm:galcoh:comparison}; applied at each finite level to a structure and its relaxation at $\Sigma$,
and passed to the limit, it gives the length identity in the proof of Theorem thm:main.
\section[Comparing Selmer structures]{Comparing Selmer structures; the abstract form of Lemma 2.1 of Diamond, Flach and Guo}\label{sec:galcoh:comparison}

This section proves the exact sequence that compares the Selmer groups of two structures $\mathcal L\subseteq\mathcal L'$, first for finite modules and then for the pair
$(W,T^D)$, and specialises it to the relaxation of a structure at a finite set $\Sigma$. The result, \Cref{thm:galcoh:relaxing-sigma}, is Lemma lem:dfg21 of the paper
with an abstract local condition in place of the Bloch-Kato one, and \Cref{cor:galcoh:length-identity} is the length identity used in the proof of Theorem thm:main and
in Remark rem:ns. \Cref{ch:blochkato} supplies what is specific to the Bloch-Kato structure (\Cref{rem:galcoh:handover-ch18}). Throughout, $\ell$ is odd.

\begin{theorem}[the comparison sequence]\label{thm:galcoh:comparison}
Let $M$ be finite of $\ell$-power order, unramified outside $S$, and $\mathcal L\subseteq\mathcal L'$ Selmer structures on $M$, so that $\mathcal L'^*\subseteq\mathcal L^*$ on $M^D$. Consider
\[
0\to H^1_{\mathcal L}(\Q,M)\to H^1_{\mathcal L'}(\Q,M)\xrightarrow{\ \operatorname{loc}\ }\bigoplus_v\mathcal L'_v/\mathcal L_v\xrightarrow{\ \theta\ }H^1_{\mathcal L^*}(\Q,M^D)^*\xrightarrow{\ \Res^*\ }H^1_{\mathcal L'^*}(\Q,M^D)^*\to0,
\]
where $\operatorname{loc}$ is localisation, $\theta\bigl((c_v)_v\bigr)(y)=\sum_v\operatorname{inv}_v(c_v\cup\Res_vy)$, and $\Res^*$ is the dual of the inclusion
$H^1_{\mathcal L'^*}\subseteq H^1_{\mathcal L^*}$. The sum is finite ($\mathcal L'_v=\mathcal L_v$ for almost all $v$), and $\theta$ is well defined because $\Res_vy\in\mathcal L^*_v=\mathcal L_v^\perp$.
\begin{enumerate}[label=(\roman*)]
\item The sequence is exact at $H^1_{\mathcal L}$ and at $H^1_{\mathcal L'}$, $\Res^*$ is surjective, and it is a complex: $\theta\circ\operatorname{loc}=0$ and $\Res^*\circ\theta=0$. In
particular $\#H^1_{\mathcal L'}(\Q,M)/\#H^1_{\mathcal L}(\Q,M)\le\prod_v\#(\mathcal L'_v/\mathcal L_v)$.
\item The sequence is exact.
\end{enumerate}
\end{theorem}

\begin{proof}
(i) The kernel of $\operatorname{loc}$ consists of the classes of $H^1_{\mathcal L'}$ whose localisations lie in $\mathcal L_v$ for all $v$, which is $H^1_{\mathcal L}$. $\Res^*$ is the dual of an
injection of finite groups, hence surjective. For $x\in H^1_{\mathcal L'}(\Q,M)$ and $y\in H^1_{\mathcal L^*}(\Q,M^D)$,
\[
\theta(\operatorname{loc}x)(y)=\sum_{v:\,\mathcal L'_v\ne\mathcal L_v}\operatorname{inv}_v(\Res_vx\cup\Res_vy)=\sum_v\operatorname{inv}_v(\Res_vx\cup\Res_vy)=0,
\]
the middle equality because at the other places $\Res_vx\in\mathcal L_v$ and $\Res_vy\in\mathcal L_v^\perp$, the last by \Cref{prop:galcoh:reciprocity}. For $y\in H^1_{\mathcal L'^*}$ and
$c_v\in\mathcal L'_v$, $c_v\cup\Res_vy=0$ because $\Res_vy\in\mathcal L'^*_v=\mathcal L_v'^\perp$; so $\Res^*\circ\theta=0$. The inequality is the injectivity of
$H^1_{\mathcal L'}/H^1_{\mathcal L}\to\bigoplus_v\mathcal L'_v/\mathcal L_v$.

(ii) \emph{Exactness at the middle term.} Enlarge $S$ so that $\mathcal L_v=\mathcal L'_v=\Hur(\Q_v,M)$ and $\mathcal L^*_v=\Hur(\Q_v,M^D)$ for $v\notin S$. In $P=\bigoplus_{v\in S}H^1(\Q_v,M)$ let $B$ be
the image of $H^1(G_S,M)$ and $A=\bigoplus_{v\in S}\mathcal L_v$; in $P'=\bigoplus_{v\in S}H^1(\Q_v,M^D)$ let $B'$ be the image of $H^1(G_S,M^D)$. By \Cref{thm:galcoh:poitou-tate},
$B^\perp=B'$, and $A^\perp=\bigoplus_{v\in S}\mathcal L^*_v$. Hence
\[
(A+B)^\perp=A^\perp\cap B'=\operatorname{loc}\bigl(H^1_{\mathcal L^*}(\Q,M^D)\bigr),
\]
since a class of $H^1(G_S,M^D)$ with localisations in $\mathcal L^*_v$ for $v\in S$ is exactly a class of $H^1_{\mathcal L^*}$ (\Cref{prop:galcoh:unramified-almost-everywhere}(i)). Now let
$(c_v)\in\bigoplus_v\mathcal L'_v/\mathcal L_v$ with $\theta((c_v))=0$, and lift it to $c=(c_v)_{v\in S}\in\bigoplus_{v\in S}\mathcal L'_v\subseteq P$. Then $c$ is orthogonal to
$(A+B)^\perp$, so $c\in(A+B)^{\perp\perp}=A+B$: $c=a+\operatorname{loc}(x)$ with $a\in A$ and $x\in H^1(G_S,M)$. Then $\Res_vx=c_v-a_v\in\mathcal L'_v$ for $v\in S$, and $x$ is
unramified outside $S$, so $x\in H^1_{\mathcal L'}(\Q,M)$ and $\operatorname{loc}(x)\equiv(c_v)$ modulo $\bigoplus\mathcal L_v$.

\emph{Exactness at $H^1_{\mathcal L^*}(\Q,M^D)^*$.} By exactness at the middle term and at $H^1_{\mathcal L'}$,
\[
\#\operatorname{im}\theta=\frac{\prod_v\#(\mathcal L'_v/\mathcal L_v)}{\#\operatorname{im}(\operatorname{loc})}=\prod_v\frac{\#\mathcal L'_v}{\#\mathcal L_v}\cdot\frac{\#H^1_{\mathcal L}(\Q,M)}{\#H^1_{\mathcal L'}(\Q,M)}
=\frac{\#H^1_{\mathcal L^*}(\Q,M^D)}{\#H^1_{\mathcal L'^*}(\Q,M^D)}=\#\ker\Res^*,
\]
the third equality by \Cref{cor:galcoh:gw-comparison}. As $\operatorname{im}\theta\subseteq\ker\Res^*$, they are equal.
\end{proof}

\noindent So the comparison sequence is proved here given \Cref{thm:galcoh:poitou-tate} (for the middle) and \Cref{thm:galcoh:greenberg-wiles} (for the fourth term); only the
complex property (i) is elementary. It can also be read off directly from the nine-term sequence; see \cite[Section 1.7]{Rubin} and \cite[Section 2.3]{MazurRubin04}.

For the pair $(W,T^D)$ we use the dual structures of \Cref{def:galcoh:dual-structure}(b).

\begin{theorem}[the comparison sequence for $W$ and $T^D$]\label{thm:galcoh:comparison-l-adic}
Let $T$ be a lattice unramified outside $S$, $\mathcal L\subseteq\mathcal L'$ Selmer structures on $W$, and $\mathcal L^*\supseteq\mathcal L'^*$ their dual structures on $T^D$. Then
\[
0\to H^1_{\mathcal L}(\Q,W)\to H^1_{\mathcal L'}(\Q,W)\to\bigoplus_v\mathcal L'_v/\mathcal L_v\to H^1_{\mathcal L^*}(\Q,T^D)^*\to H^1_{\mathcal L'^*}(\Q,T^D)^*\to0
\]
is exact, with maps defined as in \Cref{thm:galcoh:comparison}, and $\mathcal L'_v/\mathcal L_v\simeq(\mathcal L^*_v/\mathcal L'^*_v)^*$ by local duality.
\end{theorem}

\begin{proof}[Sketch of proof]
For each $k$ give $W[\ell^k]=T/\ell^kT$ the structures $\mathcal L_{(k),v}=$ preimage of $\mathcal L_v$ under $H^1(\Q_v,W[\ell^k])\to H^1(\Q_v,W)$, and $T^D/\ell^kT^D=W[\ell^k]^D$ the
structures given by the images of $\mathcal L^*_v$ under $H^1(\Q_v,T^D)\to H^1(\Q_v,T^D/\ell^kT^D)$; by the compatibility of local duality with these maps
(\Cref{thm:galcoh:tate-local-duality}) they are mutually dual, and likewise for $\mathcal L'$. Apply \Cref{thm:galcoh:comparison} at each level. Take the direct limit over $k$ for the terms on
$W$: direct limits are exact, the Selmer groups of $W[\ell^k]$ tend to those of $W$ (the error terms $H^0(\Q,W)/\ell^k$ of \Cref{prop:galcoh:tvw}(iv) die in the limit, $H^0(\Q,W)$
being torsion), and the local quotients tend to $\mathcal L'_v/\mathcal L_v$. Take the inverse limit for the terms on $T^D$: they are finite groups, so the limit is exact (Mittag-Leffler),
and $\varprojlim_kH^1_{\mathcal L^*_{(k)}}(\Q,T^D/\ell^k)=H^1_{\mathcal L^*}(\Q,T^D)$ by \Cref{ch:cohomology} (Tate's theorem on limits). Pontryagin duality exchanges the two limits. The
local isomorphism is \Cref{prop:galcoh:propagation}(ii) or, in general, the perfectness of the local pairing. See \cite[Section 1.7 and Appendix B]{Rubin},
\cite[Section 4]{Greenberg89} and the proof of \cite[Lemma 2.1]{DFG04}.
\end{proof}

\begin{theorem}[relaxing at $\Sigma$: the abstract form of Lemma 2.1 of Diamond, Flach and Guo]\label{thm:galcoh:relaxing-sigma}
Let $\ell$ be odd and $T$ a lattice unramified outside $S$. Let $(L_{V,v})_v$ be $\Q_\ell$-subspaces $L_{V,v}\subseteq H^1(\Q_v,V)$ with $L_{V,v}=\Hur(\Q_v,V)$ for $v\notin S$, and
$(L_{V^D,v})_v$ subspaces of $H^1(\Q_v,V^D)$ with $L_{V,v}^\perp=L_{V^D,v}$ for every $v$. Let $\mathcal L$ be the structure on $W$ propagated from $(L_{V,v})$ and $\mathcal L^*$ the structure
on $T^D$ propagated from $(L_{V^D,v})$, so $\mathcal L^*_p=\iota^{-1}(L_{V^D,p})$; by \Cref{prop:galcoh:dual-structure}(iii), $\mathcal L^*$ is the dual of $\mathcal L$. Let $\Sigma\subseteq S$ be a finite
set of primes $p\ne\ell$.
\begin{enumerate}[label=(\roman*)]
\item There is an exact sequence
\[
0\to H^1_{\mathcal L}(\Q,W)\to H^1_{\mathcal L^\Sigma}(\Q,W)\to\bigoplus_{p\in\Sigma}\bigl(\mathcal L^*_p\bigr)^*\to H^1_{\mathcal L^*}(\Q,T^D)^*\to H^1_{(\mathcal L^*)_\Sigma}(\Q,T^D)^*\to0 .
\]
\item Assume (a) $H^1_{L_{V^D}}(\Q,V^D)=0$. Then
\[
H^1_{\mathcal L^*}(\Q,T^D)=H^1(G_{\Q,S},T^D)_{\tors}\simeq H^0(\Q,W^D)/H^0(\Q,W^D)_{\mathrm{div}}.
\]
\item Assume (a) and (b) $H^0(\Q_p,V^D)=0$ for at least one $p\in\Sigma$. Then $H^0(\Q,V^D)=0$, $H^0(\Q,W^D)$ is finite, $H^1_{(\mathcal L^*)_\Sigma}(\Q,T^D)=0$, and there is an exact
sequence
\[
0\to H^1_{\mathcal L}(\Q,W)\to H^1_{\mathcal L^\Sigma}(\Q,W)\to\bigoplus_{p\in\Sigma\cup\{\infty\}}\bigl(\mathcal L^*_p\bigr)^*\to H^0(\Q,W^D)^*\to0,
\]
where $\mathcal L^*_\infty=H^1(\R,T^D)=0$.
\end{enumerate}
\end{theorem}

\begin{proof}
(i) Apply \Cref{thm:galcoh:comparison-l-adic} with $\mathcal L'=\mathcal L^\Sigma$. By \Cref{prop:galcoh:dual-structure}(ii) its dual is $(\mathcal L^*)_\Sigma$. The local quotients are nonzero only at
$p\in\Sigma$, where $\mathcal L^\Sigma_p/\mathcal L_p=H^1(\Q_p,W)/\mathcal L_p\simeq(\mathcal L^*_p)^*$ by \Cref{prop:galcoh:propagation}(ii).

(ii) By \Cref{prop:galcoh:selmer-l-adic}(ii), applied to $T^D$ and the structure propagated from $(L_{V^D,v})$, hypothesis (a) gives
$H^1_{\mathcal L^*}(\Q,T^D)=H^1(G_{\Q,S},T^D)_{\tors}\simeq H^0(\Q,W^D)/H^0(\Q,W^D)_{\mathrm{div}}$.

(iii) $H^0(\Q,V^D)\subseteq H^0(\Q_p,V^D)=0$ for the prime $p$ of (b). Then $H^0(\Q,W^D)$ is finite and $\delta\colon H^0(\Q,W^D)\to H^1(G_{\Q,S},T^D)_{\tors}=H^1_{\mathcal L^*}(\Q,T^D)$ is an
isomorphism (\Cref{prop:galcoh:tvw}(ii)). The strict Selmer group $H^1_{(\mathcal L^*)_\Sigma}(\Q,T^D)$ is the subgroup of the classes $\delta(x)$ with $\Res_p\delta(x)=0$ for all
$p\in\Sigma$. By naturality of connecting maps, $\Res_p\delta(x)=\delta_p(\Res_px)$, where $\delta_p\colon H^0(\Q_p,W^D)\to H^1(\Q_p,T^D)$ is the local connecting map and $\Res_px$ is $x$
viewed in $H^0(\Q_p,W^D)\supseteq H^0(\Q,W^D)$. At the prime $p$ of (b), $\delta_p$ is injective, its kernel being the image of $H^0(\Q_p,V^D)=0$ (\Cref{prop:galcoh:tvw}(ii) for
$G_{\Q_p}$). So $\delta(x)$ in the strict group forces $x=0$, and the strict group is $0$. The sequence of (i) becomes the one displayed, after identifying
$H^1_{\mathcal L^*}(\Q,T^D)^*$ with $H^0(\Q,W^D)^*$ through $\delta$; the summand at $\infty$ is $0$ by \Cref{prop:galcoh:real-place}.
\end{proof}

\begin{remark}
Without (b) the conclusion of (iii) can fail only through the last term: the sequence of (i) is still exact, with $H^1_{\mathcal L^*}(\Q,T^D)$ as in (ii), and its last term is the dual
of the group of $\delta(x)$, $x\in H^0(\Q,W^D)$, with $x\in H^0(\Q_p,W^D)_{\mathrm{div}}$ for every $p\in\Sigma$ (the kernel of $\delta_p$ is the image of $H^0(\Q_p,V^D)$, which is the
divisible part by \Cref{prop:galcoh:tvw}(ii)).
\end{remark}

\begin{corollary}[the length identity]\label{cor:galcoh:length-identity}
Under the hypotheses (a) and (b) of \Cref{thm:galcoh:relaxing-sigma},
\[
\len H^1_{\mathcal L^\Sigma}(\Q,W)=\len H^1_{\mathcal L}(\Q,W)+\sum_{p\in\Sigma}\len\mathcal L^*_p-\len H^0(\Q,W^D),
\]
both sides being finite if and only if $H^1_{\mathcal L^\Sigma}(\Q,W)$ is finite, and infinite otherwise. If $L_{V^D,p}=0$ for $p\in\Sigma$ (as in the paper), every $\mathcal L^*_p=H^1(\Q_p,T^D)_{\tors}$ is
finite, and the identity holds with all terms finite as soon as $H^1_{\mathcal L}(\Q,W)$ is finite.
\end{corollary}

\begin{proof}
In an exact sequence $0\to A\to B\to C\to D\to0$ of $\Z_\ell$-modules with $D$ finite, $B$ is finite if and only if $A$ and $C$ are (the image of $B$ in $C$ has finite index), and then
$\len B=\len A+\len C-\len D$ (\Cref{ch:cohomology}, lengths in exact sequences). Pontryagin duality preserves lengths of finite modules (\Cref{ch:modules}), so
$\len(\mathcal L^*_p)^*=\len\mathcal L^*_p$ and $\len H^0(\Q,W^D)^*=\len H^0(\Q,W^D)$. If $L_{V^D,p}=0$, $\mathcal L^*_p=\ker\iota=H^1(\Q_p,T^D)_{\tors}$, finite by \Cref{prop:galcoh:tvw}(ii) and
\Cref{thm:galcoh:local-finiteness}.
\end{proof}

\begin{paperlink}\label{pl:galcoh:dfg21}
Lemma lem:dfg21 of the paper (\cite[Lemma 2.1]{DFG04}) is \Cref{thm:galcoh:relaxing-sigma}(iii) for the Bloch-Kato structure: $L_{V,v}=H^1_f(\Q_v,V)$, $L_{V^D,v}=H^1_f(\Q_v,V^D)$,
$\mathcal L=f$ on $W$, $\mathcal L^\Sigma$ the structure of $H^1_\Sigma(\Q,W)$, and $\mathcal L^*_p=H^1_f(\Q_p,T^D)=\iota^{-1}H^1_f(\Q_p,V^D)$ exactly as the lemma defines it. The hypotheses are discharged as
follows. $L_{V,v}^\perp=L_{V^D,v}$: at $p\ne\ell$ this is \Cref{thm:galcoh:ur-orthogonal}(ii), at $\ell$ it is the theorem of Bloch and Kato (\Cref{ch:blochkato}), and at $\infty$ both
are $0$. Hypothesis (a) is $H^1_f(\Q,V^D)=0$, Theorem thm:dfg27. Hypothesis (b) holds at every $p\in\Sigma$ by the paragraph after the lemma
(\Cref{pl:galcoh:after-dfg21}): $\Hur(\Q_p,V^D)=0$, hence $H^0(\Q_p,V^D)=0$ by \Cref{prop:galcoh:unramified-computed}(iii). The paper's hypothesis $H^0(\Q,V^D)=0$ then
follows, as the revised paper notes (``so the hypothesis $H^0(\Q,V^D)=0$ is automatic here''). In the paper's situation $L_{V^D,p}=0$ for $p\in\Sigma$, so
$\mathcal L^*_p=H^1(\Q_p,T^D)_{\tors}\simeq H^0(\Q_p,W^D)$ (\Cref{cor:galcoh:ur-vanishing}). \emph{Commentary.} The paper states the lemma under the hypothesis
$H^0(\Q,V^D)=H^1_f(\Q,V^D)=0$, quoting DFG04; the book proves the version with hypothesis (b) at one prime of $\Sigma$, which is what the derivation from global duality uses
and what holds in the paper's application.

\emph{Proof of Theorem thm:main, the Selmer step.} In the proof of Theorem thm:main the paper combines Theorem thm:dfg27, Lemma lem:dfg21 and the vanishing
$H^0(\Q,W^D)=0$ of Lemma lem:h0 into the identity $\len H^1_\Sigma(\Q,W)=\len H^1_f(\Q,W)+\sum_{p\in\Sigma}\len H^1_f(\Q_p,T^D)$. This is \Cref{cor:galcoh:length-identity} with
$H^0(\Q,W^D)=0$; the finiteness of $H^1_\Sigma(\Q,W)$ is Theorem thm:dfg37 (\Cref{ch:congruence}).

\emph{Remark rem:ns.} When $\ell=3$ and the image of $\bar\rho$ is $N_s(\F_3)$, Lemma lem:h0 (proved in \Cref{ch:galrep}) gives $\len H^0(\Q,W^D)=h_3(E)=1$; for $\ell\ge5$,
$h_\ell(E)=0$, and $h_3(E)\le1$ always, with or without complex multiplication. If $H^1_f(\Q,V^D)=0$, which the paper assumes for these curves, hypotheses (a) and (b) hold,
and \Cref{cor:galcoh:length-identity} with the local lengths of Lemma lem:fpr (\Cref{ch:blochkato}) is the displayed identity of Remark rem:ns, with $-h_\ell(E)$ on the right.
\end{paperlink}

\begin{example}[a toy case: class field theory of $\Q$ again]\label{ex:galcoh:relaxing-toy}
Let $T=\Z_\ell$ with trivial action, $\ell$ odd: $V=\Q_\ell$, $W=\Q_\ell/\Z_\ell$, $T^D=\Z_\ell(1)$, $W^D=\mu_{\ell^\infty}$. Take $S=\{\infty,\ell,p\}$, $\Sigma=\{p\}$ with $p\ne\ell$, and
$L_{V,q}=\Hur(\Q_q,\Q_\ell)$ at every prime $q$ (including $\ell$), $L_{V,\infty}=0$.
\begin{itemize}
\item \emph{The dual conditions.} For $q\ne\ell$, $L_{V^D,q}=\Hur(\Q_q,\Q_\ell(1))=0$ (\Cref{thm:galcoh:ur-orthogonal}(ii); $\Frob_q$ acts on $\Q_\ell(1)$ by $q\ne1$). At $\ell$,
$H^1(\Q_\ell,\Q_\ell)$ has dimension $2$ and $H^1(\Q_\ell,\Q_\ell(1))=\widehat{\Q_\ell^\times}\otimes\Q_\ell$ has dimension $2$ (\Cref{thm:galcoh:local-euler} with $\ell=p$:
$0-\dim H^1+1=-1$); the annihilator of the line of unramified characters is the line spanned by the unit classes (local class field theory; stated). So
$L_{V^D,\ell}=\Z_\ell^\times\otimes\Q_\ell$.
\item \emph{Hypothesis (a).} $H^1(G_{\Q,S},\Q_\ell(1))=(\Z_S^\times\otimes\Z_\ell)\otimes\Q_\ell$, with $\Z_S^\times=\pm p^{\Z}\ell^{\Z}$ (\Cref{ex:galcoh:global-kummer} at each level). The
condition at $p$ kills the $p$-component (the class of $p$ in $H^1(\Q_p,\Q_\ell(1))$ is nonzero, \Cref{prop:galcoh:kummer-local}(iv)), and the condition at $\ell$ kills the
$\ell$-component (the class of $\ell$ is not a unit class). So $H^1_{L_{V^D}}(\Q,\Q_\ell(1))=0$.
\item \emph{Hypothesis (b).} $H^0(\Q_p,\Q_\ell(1))=0$, as $\Frob_p$ acts by $p\ne1$.
\item \emph{The terms.} $H^1_{\mathcal L}(\Q,W)=\Hom(G_\Q,\Q_\ell/\Z_\ell)$ unramified everywhere, which is $0$ (Minkowski). $H^1_{\mathcal L^\Sigma}(\Q,W)$ is the group of characters of
$\ell$-power order unramified outside $p$. $\mathcal L^*_p=\iota^{-1}(0)=H^1(\Q_p,\Z_\ell(1))_{\tors}=\Z_p^\times\otimes\Z_\ell\simeq\Z/\ell^{v_\ell(p-1)}$ (\Cref{prop:galcoh:kummer-local}).
$H^0(\Q,\mu_{\ell^\infty})=0$.
\end{itemize}
The sequence of \Cref{thm:galcoh:relaxing-sigma}(iii) reads $0\to0\to H^1_{\mathcal L^\Sigma}(\Q,W)\to(\Z_p^\times\otimes\Z_\ell)^*\to0\to0$: the characters of $\ell$-power order of $G_\Q$ unramified
outside $p$ form a group dual to the $\ell$-part of $(\Z/p)^\times$, as the theorem of Kronecker and Weber says. A variant with $H^0(\Q,W^D)\ne0$, the toy model of $h_3(E)=1$,
is \Cref{ex:galcoh:ex-h3}.
\end{example}

\begin{remark}[what \Cref{ch:blochkato} adds]\label{rem:galcoh:handover-ch18}
To obtain the paper's Lemma lem:dfg21 and the proof of Theorem thm:main from \Cref{thm:galcoh:relaxing-sigma} and \Cref{cor:galcoh:length-identity}, \Cref{ch:blochkato} supplies:
the Bloch-Kato local conditions $H^1_f(\Q_v,V)$, with $H^1_f=\Hur$ at $p\ne\ell$ and $H^1_f(\R,-)=0$; their exact orthogonality at $\ell$ (Bloch and Kato); the vanishing
$H^1_f(\Q,V^D)=0$ (Theorem thm:dfg27, \cite[Theorem 2.7]{DFG04}); and the splitting $\len H^1_f(\Q_p,T^D)=\len H^0(\Q_p,W^D)=v_\ell(L_p(A_f,1)^{-1})+v_\ell(\Tam_p(A))$ of Lemma
lem:fpr (\cite[I.4.2.2]{FPR94}), whose Euler factor is \Cref{prop:galcoh:unramified-computed}(iv) and whose Tamagawa factor comes from
\Cref{prop:galcoh:ur-sequence}. \Cref{ch:lfunctions} identifies $\det(1-\Frob_p^{-1}\mid(V^D)^{I_p})$ with $L_p(A_f,1)^{-1}$. This chapter supplies the duality, the local structure
(\Cref{sec:galcoh:unramified}) and the global sequence.
\end{remark}
\section{Exercises}\label{sec:galcoh:exercises}

\begin{exercise}[a twist by an even quadratic character]\label{ex:galcoh:ex-even-twist}
Let $\ell$ be odd, $d>1$ a fundamental discriminant, $\chi_d$ its (even) quadratic character and $T=\Z_\ell(\chi_d)$, so $T^D=\Z_\ell(1)(\chi_d)$. (a) For $p\nmid d\ell$ show
$\Hur(\Q_p,T)=\Z_\ell/(\chi_d(p)-1)$, which is $\Z_\ell$ or $0$, and $\#\Hur(\Q_p,T^D)=\ell^{v_\ell(p\chi_d(p)-1)}$. (b) For $p\mid d$, $p\ne\ell$, show $H^1(\Q_p,T)=0$.
(c) Show $H^0(\Q,W)=0$ and $H^0(\Q,W^D)=0$; contrast with \Cref{ex:galcoh:ex-h3}. (d) Evaluate (a) for $d=5$, $\ell=3$ at $p=2,7,11$.
\emph{Answer:} (a) \Cref{prop:galcoh:unramified-computed}(i), (iv), $\Frob$ acting by $\chi_d(p)$ on $T$ and by $p\chi_d(p)$ on $T^D$. (b) $I_p$ acts by $\chi_d|_{I_p}\ne1$ on $V$ and on
$V^D$, so both have no invariants, $H^1(\Q_p,V)=0$ by \Cref{cor:galcoh:local-duality-consequences}, and $H^1(\Q_p,T)=H^1(\Q_p,T)_{\tors}\simeq H^0(\Q_p,W)\subseteq W[2]=0$
(\Cref{prop:galcoh:tvw}(ii)). (c) $H^0(\Q,W)\subseteq W[2]=0$; an element of order $\ell$ of $W^D$ is fixed only if $\bar\chi_\ell\chi_d=1$, which forces $\bar\chi_\ell$ quadratic, so
$\ell=3$ and $\chi_d=\chi_{-3}$, which is odd. (d) $\chi_5(2)=\chi_5(7)=-1$, $\chi_5(11)=1$: $\Hur(\Q_p,T)=0,0,\Z_3$ and $\#\Hur(\Q_p,T^D)=3^{v_3(3)},3^{v_3(8)},3^{v_3(10)}=3,1,1$.
\end{exercise}

\begin{exercise}\label{ex:galcoh:ex-orders}
Compute $\#H^i(\Q_p,\Z/\ell^k)$ and $\#H^i(\Q_p,\mu_{\ell^k})$, $i=0,1,2$, for $(p,\ell,k)=(13,3,2)$, $(31,5,2)$ and $(17,3,1)$, and check the Euler characteristic and the duality
of \Cref{cor:galcoh:local-duality-consequences}(i). \emph{Answer:} with $m=\min(k,v_\ell(p-1))$ the orders are $\ell^k,\ell^{k+m},\ell^m$ and $\ell^m,\ell^{k+m},\ell^k$; here $m=1,1,0$, giving
$9,27,3$ and $3,27,9$; $25,125,5$ and $5,125,25$; $3,3,1$ and $1,3,3$.
\end{exercise}

\begin{exercise}\label{ex:galcoh:ex-real-two}
For the group $\{1,c\}$ of order $2$ the cohomology is $2$-periodic: $H^{\mathrm{odd}}=\ker(1+c)/(c-1)M$ and $H^{\mathrm{even}}=M^{c=1}/(1+c)M$ in positive degrees (\Cref{ch:cohomology}). Compute
$H^1$ and $H^2$ of $G_\R$ with coefficients in $\Z_2$ (trivial action) and in $\Z_2(1)$, and check the second against $\operatorname{Br}(\R)=\Z/2$ through the limit of the groups
$H^2(\R,\mu_{2^k})$. \emph{Answer:} $H^1(\R,\Z_2)=0$, $H^2(\R,\Z_2)=\Z/2$; $H^1(\R,\Z_2(1))=\Z/2$, $H^2(\R,\Z_2(1))=0$, which is $\varprojlim\operatorname{Br}(\R)[2^k]$ along the maps induced by
$\zeta\mapsto\zeta^2$, multiplication by $2$, hence $0$. So the paper's ``$H^1_f(\R,T^D)=0$'' genuinely needs $\ell$ odd.
\end{exercise}

\begin{exercise}\label{ex:galcoh:ex-q13}
Check the paper's identification $H^1(G,\Z_\ell(1))=\varprojlim\Q_q^\times/(\Q_q^\times)^{\ell^k}$ for $q=13$, $\ell=3$ by computing both sides, and identify the torsion with
$H^0(\Q_{13},\Q_3/\Z_3(1))$ as \Cref{prop:galcoh:tvw}(ii) predicts. \emph{Answer:} $\Z_3\times\Z/3$ on both sides ($v_3(12)=1$); the torsion is the image of the units, of order
$3=\#\mu_{3^\infty}(\Q_{13})$, since $13\equiv1\pmod3$ and $13\not\equiv1\pmod9$.
\end{exercise}

\begin{exercise}\label{ex:galcoh:ex-q37}
Let $q=37$, $\ell=3$, so $v=v_3(36)=2$. For $u=2$, $8$, $6$ compute the class of $u$ in $\F_{37}^\times\otimes\Z_3=\Z_3(1)/(36)\simeq\Z/9$ and its valuation, in the sense of
Lemma lem:tamfactor. \emph{Answer:} the cubes in $\F_{37}^\times$ are $\{1,6,8,10,11,14,23,26,27,29,31,36\}$ and the ninth powers $\{1,6,31,36\}$; so the valuations are $0$, $1$
and $2$, the last because the class of $6$ is $0$ and $v_3(0)$ is read as $v=2$.
\end{exercise}

\begin{exercise}\label{ex:galcoh:ex-uniformiser}
Let $K$ be a finite extension of $\Q_q$, $q\ne\ell$, and $\varpi$, $\varpi'=\varpi w$ two uniformisers ($w$ a unit). Let $\operatorname{pr}$ and $\operatorname{pr}'$ be the retractions
$H^1(K,\Z_\ell(1))\to\Hur(K,\Z_\ell(1))$ attached to them as in \Cref{prop:galcoh:canonical-splitting}. Show that $\operatorname{pr}'[\kappa_x]=\operatorname{pr}[\kappa_x]-v_K(x)[\kappa_w]$.
Deduce that, with $n=v_K(x)$ and $c_1$, $c_2$ the classes of $\operatorname{pr}[\kappa_x]$ and $\operatorname{pr}'[\kappa_x]$ in $\Z/\ell^{v_\ell(\#k_K-1)}$,
$\min(v_\ell(n),v_\ell(\#k_K-1),v_\ell(c_1))=\min(v_\ell(n),v_\ell(\#k_K-1),v_\ell(c_2))$: the quantity of Lemma lem:tamfactor does not depend on the uniformiser. \emph{Hint:} $x=\varpi^nu=\varpi'^n(uw^{-n})$; the two classes differ by a multiple of $n$.
\end{exercise}

\begin{exercise}\label{ex:galcoh:ex-corner}
Let $\varepsilon\colon G\to\{\pm1\}$ be a character, $\chi\colon G\to\Z_\ell^\times$ a character and $\kappa$ a cocycle with values in $\Z_\ell(\chi)$, so $\kappa(gh)=\kappa(g)+\chi(g)\kappa(h)$.
Show that $g\mapsto\sm{\chi\varepsilon(g)}{\varepsilon\kappa(g)}0{\varepsilon(g)}$ is a homomorphism and that $g\mapsto\sm{\chi\varepsilon(g)}{\kappa(g)}0{\varepsilon(g)}$ is one only if $\kappa(gh)=\varepsilon(h)\kappa(g)+\chi(g)\varepsilon(g)\kappa(h)$
for all $g,h$, which fails for a Kummer cocycle when $\varepsilon\ne1$. What changes in $c'$ if the paper's original matrix is used? \emph{Answer:} the upper right entry of the
product is $\chi\varepsilon(g)\varepsilon\kappa(h)+\varepsilon\kappa(g)\varepsilon(h)=\varepsilon(gh)\kappa(gh)$; with the corrected matrix $c'$ is multiplied by $\varepsilon(\Frob)=\pm1$, which changes no valuation.
\end{exercise}

\begin{exercise}\label{ex:galcoh:ex-pari-tate}
(PARI) For \curve{874a1} at $q=19$ and \curve{1066a1} at $q=13$ ($\ell=3$), compute the Tate parameter with \texttt{ellinit(E,O(q\^{}30)).tate}, its valuation $n$ and the reduction
$\bar u$ of its unit part; decide whether $\bar u$ is a cube and, for $q=19$, a ninth power; deduce the extra term $\min(v_3(n),v_3(q-1),v_3(c))$ and compare
$\sum_qt_3(E,q)+\text{extra}$ with $v_3(\degphi)$. \emph{Answer:} \curve{874a1} $=[1,-1,0,-19,-13]$: $n=3$, $\bar u=12$, a cube but not a ninth power, $a_{19}=-1$, extra $1$; the
other bad primes $2$, $23$ are of type $\mathrm I_1$; bound $2$, and $\degphi=144$. \curve{1066a1} $=[1,0,1,-1693,26680]$: $n=3$, $\bar u=5$, a cube mod $13$, extra $1$; the
other bad primes are $2$ ($\mathrm I_7$) and $41$ ($\mathrm I_2$); bound $2$, and $\degphi=504=2^3\cdot3^2\cdot7$. Both have surjective mod $3$ image, and the bound is attained.
\end{exercise}

\begin{exercise}\label{ex:galcoh:ex-good-prime}
Let $E$ have good reduction at $p\ne\ell$, $a=a_p(E)$, $V=\ad V_\ell E$ and $V^D=V(1)$. Show that $\Hur(\Q_p,V^D)=0$ and that
\[
\#\Hur(\Q_p,T^D)=\ell^{\,v_\ell\left((p-1)\left((p+1)^2-a^2\right)\right)},
\]
and compute $\det(1-\Phi\mid V^D)$ and $\det(1-\Frob\mid V^D)$. Evaluate for \curve{11a1} at $p=2$, $\ell=5$. \emph{Answer:} the arithmetic Frobenius has eigenvalues
$p\alpha/\beta$, $p$, $p\beta/\alpha$ on $V^D$, of absolute value $p$; $\det(1-\Frob)=(1-p)\bigl((p+1)^2-a^2\bigr)$ and $\det(1-\Phi)=(1-p^{-1})\bigl((p+1)^2-a^2\bigr)/p^2$, of the same
$\ell$-adic valuation. For \curve{11a1}, $a_2=-2$ and $(p-1)((p+1)^2-a^2)=5$: the group has order $5$.
\end{exercise}

\begin{exercise}\label{ex:galcoh:ex-type-iv}
Let $E$ have reduction of type $\mathrm{IV}$ or $\mathrm{IV}^*$ at $q\ne3$, $\ell\ne q$ odd. Using Lemma lem:IV of the paper ($(\ad V_\ell E)^{I_q}$ is a line on which $\Frob_q$ acts by
$\chi_{-3}(q)$), show that $\#\Hur(\Q_q,T^D)=\ell^{v_\ell(q-\chi_{-3}(q))}$, and check the paper's value $e_3(E,2)=1$ at $q=2$. \emph{Hint:} \Cref{prop:galcoh:unramified-computed}(iv) with
eigenvalue $q\chi_{-3}(q)$; $\chi_{-3}(2)=-1$ and $v_3(3)=1$.
\end{exercise}

\begin{exercise}\label{ex:galcoh:ex-gw-ell}
Apply \Cref{thm:galcoh:greenberg-wiles} to $M=\Z/\ell$ ($\ell$ odd) with $\mathcal L_p=\Hur(\Q_p,\Z/\ell)$ for $p\ne\ell$, $\mathcal L_\ell=H^1(\Q_\ell,\Z/\ell)$ and $\mathcal L_\infty=0$. Compute both sides
and identify the cyclic extensions of $\Q$ of degree $\ell$ unramified outside $\ell$. \emph{Answer:} the dual Selmer group is $0$ (it is strict at $\ell$); the right side is
$\ell\cdot\frac1\ell\cdot\frac{\#H^1(\Q_\ell,\Z/\ell)}{\#H^0(\Q_\ell,\Z/\ell)}=\ell^2/\ell=\ell$, with $\#H^1(\Q_\ell,\Z/\ell)=\ell^2$ by \Cref{thm:galcoh:local-euler}. So there are $\ell$ homomorphisms, and
one such field: the subfield of degree $\ell$ of $\Q(\zeta_{\ell^2})$.
\end{exercise}

\begin{exercise}\label{ex:galcoh:ex-two-primes}
In \Cref{ex:galcoh:relaxing-toy} take $\Sigma=\{p_1,p_2\}$, two primes different from $\ell$. Check the hypotheses of \Cref{thm:galcoh:relaxing-sigma} and show that the characters of
$G_\Q$ of $\ell$-power order unramified outside $\{p_1,p_2\}$ form a group isomorphic to $\Z/\ell^{v_\ell(p_1-1)}\times\Z/\ell^{v_\ell(p_2-1)}$.
\end{exercise}

\begin{exercise}[the toy model of $h_3(E)$]\label{ex:galcoh:ex-h3}
Let $T=\Z_3(\chi_{-3})$, the module $\Z_3$ on which $G_\Q$ acts through the quadratic character $\chi_{-3}$. Show that $W^D=\mu_{3^\infty}\otimes\chi_{-3}$ and
$H^0(\Q,W^D)\simeq\Z/3$: use $\bar\chi_3^{\mathrm{cyc}}=\chi_{-3}$ for the elements of order $3$, and an element $\sigma$ with $\chi_3^{\mathrm{cyc}}(\sigma)=4$ for those of order $9$. Show that for
$\ell\ge5$ and every quadratic character $\psi$, $H^0(\Q,\mu_{\ell^\infty}\otimes\psi)=0$. Relate this to Lemma lem:h0 ($h_\ell(E)=0$ for $\ell\ge5$, $h_3(E)\le1$).
\emph{Hint:} $\sigma$ fixes $\zeta_3$ (as $4\equiv1\pmod3$), so $\chi_{-3}(\sigma)=1$, and $\sigma$ acts on $\mu_9\otimes\chi_{-3}$ by $4\not\equiv1\pmod9$. For $\ell\ge5$, $\bar\chi_\ell\psi$ is nontrivial because
$\bar\chi_\ell$ has order $\ell-1\ge4$; the determinant argument of Lemma lem:h0 is the same mechanism in rank two.
\end{exercise}

\begin{exercise}\label{ex:galcoh:ex-q2}
($\ell=p=2$.) Show that $\#\Q_2^\times/(\Q_2^\times)^2=8$, and check \Cref{thm:galcoh:local-euler} for $M=\mu_2$ over $\Q_2$. \emph{Answer:} $\Q_2^\times=2^{\Z}\times\{\pm1\}\times(1+4\Z_2)$ and
$1+4\Z_2\simeq\Z_2$, so the quotient is $(\Z/2)^3$; $\#H^0=\#H^2=2$, and $2\cdot2/8=\frac12=|2|_2$.
\end{exercise}

\begin{exercise}\label{ex:galcoh:ex-77c1}
Let $E=\curve{77c1}$ (nonsplit multiplicative of type $\mathrm I_3$ at $7$). Show that $E(\Q_7)[3]=0$ and deduce that $H^i(\Q_7,E[3])=0$ for $i=0,1,2$. \emph{Hint:}
$E(\Q_7)/E_0(\Q_7)$ has order $c_7=1$, $E_0/E_1\simeq\tilde E_{\mathrm{ns}}(\F_7)$ has order $7+1=8$, $E_1(\Q_7)$ is pro-$7$; then use
\Cref{cor:galcoh:local-duality-consequences} with $E[3]^D\simeq E[3]$.
\end{exercise}

\begin{exercise}\label{ex:galcoh:ex-good-h1}
Let $E$ have good reduction at $p\ne\ell$. Show that $H^1(\Q_p,T_\ell E)=\Hur(\Q_p,T_\ell E)$, finite of order $\ell^{v_\ell(\#\tilde E(\F_p))}$, and that it is isomorphic to $E(\Q_p)[\ell^\infty]$.
Check with \curve{11a1}, $p=2$, $\ell=5$. \emph{Hint:} $\dim H^1(\Q_p,V)=0$ by \Cref{cor:galcoh:local-duality-consequences}(ii) ($V$ and $V^D\simeq V$ have no invariants); $\det(1-\Frob\mid V)=1-a_p+p$;
for \curve{11a1}, $\#\tilde E(\F_2)=5$.
\end{exercise}

\begin{exercise}\label{ex:galcoh:ex-cft-orth}
(Harder.) Prove \Cref{thm:galcoh:ur-orthogonal}(i) for $M=\Z/\ell$, $M^D=\mu_\ell$, $p\ne\ell$, directly from local class field theory (\Cref{ch:local}): for a character
$\chi\colon G_{\Q_p}\to\Z/\ell$ and $a\in\Q_p^\times$, $\operatorname{inv}(\chi\cup\delta(a))=\pm\frac1\ell\chi(\operatorname{rec}(a))$, where $\operatorname{rec}$ is the reciprocity map. Deduce that unramified characters are
orthogonal to unit classes, and count.
\end{exercise}

\begin{exercise}[Greenberg-Wiles on the Kummer side]\label{ex:galcoh:ex-gw-mu}
Let $\ell$ be odd, $M=\mu_\ell$, and $\mathcal L$ the structure with $\mathcal L_p=\Hur(\Q_p,\mu_\ell)$ for $p\ne\ell$, $\mathcal L_\ell=H^1(\Q_\ell,\mu_\ell)$, $\mathcal L_\infty=0$. Using global Kummer
theory and \Cref{prop:galcoh:unit-classes}, show that $H^1_{\mathcal L}(\Q,\mu_\ell)$ is the image of $\ell^{\Z}$ in $\Q^\times/(\Q^\times)^\ell$, of order $\ell$. Identify the dual structure on
$\Z/\ell$ and show $H^1_{\mathcal L^*}(\Q,\Z/\ell)=0$. Check \Cref{thm:galcoh:greenberg-wiles}, and compare with \Cref{ex:galcoh:ex-gw-ell}, where the roles of $M$ and $M^D$ are exchanged.
\emph{Answer:} $a\in\Q^\times$ is a unit times an $\ell$-th power at each $p\ne\ell$ iff $\ell\mid v_p(a)$ for all $p\ne\ell$, iff $a\in\pm\ell^{\Z}(\Q^\times)^\ell$, and $-1=(-1)^\ell$. The dual
structure is unramified at $p\ne\ell$ (\Cref{thm:galcoh:ur-orthogonal}) and $0$ at $\ell$; a character in it is unramified everywhere, hence trivial (Minkowski). Right side:
$\frac{\#H^0(\Q,\mu_\ell)}{\#H^0(\Q,\Z/\ell)}=\frac1\ell$, the factors at $p\ne\ell$ and at $\infty$ are $1$ (\Cref{prop:galcoh:unramified-computed}(ii); $H^0(\R,\mu_\ell)=1$), and at
$\ell$ the factor is $\#H^1(\Q_\ell,\mu_\ell)/\#H^0(\Q_\ell,\mu_\ell)=\ell^2$ (\Cref{ex:galcoh:local-orders}(b)); so both sides equal $\ell$.
\end{exercise}
\section{Notes and references}\label{sec:galcoh:notes}

\emph{Galois cohomology of fields.} Serre's \emph{Galois Cohomology} \cite{SerreGaloisCoh} (Chapter I for profinite groups, Chapter II for number fields and local fields) and
his \emph{Local Fields} \cite{SerreLocalFields} (Chapters X to XIV: Hilbert 90, Kummer theory, the Brauer group, local class field theory in cohomological form) are the
sources for \Cref{sec:galcoh:field}; Neukirch, Schmidt and Wingberg \cite{NSW} (Chapters I, II, VI to VIII) is the comprehensive reference, and the articles of Serre and
Tate in Cassels and Fr\"ohlich \cite{CasselsFrohlich} remain the best short account of local and global class field theory with cohomology. Milne's \emph{Class Field
Theory} \cite{MilneCFT} is a free alternative.

\emph{Continuous cohomology and $\ell$-adic coefficients.} The passage from finite to $\ell$-adic coefficients (\Cref{prop:galcoh:tvw}) follows Tate \cite[Section 2]{Tate95} and
Rubin \cite[Sections 1.1 to 1.3 and Appendix B]{Rubin}. Rubin's first chapter is the closest model for \Cref{sec:galcoh:unramified,sec:galcoh:structures,sec:galcoh:comparison}:
unramified and propagated local conditions, Selmer groups defined by local conditions, and the comparison exact sequence. The term ``Selmer structure'', and the
relaxed and strict structures, are from Mazur and Rubin \cite[Chapter 2]{MazurRubin04}.

\emph{Duality.} Tate announced local and global duality and the Euler characteristic formulas at the Stockholm congress \cite{TateDuality}; complete proofs are in
\cite[Chapter II, Section 5]{SerreGaloisCoh} (local), and in Milne \cite[Chapter I]{MilneADT} and \cite[Chapters VII and VIII]{NSW} (local and global, including the
Poitou-Tate sequence). The formula of \Cref{thm:galcoh:greenberg-wiles} is due to Greenberg and to Wiles; it is stated in \cite[Section 1]{Wiles95} and proved in
\cite[Section 2]{DDT}. Greenberg's survey \cite{Greenberg89} treats Selmer groups of $\ell$-adic representations and the comparison sequence in the form used by Diamond,
Flach and Guo \cite[Section 2]{DFG04}.

\emph{Elliptic curves.} \Cref{sec:galcoh:selmer-ec} follows Silverman \cite[Chapters VIII and X]{SilvermanAEC}; the Tate curve and its Galois representation are in
\cite[Chapter V]{SilvermanATAEC} and \cite[Chapter IV, Appendix]{SerreAbelian}. Tate's theorem that the local Kummer image is its own annihilator is in
\cite[Chapter I, Section 3]{MilneADT}.

\emph{The paper's sources.} Lemma lem:dfg21 is \cite[Lemma 2.1]{DFG04} and Theorem thm:dfg27 is \cite[Theorem 2.7]{DFG04}, as the paper cites them; both numbers and statements
were checked against the published text (Lemma 2.1 on p.~696, with the hypotheses that $\Sigma$ is nonempty, contains every prime at which $H^1_f(\Q_p,W)\ne H^1_{\ur}(\Q_p,W)$, and
$H^0(\Q,V^D)=H^1_f(\Q,V^D)=0$; Theorem 2.7 on p.~699). The formula of Fontaine and Perrin-Riou is \cite[I.4.2.2]{FPR94}, and the Bloch-Kato conditions are in \cite[Section 3]{BlochKato};
these two numbers are the paper's.

\emph{Citations.} Every reference in this chapter to a book is to a chapter or section, not to a numbered theorem. Theorem numbers we recall but have not checked
against the printed editions, and which the reader may find useful as unverified pointers, are: Serre, \emph{Galois Cohomology}, II.5.2 (local duality) and II.5.7
(Euler characteristic); Milne, \emph{Arithmetic Duality Theorems}, I.2.3, I.2.6, I.2.8 and I.4.10; Wiles \cite{Wiles95}, Proposition 1.6; Darmon, Diamond and Taylor
\cite{DDT}, Theorem 2.19; Rubin \cite{Rubin}, Theorem 1.7.3 and Lemma 1.3.2; Bloch and Kato \cite{BlochKato}, Proposition 3.8; Silverman \cite{SilvermanAEC}, X.4.2.

\emph{Computations.} The numbers in \Cref{ex:galcoh:local-orders,ex:galcoh:tate-kummer-examples,ex:galcoh:selmer-11a1} and in the exercises were computed with PARI/GP \cite{PARI}
2.17.2 through \texttt{cypari2} (\texttt{ellinit}, \texttt{ellinit(E,O(q\^{}30)).tate}, \texttt{elllocalred}, \texttt{ellap}, \texttt{ellmoddegree}, \texttt{elldivpol},
\texttt{polrootspadic}, \texttt{ispower}) and compared with Cremona's tables \cite{Cremona97} (files \texttt{allcurves}, \texttt{allbsd}, \texttt{degphi}, \texttt{galrep},
conductor below $10{,}000$).
